% Luận văn / đồ án tốt nghiệp: máy học dự đoán bệnh đái tháo đường (cấu trúc theo chương)
% Biên dịch: XeLaTeX hoặc LuaLaTeX (hỗ trợ tiếng Việt Unicode)

\documentclass[14pt,a4paper]{extreport}

% Gói tiếng Việt (Unicode)
\usepackage{fontspec}
\usepackage{polyglossia}
\setmainlanguage{vietnamese}
\setmainfont{Times New Roman}

% Layout và hình thức
\usepackage[margin=2.5cm]{geometry}
\usepackage{graphicx}
\usepackage{booktabs}
\usepackage{framed}
\usepackage[normalem]{ulem}
\usepackage{indentfirst}
\usepackage{setspace}
\onehalfspacing
\setlength{\parindent}{1.25cm}
\usepackage{tocloft}
% Mục lục: hiển thị "Chương n" thay vì chỉ số n
\setlength{\cftchapnumwidth}{7.2em}
\renewcommand{\cftchappresnum}{Chương~}
\renewcommand{\cftchapaftersnum}{}
\setlength{\cftsecnumwidth}{5.2em}
\setlength{\cftsubsecnumwidth}{2.8em}

% Định dạng tiêu đề chương: số chương + tiêu đề cùng dòng (tránh xuống hàng tách khỏi "Chương n")
\usepackage{titlesec}
\titleformat{\chapter}[hang]
  {\normalfont\huge\bfseries\raggedright}
  {\chaptername\ \thechapter.}
  {0.75em}
  {}
\titlespacing*{\chapter}{0pt}{3.5ex plus 1ex minus .2ex}{2.3ex plus .2ex}
% Chương không đánh số (Tóm tắt, Phụ lục): giữ kiểu khối căn giữa như mặc định
\titleformat{name=\chapter,numberless}[display]
  {\normalfont\huge\bfseries\centering}
  {}
  {0pt}
  {}

% Trích dẫn và tài liệu tham khảo
\usepackage{csquotes}
\DeclareQuoteStyle{vietnamese}
  {\guillemotleft}{\guillemotright}
  {\textquotedblleft}{\textquotedblright}
\usepackage[backend=biber,style=ieee,sorting=none,bibencoding=utf8,giveninits=false]{biblatex}
\addbibresource{references.bib}
\DeclareFieldFormat{labelnumberwidth}{\mkbibbrackets{#1}\addperiod}
\ExecuteBibliographyOptions{giveninits=false}
\setlength{\emergencystretch}{5em}
\AtBeginBibliography{\sloppy}

\DeclareBibliographyDriver{book}{%
  \usebibmacro{bibindex}%
  \usebibmacro{begentry}%
  \printnames{author}%
  \setunit{\addcomma\space}%
  \printfield{year}%
  \setunit{\addcomma\space}%
  \printfield{title}%
  \setunit{\addcomma\space}%
  \printlist{publisher}%
  \iflistundef{location}{}{\setunit{\addcomma\space}\printlist{location}}%
  \usebibmacro{finentry}%
}

\DeclareBibliographyDriver{incollection}{%
  \usebibmacro{bibindex}%
  \usebibmacro{begentry}%
  \printnames{author}%
  \setunit{\addcomma\space}%
  \printfield{year}%
  \setunit{\addcomma\space}%
  \printfield{title}%
  \setunit{\addcomma\space}%
  \printfield{booktitle}%
  \setunit{\addcomma\space}%
  \printlist{publisher}%
  \usebibmacro{finentry}%
}

\DeclareBibliographyDriver{periodical}{%
  \usebibmacro{bibindex}%
  \usebibmacro{begentry}%
  \printfield{title}%
  \setunit{\addcomma\space}%
  \printlist{publisher}%
  \iflistundef{location}{}{\setunit{\addcomma\space}\printlist{location}}%
  \usebibmacro{finentry}%
}

\usepackage[unicode,colorlinks=true,linkcolor=blue,urlcolor=blue,citecolor=blue,
  bookmarks=true,bookmarksopen=true,bookmarksnumbered=true,bookmarksdepth=3,
  pdfencoding=unicode]{hyperref}
\usepackage[open,numbered,depth=3]{bookmark}
\usepackage{tikz}
\usetikzlibrary{arrows.meta,positioning}
\usepackage{amsmath}
\usepackage{amssymb}

% Thông tin bìa (có thể sửa trực tiếp tại đây)
\newcommand{\hocphan}{Khoa học máy tính / Nghiên cứu khoa học}
\newcommand{\tenduan}{Nghiên cứu và đánh giá hiệu quả của các thuật toán máy học trong việc dự đoán bệnh đái tháo đường}
\newcommand{\giangvien}{TS. Đinh Thị Thu Hương}
\newcommand{\hocvien}{Nông Văn Tình}
\newcommand{\mahocvien}{8480101250028}
\newcommand{\lop}{Khoa học máy tính}
\newcommand{\malop}{25MCS2}
\newcommand{\nganh}{Khoa học máy tính}
\newcommand{\manganh}{8480101}
\newcommand{\thangnam}{.... / .....}
\newcommand{\noidungnghiencuu}{\tenduan}

\begin{document}

%==============================================================================
% BÌA 1 (trang 1 - theo mẫu BIA DO AN CUOI MON MCS.docx)
%==============================================================================
\thispagestyle{empty}
\begin{center}
\vspace*{0.3cm}
{\fontsize{14}{18}\selectfont\bfseries BỘ GIÁO DỤC VÀ ĐÀO TẠO}\\[0.5cm]
{\fontsize{14}{18}\selectfont\bfseries TRƯỜNG ĐẠI HỌC TƯ THỤC QUỐC TẾ SÀI GÒN}\\[0.5cm]
\includegraphics[height=1.8cm]{image/logo.png}\\[0.3cm]

\vspace{0.8cm}
{\fontsize{18}{22}\selectfont\bfseries BÀI ĐỒ ÁN CUỐI KỲ}\\[1.2cm]

{\fontsize{14}{18}\selectfont\bfseries TÊN ĐỒ ÁN:}\\[0.3cm]
{\fontsize{14}{18}\selectfont\bfseries \tenduan}\\[1.5cm]

\begin{minipage}[t]{0.85\textwidth}
\fontsize{13}{17}\selectfont
\bfseries Giảng viên giảng dạy:\quad \giangvien\\[0.4cm]
\bfseries Học viên thực hiện:\quad \hocvien\\[0.4cm]
\bfseries Mã học viên:\quad \mahocvien\\[0.4cm]
\bfseries Lớp:\quad \lop \hspace{2cm} \bfseries Mã số lớp:\quad \malop\\[1.2cm]
\end{minipage}

\vfill
{\fontsize{13}{17}\selectfont\itshape TP. Hồ Chí Minh, tháng \thangnam}\\[0.6cm]
\end{center}

%==============================================================================
% BÌA 2 (trang 2 - theo mẫu BIA DO AN CUOI MON MCS.docx)
%==============================================================================
\newpage
\thispagestyle{empty}
{\setlength{\fboxrule}{4.5pt}\setlength{\fboxsep}{12pt}%
\noindent\fbox{\begin{minipage}{\dimexpr\textwidth-2\fboxsep-2\fboxrule}
\begin{center}
\vspace*{0.3cm}
{\fontsize{14}{18}\selectfont\bfseries BỘ GIÁO DỤC VÀ ĐÀO TẠO}\\[0.5cm]
{\fontsize{14}{18}\selectfont\bfseries TRƯỜNG ĐẠI HỌC TƯ THỤC QUỐC TẾ SÀI GÒN}\\[0.8cm]
{\fontsize{14}{18}\selectfont\bfseries *********}\\[1cm]

\includegraphics[width=11cm]{image/full-name-and-logo.png}\\[0.7cm]

{\fontsize{18}{22}\selectfont\bfseries BÀI CUỐI KỲ}\\[1.2cm]

\begin{tabular}{l l}
{\fontsize{14}{18}\selectfont\bfseries Ngành:} & \nganh \\[0.3cm]
{\fontsize{14}{18}\selectfont\bfseries Mã ngành:} & \manganh \\[0.3cm]
{\fontsize{14}{18}\selectfont\bfseries Mã lớp:} & \malop \\[0.3cm]
\end{tabular}

\end{center}
\begin{flushleft}
{\fontsize{14}{18}\selectfont\bfseries \uline{Nội dung nghiên cứu:}}\\[0.3cm]
\hspace{1cm}\noidungnghiencuu
\end{flushleft}
\vspace{1.2cm}

\begin{center}
\vfill
{\fontsize{13}{17}\selectfont\itshape TP. Hồ Chí Minh, tháng \thangnam}\\[0.8cm]
\end{center}

\end{minipage}}}

\begin{center}
  % Bảng chấm điểm
\begin{tabular}{|p{2.5cm}|p{3cm}|p{5cm}|p{3.5cm}|}
\hline
\multicolumn{2}{|c|}{\textbf{Điểm}} 
& \multicolumn{1}{c|}{\textbf{Nhận xét}} 
& \multicolumn{1}{c|}{\textbf{Giảng viên ký}} \\
\cline{1-2}
\textbf{Bằng số} & \textbf{Bằng chữ} & & \\
\hline
&  &  &  \\[2.5cm]
\hline
\end{tabular}
\end{center}

\newpage
\tableofcontents
\newpage
\setcounter{page}{1}

%==============================================================================
\chapter*{Tóm tắt (Abstract)}
\addcontentsline{toc}{chapter}{Tóm tắt (Abstract)}
%==============================================================================

\begin{quote}
\itshape
\noindent
Đái tháo đường hiện là một gánh nặng y tế toàn cầu, trong đó việc phát hiện sớm đóng vai trò then chốt giúp giảm thiểu biến chứng và chi phí điều trị. Nghiên cứu này thực hiện so sánh bốn thuật toán phân loại cổ điển bao gồm: hồi quy logistic, SVM (kernel RBF), rừng ngẫu nhiên (Random Forest) và KNN trên bài toán phân loại nhị phân nguy cơ đái tháo đường (tập dữ liệu Pima Indians Diabetes, $n=768$).

\noindent
Quy trình xử lý (pipeline) được thiết kế đảm bảo tính kiểm soát, khả năng tái lập và chống rò rỉ dữ liệu (\textit{data leakage}): việc xử lý giá trị thiếu bằng trung vị và chuẩn hóa đặc trưng chỉ được thực hiện trong phạm vi tập huấn luyện. Các siêu tham số được tối ưu hóa dựa trên điểm $F_1$ thông qua phương pháp xác thực chéo phân tầng 5-fold ($5\text{-fold stratified cross-validation}$) trên tập huấn luyện. Cuối cùng, mô hình được đánh giá trên tập kiểm tra độc lập ($hold\text{-}out$ $20\%$) có phân tầng, thực hiện lặp lại với 5 giá trị hạt giống ($seed$) khác nhau để đảm bảo tính khách quan.

\noindent
\textbf{Rừng ngẫu nhiên} (Random Forest - RF) đạt điểm $F_1$ trung bình cao nhất ($0{,}614 \pm 0{,}025$; CI 95\%: $[0{,}583;\,0{,}645]$) cùng chỉ số Brier trung bình thấp nhất ($0{,}163 \pm 0{,}012$). Trong khi đó, \textbf{hồi quy logistic} dẫn đầu về độ chính xác (Accuracy) và ROC-AUC trung bình (lần lượt là $0{,}748 \pm 0{,}018$ và $0{,}824 \pm 0{,}031$), đồng thời có chỉ số Brier xấp xỉ RF. Kết quả này phản ánh sự đánh đổi giữa khả năng mô hình hóa phi tuyến phức tạp và năng lực xếp hạng toàn cục ổn định của mô hình tuyến tính trên cỡ mẫu hạn chế. Chỉ số PR-AUC trung bình cũng ghi nhận kết quả tốt nhất ở mô hình RF; nghiên cứu bổ sung đánh giá bằng bootstrap ROC-AUC và Cliff's delta. Kiểm định Wilcoxon signed-rank trên điểm $F_1$ (so sánh cặp RF--KNN, $n=5$ lần chia tập) cho thấy sự khác biệt không có ý nghĩa thống kê ($p=0{,}44$); \emph{kết quả này cần được xem xét thận trọng do lực thống kê (statistical power) thấp khi số lần lặp hạn chế, không loại trừ khả năng tồn tại những chênh lệch thực tế nhỏ}. Về mặt lâm sàng, các đặc trưng Glucose và BMI đóng vai trò quan trọng nhất, nhất quán qua hệ số hồi quy, độ quan trọng đặc trưng trong RF và tầm quan trọng hoán vị (\textit{permutation importance}). Tóm lại, một quy trình đánh giá nghiêm ngặt là điều kiện tiên quyết để so sánh khách quan các mô hình trên dữ liệu y tế dạng bảng quy mô nhỏ; dù mô hình phi tuyến có khả năng khai thác các tương tác đặc trưng tốt hơn, nhưng ưu thế của chúng so với các mô hình đơn giản không phải lúc nào cũng đủ mạnh về mặt thống kê.

\noindent
Nghiên cứu này còn tồn tại một số hạn chế như quy mô mẫu nhỏ, đặc điểm quần thể Pima có tính đồng nhất cao và thiếu bước kiểm định ngoại mẫu (\textit{external validation}). Do đó, các kết luận rút ra chủ yếu mang giá trị về mặt phương pháp luận và chưa thể khái quát hóa trực tiếp cho các bối cảnh lâm sàng rộng lớn hơn.
\end{quote}

\noindent\textbf{Từ khóa (Keywords):} machine learning; diabetes prediction; clinical decision support; tabular biomedical data; binary classification; random forest; model evaluation protocol; reproducibility.

\chapter{Giới thiệu}
\label{sec:intro}

\section{Lý do đề tài}

Đái tháo đường (\textit{diabetes mellitus}) là bệnh lý mãn tính đang bùng nổ trên toàn cầu, gây quá tải cho hệ thống y tế và bào mòn chất lượng sống của người bệnh. Các hướng dẫn lâm sàng hiện nay đều nhấn mạnh việc tầm soát và can thiệp sớm là yếu tố sống còn để ngăn chặn các biến chứng nguy hiểm như tim mạch, suy thận hay tổn thương thần kinh. Dù vậy, chẩn đoán sớm vẫn là một thách thức lớn khi các triệu chứng ban đầu thường mơ hồ, đi kèm với cơ chế bệnh sinh phức tạp bởi nhiều yếu tố sinh học đan xen.

Trước bối cảnh đó, học máy (\textit{machine learning}) đã mở ra hướng đi đầy triển vọng trong việc hỗ trợ dự đoán và sàng lọc từ dữ liệu lâm sàng. Khả năng khai phá những quy luật ẩn sâu trong dữ liệu giúp các thuật toán này thiết lập nên những mô hình phân loại có độ chính xác cao. Bên cạnh đó, việc khai thác các bộ dữ liệu y tế công khai không chỉ giúp tối ưu hóa việc đối chiếu các mô hình mà còn mang lại giá trị thực tiễn lớn về mặt phương pháp luận. \emph{Cần lưu ý rằng các công cụ này chỉ đóng vai trò hỗ trợ ra quyết định và không thay thế cho quy trình chẩn đoán chuyên môn của đội ngũ bác sĩ.}

Tuy nhiên, hiệu năng của các thuật toán thường biến động lớn, phụ thuộc chặt chẽ vào đặc thù dữ liệu, quy trình tiền xử lý và giao thức thực nghiệm. Nhiều nghiên cứu trước đây đã đưa ra những kết quả thiếu nhất quán hoặc khó đối chiếu trực tiếp, chủ yếu do sự khác biệt trong cách xử lý đầu vào và tiêu chí đánh giá. Thực trạng này đặt ra yêu cầu cấp thiết về một nghiên cứu đối chiếu có hệ thống trên cùng một quy trình (\textit{pipeline}) chuẩn hóa, nhằm đánh giá khách quan hiệu quả của các mô hình trong bối cảnh dữ liệu y tế dạng bảng có quy mô hạn chế. Các lập luận chi tiết và phân tích chuyên sâu sẽ được trình bày cụ thể tại Chương~\ref{sec:literature}.

Từ những lập luận trên, nghiên cứu này xây dựng một quy trình thực nghiệm chuẩn hóa nhằm đánh giá toàn diện hiệu năng của các thuật toán học máy truyền thống trong bài toán phân loại nhị phân nguy cơ đái tháo đường. Kết quả kỳ vọng không chỉ đưa ra cái nhìn đối chiếu khách quan giữa các mô hình, mà còn góp phần khẳng định tầm quan trọng của giao thức thực nghiệm đối với độ tin cậy và khả năng tổng quát hóa của mô hình. Đây chính là nền tảng cho các mục tiêu và đóng góp cụ thể sẽ được trình bày chi tiết trong các chương tiếp theo.

\section{Mục tiêu nghiên cứu}
Nghiên cứu này tập trung vào bài toán \emph{phân loại nhị phân}: dự báo nguy cơ khởi phát đái tháo đường dựa trên tập hợp các đặc trưng lâm sàng và chỉ số sinh hiệu. Trọng tâm của nghiên cứu là xây dựng một giao thức thực nghiệm chặt chẽ nhằm huấn luyện và đối chiếu các mô hình học máy, qua đó loại bỏ các sai số lạc quan hệ thống phát sinh từ hiện tượng rò rỉ dữ liệu (\textit{data leakage}) hoặc sự thiếu nhất quán trong quy trình tiền xử lý giữa các mô hình so sánh.

Từ bối cảnh và các khoảng trống nghiên cứu đã xác định, \textbf{mục tiêu tổng quát} của đề tài là xây dựng quy trình thực nghiệm và đánh giá có hệ thống hiệu năng của các thuật toán học máy truyền thống trên dữ liệu y sinh dạng bảng quy mô hạn chế (thực nghiệm trên tập dữ liệu Pima Indians Diabetes, $n=768$).

\textbf{Mục tiêu cụ thể} được hiện thực hóa qua các nhiệm vụ trọng tâm sau:

\begin{enumerate}
  \item \textbf{Xây dựng quy trình thực nghiệm chuẩn hóa và có khả năng tái lập}: Thiết lập luồng xử lý (\textit{pipeline}) tích hợp xuyên suốt từ tiền xử lý dữ liệu (xử lý giá trị thiếu bằng trung vị, chuẩn hóa đặc trưng) đến giai đoạn huấn luyện và \textbf{đánh giá độc lập} (\textit{stratified hold-out} lặp lại với nhiều hạt giống ngẫu nhiên). Quy trình này đảm bảo tính khách quan trong đối chiếu và triệt tiêu hiện tượng rò rỉ dữ liệu (\textit{data leakage}). Toàn bộ \textbf{tám đặc trưng gốc} được giữ nguyên nhằm bảo toàn giá trị thông tin theo mô tả tại Chương~\ref{sec:methods}.
  \item \textbf{Triển khai và tối ưu hóa các bộ phân loại}: Thực hiện huấn luyện và tinh chỉnh siêu tham số cho bốn thuật toán tiêu biểu: Hồi quy Logistic, Máy vectơ hỗ trợ (SVM) nhân RBF, $k$-láng giềng gần nhất (KNN) và Rừng ngẫu nhiên (Random Forest). Quá trình tối ưu hóa áp dụng \textbf{xác thực chéo phân tầng 5-fold} trên tập huấn luyện, lấy điểm $F_1$ làm tiêu chí mục tiêu chủ đạo.
  \item \textbf{Đánh giá và đối chiếu hiệu năng đa chiều}: Sử dụng hệ thống chỉ số toàn diện bao gồm Accuracy, Precision, Recall, $F_1$-score, ROC--AUC và Brier score (nhằm đánh giá độ hiệu chuẩn xác suất). Hiệu năng và độ ổn định của mô hình được tổng hợp qua thống kê mô tả (trung bình $\pm$ độ lệch chuẩn) trên nhiều lần phân chia dữ liệu, giúp loại bỏ các kết luận chủ quan từ thực nghiệm đơn lẻ.
  \item \textbf{Phân tích mức độ ảnh hưởng của đặc trưng}: Đánh giá tầm quan trọng của các biến đầu vào đối với kết quả dự báo (dựa trên hệ số chuẩn hóa của mô hình tuyến tính và \textit{feature importance} của mô hình phi tuyến), từ đó xác định các yếu tố then chốt và đối chiếu với các bằng chứng lâm sàng.
  \item \textbf{Kiểm định thống kê và thảo luận thực tiễn}: Áp dụng các kiểm định phi tham số (như Wilcoxon signed-rank) để xác nhận ý nghĩa thống kê của sự khác biệt hiệu năng; đồng thời thảo luận về giới hạn tổng quát hóa và điều kiện triển khai thực tế. Nghiên cứu tái khẳng định các công cụ này chỉ đóng vai trò hỗ trợ tầm soát, không thay thế chẩn đoán chuyên môn.
\end{enumerate}

Thông qua lộ trình trên, nghiên cứu hướng tới việc đề xuất \textbf{một khung tham chiếu} chuẩn mực cho việc đánh giá và công bố kết quả học máy trên dữ liệu y sinh dạng bảng quy mô nhỏ. Chi tiết về phương pháp triển khai và kết quả thực nghiệm sẽ lần lượt được trình bày tại Chương~\ref{sec:methods} và Chương~\ref{sec:results}.

\section{Câu hỏi nghiên cứu}

Để hiện thực hóa các mục tiêu đã đề ra và xây dựng hệ thống thực nghiệm có khả năng kiểm chứng cao, nghiên cứu tập trung giải quyết hai câu hỏi nghiên cứu (\textit{Research Questions} - RQ) trọng tâm sau:

\begin{itemize}
  \item \textbf{RQ1:} Trong điều kiện quy trình tiền xử lý và đặc tính dữ liệu đã thiết lập, thuật toán phân loại nào thể hiện sự vượt trội và tính ổn định cao nhất dựa trên các tiêu chí Accuracy, $F_1$-score và ROC--AUC?
  \item \textbf{RQ2:} Những đặc trưng lâm sàng nào đóng vai trò tiên lượng then chốt trong mô hình, và mức độ quan trọng của các biến số này có sự tương đồng với các bằng chứng y văn hiện hành hay không?
\end{itemize}


\section{Đóng góp của nghiên cứu}

Bên cạnh việc triển khai và công bố hiệu năng của các mô hình cụ thể, nghiên cứu này đóng góp \textbf{giá trị thực tiễn về mặt phương pháp luận} thông qua việc chuẩn hóa giao thức thực nghiệm và quy trình diễn giải kết quả trên dữ liệu y sinh dạng bảng quy mô nhỏ. Đóng góp trọng tâm của luận văn \textbf{không nằm ở việc xây dựng thuật toán mới}, mà tập trung thiết lập một \textbf{khung đánh giá nghiêm ngặt, có khả năng tái lập và ngăn chặn rò rỉ dữ liệu} (\textit{leakage-free}). Cụ thể, luận văn nhấn mạnh ba khía cạnh then chốt: (i) khung thực nghiệm chuẩn cho dữ liệu y sinh cỡ mẫu hạn chế; (ii) phân tích hệ thống về sự đánh đổi giữa hiệu năng và độ ổn định; (iii) khẳng định rằng \textbf{giao thức đánh giá quyết định trực tiếp đến thứ hạng mô hình}, và các chênh lệch về chỉ số trung bình thường \textbf{không mang ý nghĩa thống kê} khi quy trình thực nghiệm được kiểm soát chặt chẽ.

Các đóng góp cụ thể được tóm lược như sau:

\begin{enumerate}
  \item \textbf{Đề xuất quy trình huấn luyện--đánh giá thống nhất}: Xây dựng luồng xử lý (\textit{pipeline}) đảm bảo các bước tiền xử lý chỉ được thiết lập (\texttt{fit}) dựa trên dữ liệu huấn luyện. Quy trình tách biệt rõ ràng giữa tinh chỉnh siêu tham số và đánh giá độc lập (\textit{stratified hold-out}), kết hợp lặp lại với nhiều hạt giống ngẫu nhiên (\textit{seed}) nhằm \textbf{triệt tiêu rò rỉ thông tin} và đảm bảo tính khách quan (chi tiết tại Chương~\ref{sec:methods}).
  \item \textbf{Thực hiện đối chiếu hệ thống trên cùng một hệ quy chiếu}: So sánh hiệu năng của bốn thuật toán tiêu biểu trong cùng một điều kiện thực nghiệm đồng nhất, thay vì đối chiếu gián tiếp với các kết quả có giao thức khác biệt từ các công trình trước đó.
  \item \textbf{Áp dụng chiến lược đánh giá đa chỉ số}: Báo cáo đồng thời các chỉ số Accuracy, Precision, Recall, $F_1$-score, ROC--AUC và Brier score (đánh giá độ hiệu chuẩn). Nghiên cứu chú trọng vào \textbf{độ ổn định} thông qua thống kê mô tả và khoảng tin cậy, giúp loại bỏ các kết luận thiếu căn cứ từ các lần phân chia dữ liệu ngẫu nhiên đơn lẻ.
  \item \textbf{Tăng cường kiểm định thống kê}: Sử dụng kiểm định Wilcoxon signed-rank để phân biệt rõ rệt giữa \emph{cải thiện thực chất} và \emph{biến động do ngẫu nhiên}, từ đó củng cố tính xác thực cho các nhận định về hiệu năng mô hình.
  \item \textbf{Cung cấp góc nhìn diễn giải đa mô hình}: Sử dụng hệ số hồi quy, tầm quan trọng đặc trưng dựa trên độ tạp chất (\textit{impurity-based}) và tầm quan trọng hoán vị (\textit{permutation importance}) để đối chiếu công bằng giữa các họ mô hình, kết hợp với phân tích sai số (\textit{error analysis}) để làm rõ ý nghĩa lâm sàng (Chương~\ref{sec:results}, Chương~\ref{sec:discussion}).
  \item \textbf{Xác định giới hạn và khả năng tổng quát hóa}: Chỉ rõ các hạn chế cố hữu của dữ liệu quy mô nhỏ, đồng thời khẳng định vai trò của \textbf{quy trình đánh giá chặt chẽ} trong việc xác lập độ tin cậy và tiềm năng ứng dụng thực tế của mô hình.
\end{enumerate}

\section{Định nghĩa bài toán}

Nghiên cứu này tiếp cận bài toán dưới dạng \emph{phân loại nhị phân có giám sát}. Xét tập dữ liệu $\mathcal{D} = \{(\mathbf{x}_i, y_i)\}_{i=1}^{N}$, trong đó $\mathbf{x}_i \in \mathbb{R}^d$ là vectơ đặc trưng của mẫu thứ $i$ ($d=8$ đối với tập dữ liệu Pima) và $y_i \in \{0,1\}$ là nhãn mục tiêu tương ứng. Theo quy ước, giá trị $y_i=1$ đại diện cho lớp dương tính (có nguy cơ hoặc đã khởi phát đái tháo đường) và $y_i=0$ đại diện cho lớp âm tính (không có nguy cơ).

Mục tiêu cốt lõi của nghiên cứu là xây dựng một hàm dự báo có khả năng tổng quát hóa tốt trên những dữ liệu chưa từng xuất hiện trong quá trình huấn luyện. Hàm số này được biểu diễn dưới dạng một ánh xạ $f:\mathbb{R}^d \to \{0,1\}$, hoặc thông qua việc ước lượng xác suất hậu nghiệm $\hat{P}(y=1 \mid \mathbf{x})$ kết hợp với một ngưỡng quyết định (\textit{decision threshold}) phù hợp. Hiệu năng của mô hình được đánh giá định lượng trên tập kiểm thử độc lập thông qua hệ thống chỉ số toàn diện bao gồm: Accuracy, Precision, Recall, $F_1$-score và ROC--AUC. Các chi tiết kỹ thuật về quy trình tiền xử lý, huấn luyện và giao thức thực nghiệm sẽ được trình bày cụ thể tại Chương~\ref{sec:methods}.

\section{Cấu trúc luận văn}

Nội dung của luận văn được tổ chức thành sáu chương chính với lộ trình triển khai như sau:

\textbf{Chương~\ref{sec:intro} (Mở đầu):} Phác thảo bức tranh tổng quan về bối cảnh nghiên cứu, từ đó xác định mục tiêu, các câu hỏi nghiên cứu trọng tâm và những đóng góp cốt lõi của đề tài.

\textbf{Chương~\ref{sec:literature} (Cơ sở lý thuyết):} Hệ thống hóa nền tảng lý thuyết về ứng dụng học máy trong dự báo y sinh; tổng quan các công trình liên quan đến dự đoán nguy cơ đái tháo đường trên dữ liệu dạng bảng, đồng thời phân tích những thách thức trong việc đối chiếu hiệu năng giữa các nghiên cứu hiện hành.

\textbf{Chương~\ref{sec:methods} (Phương pháp nghiên cứu):} Mô tả chi tiết khung thực nghiệm đề xuất, đặc điểm tập dữ liệu Pima, quy trình tiền xử lý và các thuật toán triển khai. Chương này tập trung vào chiến lược tối ưu hóa siêu tham số, giao thức đánh giá đa tầng (\textit{nested cross-validation} và \textit{stratified hold-out}) cùng các kỹ thuật giải thích mô hình.

\textbf{Chương~\ref{sec:results} (Kết quả thực nghiệm):} Công bố kết quả thực nghiệm chi tiết, so sánh hiệu năng giữa các mô hình dựa trên hệ thống chỉ số đa chiều và độ hiệu chuẩn xác suất, kết hợp với các phân tích trực quan hóa dữ liệu.

\textbf{Chương~\ref{sec:discussion} (Thảo luận):} Phân tích chuyên sâu các kết quả thu được nhằm giải quyết triệt để hai câu hỏi nghiên cứu RQ1 và RQ2. Chương này cũng thảo luận về hàm ý thực tiễn, phân tích sai số, các hạn chế hiện tại và đề xuất định hướng phát triển tương lai.

\textbf{Chương~\ref{sec:conclusion} (Kết luận):} Tổng kết những thành tựu đạt được và đưa ra các nhận định khái quát về giá trị đóng góp của nghiên cứu.

Cuối cùng, phần \textbf{Phụ lục} cung cấp chi tiết về cấu hình siêu tham số, kết quả thực nghiệm theo từng hạt giống (\textit{seed}), môi trường thực thi và hướng dẫn tái lập mã nguồn nhằm đảm bảo tính minh bạch và khả năng kiểm chứng của toàn bộ nghiên cứu.

\chapter{Cơ sở lý thuyết}
\label{sec:literature}

\section{Tổng quan nghiên cứu}

Trong lĩnh vực dự báo nguy cơ đái tháo đường từ dữ liệu lâm sàng dạng bảng, tập dữ liệu \emph{Pima Indians Diabetes} từ lâu đã trở thành một chuẩn tham chiếu (\textit{benchmark}) kinh điển để đánh giá và đối chiếu hiệu năng giữa các thuật toán.\cite{smith1988,uci-pima} Thay vì liệt kê rời rạc các công trình đơn lẻ, các hướng tiếp cận hiện hành có thể được hệ thống hóa thành các nhóm trọng tâm sau:

\textbf{(i) Các mô hình học máy truyền thống trên tập dữ liệu Pima:} Từ những phương pháp sơ khởi như thuật toán ADAP trong nghiên cứu gốc đến các bộ phân loại phổ biến hiện nay như Hồi quy Logistic, SVM, $k$-NN và Rừng ngẫu nhiên; nhiều công trình đã sử dụng Pima làm tập dữ liệu thực nghiệm tiêu chuẩn sau khi thực hiện các bước tiền xử lý cơ bản.\cite{smith1988,sisodia2018,chang2022pima,miao2021pima} Nhóm nghiên cứu này đóng vai trò thiết lập các mô hình cơ sở (\textit{baseline}) quan trọng để đánh giá mức độ cải thiện của các phương pháp đề xuất mới.

\textbf{(ii) Các phương pháp học máy kết hợp (\textit{Ensemble}) và xử lý mất cân bằng dữ liệu:} Một hướng tiếp cận khác tập trung khai thác sức mạnh của các mô hình tổ hợp (như XGBoost) hoặc áp dụng các kỹ thuật lấy mẫu lại (như SMOTE) nhằm cải thiện độ nhạy (\textit{sensitivity}) đối với lớp dương tính. Các công trình này thường công bố độ chính xác hoặc chỉ số AUC vượt trội trong những thiết lập thực nghiệm riêng biệt.\cite{kirbas2025shap,zhang2025pima}

\textbf{(iii) Tính diễn giải của mô hình (\textit{Model Interpretability}):} Song song với nỗ lực tối ưu hóa hiệu năng, xu hướng nghiên cứu hiện đại nhấn mạnh vào khả năng giải thích các quyết định dự báo thông qua phân tích tầm quan trọng đặc trưng hoặc sử dụng các kỹ thuật như SHAP/LIME. Cách tiếp cận này giúp kết nối kết quả dự báo với các biến số lâm sàng thực tế, vốn là một yêu cầu then chốt trong bối cảnh ứng dụng y tế.\cite{kirbas2025shap,chang2022pima}

Mặc dù đạt được nhiều bước tiến, sự đa dạng trong quy trình tiền xử lý và giao thức thực nghiệm đã dẫn đến một hệ thống kết quả \textbf{thiếu tính nhất quán}, ngay cả khi thực hiện trên cùng một tập dữ liệu công khai. Thực trạng này đặt ra yêu cầu cấp thiết về một nghiên cứu đối chiếu có tính phản biện dựa trên một khung thực nghiệm chuẩn hóa và có khả năng tái lập cao, nội dung này sẽ được trình bày chi tiết trong các mục tiếp theo.

\section{Đối chiếu định lượng}

Bên cạnh sự đa dạng về phương pháp luận, các công trình nghiên cứu sử dụng tập dữ liệu Pima còn ghi nhận \textbf{sự biến thiên đáng kể trong kết quả định lượng}. Mặc dù Accuracy, Precision, Recall, $F_1$-score và ROC--AUC là những chỉ số phổ biến, song \textbf{độ đầy đủ và quy cách trình bày} (như định nghĩa ngưỡng quyết định, phương pháp tính trung bình \textit{macro} hay \textit{weighted}) lại \textbf{thiếu tính nhất quán} giữa các bài báo. Thực trạng này gây khó khăn đáng kể cho việc tổng hợp dữ liệu thành một hệ thống đối chiếu chuẩn hóa.\cite{chang2022pima,kirbas2025shap}

Tổng quan các công trình liên quan cho thấy các mô hình \textbf{tổ hợp} (\textit{ensemble}) và \textbf{thúc đẩy} (\textit{boosting}), mà điển hình là Rừng ngẫu nhiên và XGBoost, thường chiếm ưu thế về độ chính xác trong nhiều thiết lập thực nghiệm. Cụ thể, mô hình RF thường đạt khoảng $82$--$83\%$, trong khi các biến thể boosting có thể chạm ngưỡng $\approx 85\%$ Accuracy và ROC--AUC $\approx 0{,}91$ theo báo cáo của Zhang và cộng sự.\cite{kirbas2025shap,zhang2025pima} Mặt khác, các thuật toán truyền thống như Hồi quy Logistic, SVM hay $k$-NN vẫn duy trì vai trò là các mô hình cơ sở quan trọng, đồng thời sở hữu khả năng \textbf{cạnh tranh hiệu năng mạnh mẽ} nếu được tối ưu hóa siêu tham số trong cùng một quy trình xử lý. Thứ hạng của các mô hình này đặc biệt nhạy cảm với kỹ thuật chuẩn hóa đặc trưng, phương thức xử lý mất cân bằng và chiến lược phân chia dữ liệu.\cite{chang2022pima,sisodia2018,miao2021pima} \textbf{Bảng~\ref{tab:related-numbers}} tóm tắt một số kết quả công bố gần đây nhằm mục đích minh họa, tuy nhiên cần lưu ý rằng các số liệu này không tương đồng về điều kiện thực nghiệm trực tiếp.

\begin{table}[htbp]
  \centering
  \caption{Đối chiếu kết quả công bố trong một số nghiên cứu gần đây trên tập dữ liệu Pima (Lưu ý: Các giá trị chỉ mang tính tham khảo do sự khác biệt về điều kiện thực nghiệm).}
  \label{tab:related-numbers}
  \small
  \setlength{\tabcolsep}{4pt}
  \begin{tabular}{llp{4.5cm}p{4cm}}
    \toprule
    Nguồn & Thuật toán tiêu biểu & Đặc điểm thực nghiệm & Hạn chế đối chiếu \\
    \midrule
    \cite{smith1988} & ADAP & Mô hình nền tảng trên đặc trưng gốc. & Giao thức thực nghiệm sơ khai. \\
    \cite{sisodia2018} & Đa mô hình ML & So sánh truyền thống trên dữ liệu chẩn đoán. & Khác biệt trong tiền xử lý. \\
    \cite{chang2022pima} & Đa mô hình ML & Đánh giá diện rộng các thuật toán. & Phụ thuộc thiết lập cục bộ. \\
    \cite{kirbas2025shap} & RF, XGBoost & Acc $\approx 82$--$83\%$; giải thích SHAP. & Có áp dụng SMOTE. \\
    \cite{zhang2025pima} & XGBoost, RF & Acc $\approx 85\%$; AUC $\approx 0{,}91$. & Thiết lập thực nghiệm riêng biệt. \\
    \bottomrule
  \end{tabular}
\end{table}

Tuy nhiên, các số liệu nêu trên \textbf{không nên được xem là kết quả đối đầu trực tiếp} (\textit{head-to-head}) nếu tách rời khỏi bối cảnh thực nghiệm cụ thể. Hiệu năng của mô hình vốn phụ thuộc mật thiết vào chiến lược tiền xử lý như cách thức xử lý giá trị thiếu hay chuẩn hóa đặc trưng, bên cạnh phương thức phân đoạn dữ liệu, không gian siêu tham số và việc áp dụng các kỹ thuật cân bằng dữ liệu.\cite{chang2022pima,pedregosa2011,kirbas2025shap} Ngay cả khi thực hiện trên cùng một tập dữ liệu, sự biến thiên của các yếu tố này vẫn có thể tạo ra những sai lệch đáng kể giữa các công bố nghiên cứu.

Do đó, việc khẳng định ưu thế của một thuật toán chỉ dựa trên \textbf{chỉ số đơn lẻ} hoặc so sánh số liệu từ \textbf{các nghiên cứu có giao thức khác nhau} dễ dẫn đến những diễn giải thiếu khách quan. Thay vào đó, một cách tiếp cận tin cậy hơn là đặt các mô hình vào \textbf{cùng một điều kiện thực nghiệm}, đồng thời sử dụng \textbf{nhất quán một luồng xử lý} (\textit{pipeline}) và \textbf{chung giao thức đánh giá}. Song song với đó, việc báo cáo \textbf{đa dạng chỉ số} kèm mức độ biến thiên qua các lần lặp là vô cùng cần thiết. Định hướng này chính là nền tảng cho phương pháp nghiên cứu của luận văn (Chương~\ref{sec:methods}--Chương~\ref{sec:results}), xuất phát từ quan điểm cho rằng hiệu năng thực tế là kết quả của toàn bộ cấu trúc pipeline đánh giá.\cite{miao2021pima,pedregosa2011}


\section{Hạn chế tồn tại}

Mặc dù tập dữ liệu \emph{Pima Indians Diabetes} đã được khai thác rộng rãi, việc \textbf{đối chiếu và diễn giải kết quả giữa các công trình} vẫn gặp nhiều trở ngại do sự \textbf{thiếu nhất quán trong thiết kế thực nghiệm và quy cách công bố}. Nhiều nghiên cứu chưa cung cấp đầy đủ thông tin về hạt giống ngẫu nhiên (\textit{random seeds}), phiên bản thư viện hoặc biên độ biến thiên của các chỉ số qua nhiều lần thực thi. Sự thiếu sót này vô hình trung tạo ra rào cản lớn cho việc \textbf{tái lập kết quả và kiểm chứng độc lập}.\cite{chang2022pima,pedregosa2011}

Một vấn đề cốt lõi khác nằm ở sự \textbf{thiếu phân tách tường minh} giữa giai đoạn \emph{lựa chọn mô hình/tinh chỉnh siêu tham số} (\textit{model selection}) và giai đoạn \emph{đánh giá hiệu năng cuối cùng} (\textit{final evaluation}). Trong trường hợp tập kiểm thử bị sử dụng lặp lại để tối ưu hóa tham số, các ước lượng hiệu năng thường có xu hướng \textbf{lạc quan quá mức} (\textit{optimistic bias}), dẫn đến những sai lệch khi phản ánh khả năng tổng quát hóa trên dữ liệu thực tế.\cite{pedregosa2011}

Ngoài ra, quy trình tiền xử lý trong nhiều công trình vẫn \textbf{chưa được đặc tả chi tiết} hoặc chưa tuân thủ nghiêm ngặt nguyên tắc ngăn chặn \emph{rò rỉ dữ liệu} (\textit{data leakage}). Cụ thể là việc tính toán các thông số thống kê để xử lý giá trị thiếu hoặc chuẩn hóa trên \textbf{toàn bộ} tập dữ liệu trước khi phân tách huấn luyện/kiểm thử sẽ vô tình đưa thông tin từ tập kiểm thử vào quá trình học của mô hình, từ đó làm mất đi tính khách quan của các chỉ số đánh giá.\cite{chang2022pima,pedregosa2011}

Bên cạnh đó, việc lựa chọn \textbf{hệ chỉ số chưa thực sự sát với bối cảnh y tế} cũng là một hạn chế đáng kể. Việc quá chú trọng vào độ chính xác (Accuracy) mà xem nhẹ độ nhạy (Recall) hoặc ROC--AUC có thể dẫn đến các kết luận sai lệch trong bài toán \textbf{tầm soát bệnh lý}. Trong kịch bản này, việc bỏ sót ca dương tính (\textit{false negative}) thường gây ra hậu quả nghiêm trọng hơn so với các báo động giả (\textit{false positive}), vốn là một yếu tố then chốt cần được ưu tiên cân nhắc khi đánh giá hiệu năng trên dữ liệu Pima.\cite{chang2022pima,kirbas2025shap}

Những hạn chế nêu trên không chỉ làm suy giảm \textbf{giá trị khoa học nội tại} (\textit{internal validity}) của từng công trình mà còn cản trở \textbf{khả năng tổng hợp và đối chiếu khách quan} giữa các nghiên cứu. Thực trạng này càng củng cố yêu cầu cấp thiết về một phương pháp tiếp cận \textbf{hệ thống, minh bạch và có khả năng tái lập}. Các nguyên tắc này sẽ được hiện thực hóa trong quy trình thực nghiệm tại Chương~\ref{sec:methods}, đồng thời được kiểm chứng thông qua các phân tích về tính giá trị (\textit{validity}) tại Chương~\ref{sec:results} và Chương~\ref{sec:discussion}.\cite{pedregosa2011}

\section{Mục tiêu cải tiến}

Tổng hợp từ các phân tích về bối cảnh, dữ liệu định lượng và những nhận định mang tính phản biện, có thể khẳng định rằng dù tập dữ liệu \emph{Pima Indians Diabetes} đã được khai thác rộng rãi, vẫn tồn tại những \textbf{khoảng trống nghiêm trọng} trong thiết kế thực nghiệm và quy cách công bố kết quả. Nhiều công trình hiện nay ghi nhận mức hiệu năng có phần \textbf{lạc quan quá mức}, vốn là hệ quả của giao thức đánh giá thiếu nhất quán hoặc hiện tượng rò rỉ dữ liệu (\textit{data leakage}) tiềm ẩn. Thực trạng này làm suy giảm tính tin cậy của các kết luận về hiệu năng \emph{tương đối} giữa các thuật toán khi thực hiện đối chiếu liên nghiên cứu.\cite{chang2022pima,pedregosa2011}

Cụ thể, các hạn chế trọng tâm bao gồm:

\begin{itemize}
    \item \textbf{Thiếu quy trình chuẩn hóa}: Sự vắng mặt của một hệ quy chiếu chung về tiền xử lý, không gian siêu tham số và phân đoạn dữ liệu khiến các kết quả công bố \textbf{khó đối chiếu trực tiếp}, từ đó làm hạn chế khả năng kế thừa cũng như tổng hợp tri thức từ các nghiên cứu đi trước.\cite{chang2022pima,miao2021pima}
    \item \textbf{Thách thức về tính tái lập (\textit{Reproducibility})}: Việc thiếu đặc tả các thiết lập kỹ thuật như hạt giống ngẫu nhiên hay phiên bản thư viện, kết hợp với việc không báo cáo \textbf{biên độ biến thiên} của chỉ số qua nhiều lần thực thi khiến kết quả dễ bị nhiễu bởi yếu tố ngẫu nhiên trong phân chia dữ liệu.\cite{pedregosa2011}
    \item \textbf{Sự chồng lấn giữa lựa chọn và đánh giá mô hình}: Nguyên tắc \textbf{tách biệt tường minh giữa tối ưu hóa và đánh giá độc lập} chưa được tuân thủ nghiêm ngặt. Trong trường hợp dữ liệu kiểm thử tham gia trực tiếp hoặc gián tiếp vào quá trình tinh chỉnh, mô hình sẽ bị thiên kiến và không phản ánh đúng năng lực tổng quát hóa thực tế.\cite{pedregosa2011}
    \item \textbf{Hệ chỉ số chưa sát thực tiễn lâm sàng}: Việc lạm dụng độ chính xác (Accuracy) trong khi thiếu vắng các chỉ số đo lường sai số loại I và loại II (\textit{False Positive/Negative}) có thể gây ra những diễn giải sai lệch trong bài toán \textbf{tầm soát y tế}, nơi mà chi phí của việc bỏ lỡ ca bệnh là vô cùng lớn.\cite{chang2022pima,kirbas2025shap}
\end{itemize}

\textbf{Để khắc phục các hạn chế nêu trên, nghiên cứu này} tập trung xây dựng một \textbf{khung thực nghiệm có kiểm soát chặt chẽ và khả năng tái lập cao}. Các mô hình truyền thống được đặt trong \textbf{cùng một điều kiện thực nghiệm} với luồng xử lý (\textit{pipeline}) nhất quán nhằm triệt tiêu rò rỉ dữ liệu, đồng thời phân định rạch ròi giữa \textbf{vòng lặp tối ưu nội bộ} và \textbf{đánh giá độc lập} qua nhiều hạt giống ngẫu nhiên. Nghiên cứu báo cáo \textbf{hệ chỉ số toàn diện} kèm thống kê mô tả, áp dụng \textbf{kiểm định thống kê} để xác định ý nghĩa của các khác biệt hiệu năng, đồng thời thực hiện \textbf{phân tích sai số} (\textit{error analysis}) để làm rõ giới hạn của mô hình. Chi tiết triển khai được trình bày tại Chương~\ref{sec:methods}, trong khi kết quả và các thảo luận về tính giá trị sẽ được phân tích sâu tại Chương~\ref{sec:results} và Chương~\ref{sec:discussion}.\cite{pedregosa2011}

\chapter{Phương pháp nghiên cứu}
\label{sec:methods}

\section{Quy trình thực nghiệm}

Nghiên cứu này xây dựng một \textbf{quy trình học máy toàn diện} nhằm thiết lập, ước lượng và đối chiếu các mô hình phân loại nhị phân trên dữ liệu lâm sàng dạng bảng. Quy trình được thiết kế với mục tiêu tối ưu hóa tính \textbf{nhất quán} giữa các thuật toán, đồng thời đảm bảo \textbf{khả năng tái lập} và nâng cao \textbf{độ tin cậy} cho các kết quả đánh giá độc lập (\textit{out-of-sample evaluation}).

Luồng xử lý (\textit{pipeline}) được cấu trúc theo trình tự tuyến tính từ giai đoạn \textbf{Dữ liệu thô, Tiền xử lý, Huấn luyện mô hình cho đến Đánh giá} (Hình~\ref{fig:pipeline}). Tại giai đoạn khởi đầu, dữ liệu được làm sạch, xử lý giá trị thiếu và chuẩn hóa đặc trưng. Các thao tác này được thực hiện nghiêm ngặt \emph{trong phạm vi tập huấn luyện} hoặc từng \textit{fold} huấn luyện nhằm triệt tiêu hoàn toàn hiện tượng rò rỉ dữ liệu (\textit{data leakage}). Tiếp theo, một hệ quy chiếu thực nghiệm thống nhất được áp dụng cho các mô hình đại diện cho những \textbf{giả định học tập (\textit{learning paradigms}) khác nhau}, bao gồm mô hình tuyến tính (Hồi quy Logistic), tối ưu hóa lề (SVM), dựa trên khoảng cách (KNN) và tổ hợp cây quyết định (Rừng ngẫu nhiên). Sau giai đoạn huấn luyện, hiệu năng được đo lường đa chiều thông qua hệ chỉ số Accuracy, Precision, Recall, $F_1$-score, ROC--AUC và \textbf{Brier score} nhằm đánh giá độ hiệu chuẩn xác suất. Cách tiếp cận này giúp phản ánh toàn diện các khía cạnh của bài toán, đặc biệt là trong bối cảnh \textbf{tầm soát y tế}, nơi mà sai lầm loại II (\textit{false negative}) thường dẫn đến hệ quả nghiêm trọng hơn so với sai lầm loại I (\textit{false positive}).

Quy trình thực nghiệm đảm bảo sự \textbf{tách biệt hoàn toàn} giữa tập huấn luyện và tập kiểm thử, đồng thời phân lập rạch ròi \textbf{vòng lặp tối ưu siêu tham số} khỏi giai đoạn \textbf{đánh giá hiệu năng cuối cùng} (Mục~\ref{subsec:metrics}). Trọng tâm của nghiên cứu không chỉ dừng lại ở việc tối ưu hóa một chỉ số đơn lẻ, mà còn hướng tới tính \textbf{minh bạch}, khả năng \textbf{diễn giải} cũng như tính \textbf{tái lập} của các mô hình dự báo.

\begin{figure}[htbp]
  \centering
  \begin{tikzpicture}[
    node distance=6mm and 10mm,
    box/.style={rectangle, draw, rounded corners=3pt, minimum height=10mm, minimum width=25mm, align=center, font=\small\sffamily, fill=gray!5},
    arr/.style={-{Stealth[length=2.5mm]}, thick, draw=black!70}
  ]
    \node[box] (data) {Dữ liệu thô};
    \node[box, right=of data] (pre) {Tiền xử lý};
    \node[box, right=of pre] (fs) {Kỹ thuật\\đặc trưng};
    \node[box, right=of fs] (m) {Huấn luyện\\mô hình};
    \node[box, right=of m] (e) {Đánh giá\\hiệu năng};
    
    \draw[arr] (data) -- (pre);
    \draw[arr] (pre) -- (fs);
    \draw[arr] (fs) -- (m);
    \draw[arr] (m) -- (e);
  \end{tikzpicture}
  \caption{Sơ đồ quy trình nghiên cứu đề xuất. Toàn bộ chuỗi xử lý được lặp lại với 05 hạt giống ngẫu nhiên (\textit{seeds}) đi kèm giao thức xác thực chéo nội bộ cố định (Mục~\ref{subsec:protocol}).}
  \label{fig:pipeline}
\end{figure}

\section{Mô tả tập dữ liệu}
\label{subsec:dataset}

Nghiên cứu này sử dụng bộ dữ liệu \textbf{Pima Indians Diabetes} làm cơ sở để xây dựng và đánh giá các mô hình dự báo nguy cơ đái tháo đường từ các đặc trưng lâm sàng dạng bảng. Đây vốn là tập dữ liệu \textbf{chuẩn tham chiếu} (\textit{benchmark}) điển hình trong lĩnh vực học máy y sinh, tạo điều kiện thuận lợi cho việc đối chiếu khách quan các kết quả thực nghiệm với hệ thống y văn hiện hành.\cite{smith1988,uci-pima}

Tập dữ liệu bao gồm \textbf{768} mẫu, trong đó theo mô tả gốc, các quan sát này được thu thập từ đối tượng là \textbf{phụ nữ} thuộc cộng đồng người Pima tại Hoa Kỳ.\cite{smith1988} Mỗi mẫu dữ liệu được đặc trưng bởi \textbf{08 biến định lượng} bao gồm số lần mang thai, nồng độ glucose huyết tương, huyết áp tâm trương, độ dày nếp gấp da, chỉ số insulin huyết thanh, chỉ số khối cơ thể (BMI), chức năng phả hệ đái tháo đường (\textit{diabetes pedigree function}) và tuổi. Bên cạnh đó, \textbf{biến mục tiêu} là nhãn nhị phân biểu thị trạng thái có hoặc không mắc bệnh đái tháo đường dựa trên tiêu chí chẩn đoán của nghiên cứu gốc.

Về \textbf{phân bố lớp}, tập dữ liệu ghi nhận tình trạng \textbf{mất cân bằng nhẹ} khi lớp âm tính chiếm \textbf{500} mẫu ($\approx 65,1\%$), còn lớp dương tính chiếm \textbf{268} mẫu ($\approx 34,9\%$). Đặc tính này chính là cơ sở then chốt để lựa chọn hệ chỉ số đánh giá, bởi lẽ việc sử dụng \textbf{Accuracy} có thể dẫn đến những nhận định sai lệch về hiệu năng thực tế trên lớp thiểu số. Do đó, nghiên cứu ưu tiên sử dụng các chỉ số thay thế như Precision, Recall, $F_1$-score và ROC--AUC (Mục~\ref{subsec:metrics}).

Một đặc điểm kỹ thuật quan trọng cần lưu ý là sự xuất hiện của giá trị \textbf{0} tại một số biến sinh lý như Glucose, Huyết áp, Độ dày da, Insulin và BMI. Đây vốn là những giá trị \textbf{không hợp lệ về mặt lâm sàng} và được xác định là các \textbf{giá trị thiếu đã qua mã hóa}. Quy trình xử lý hệ thống cho các trường hợp này sẽ được trình bày chi tiết tại Mục~\ref{subsec:preprocess}.

\textbf{Lý do lựa chọn tập dữ liệu Pima} tập trung vào hai khía cạnh trọng tâm. \textbf{Thứ nhất}, tính \textbf{công khai} và phổ biến của bộ dữ liệu này giúp đảm bảo tính \textbf{minh bạch} cũng như khả năng \textbf{tái lập} quy trình thực nghiệm.\cite{uci-pima,pedregosa2011} \textbf{Thứ hai}, mặc dù quy mô mẫu chỉ ở mức trung bình, bài toán vẫn bảo toàn nguyên vẹn giá trị phương pháp luận và được xem là mô hình lý tưởng để minh chứng cho các \textbf{giao thức thực nghiệm} chặt chẽ trên dữ liệu dạng bảng.

\textbf{Các hạn chế của dữ liệu}: Cỡ mẫu hạn chế có thể làm tăng phương sai ước lượng và rủi ro quá khớp (\textit{overfitting}) trong quá trình tối ưu hóa mô hình. Đồng thời, do dữ liệu chỉ đại diện cho một quần thể dân tộc đặc thù, các kết luận nghiên cứu \textbf{không} hướng tới việc tổng quát hóa trực tiếp cho các cộng đồng khác nếu chưa có bước kiểm định độc lập (\textit{external validation}). Chính vì vậy, các kết quả thực nghiệm cần được diễn giải một cách thận trọng trong phạm vi các giới hạn này (chi tiết tại Chương~\ref{sec:discussion}).\cite{uci-pima,smith1988}


\section{Tiền xử lý dữ liệu}
\label{subsec:preprocess}

Quy trình tiền xử lý được thiết lập nhằm \textbf{chuẩn hóa chất lượng dữ liệu đầu vào}, qua đó giảm thiểu nhiễu từ các giá trị thiếu cũng như sự khác biệt về thang đo giữa các biến. Đồng thời, việc cố định \textbf{trình tự các bước xử lý} giúp đảm bảo mọi mô hình đều được đối chiếu trên một hệ quy chiếu đồng nhất. Quy trình này bao gồm ba giai đoạn logic trọng tâm, bao gồm: \textbf{xử lý giá trị thiếu}, \textbf{chuẩn hóa đặc trưng}, và \textbf{phân đoạn dữ liệu kết hợp xác thực chéo nội bộ} nhằm tối ưu hóa siêu tham số. Toàn bộ các thao tác này đều tuân thủ nghiêm ngặt nguyên tắc \textbf{ngăn chặn rò rỉ dữ liệu} (\textit{data leakage}), vốn là yếu tố then chốt để đảm bảo sự tách biệt tường minh giữa giai đoạn lựa chọn mô hình và đánh giá hiệu năng cuối cùng (Mục~\ref{subsec:metrics}).\cite{pedregosa2011}

\subsection{Xử lý giá trị thiếu}

Như đã phân tích tại Mục~\ref{subsec:dataset}, các giá trị \textbf{0} không hợp lệ về mặt sinh lý tại các biến như Glucose, Huyết áp, Độ dày da, Insulin và BMI sẽ được chuyển đổi thành \textbf{giá trị thiếu} (\texttt{NaN}) trước khi thực hiện các bước xử lý tiếp theo.

Nghiên cứu áp dụng phương pháp \textbf{gán giá trị thiếu bằng trung vị} (\textit{median imputation}). Giá trị này được ước lượng \emph{duy nhất trên tập huấn luyện} hoặc phần huấn luyện của từng \textit{fold} trong quá trình xác thực chéo, sau đó được dùng để thay thế các vị trí trống trên tập kiểm thử hoặc các \textit{fold} kiểm định tương ứng. So với số trung bình, trung vị sở hữu ưu điểm là \textbf{ít nhạy cảm với các giá trị ngoại lai} (\textit{outliers}), vốn là một đặc điểm phổ biến trong dữ liệu y sinh, từ đó giúp duy trì tính ổn định cho phân phối thực nghiệm của các đặc trưng. Quy trình này được triển khai thông qua công cụ \texttt{SimpleImputer(strategy='median')} tích hợp trong một \texttt{Pipeline} thống nhất nhằm đảm bảo tính tự động hóa và ngăn ngừa triệt để các sai sót thực nghiệm.

\subsection{Chuẩn hóa đặc trưng}

Do các biến số lâm sàng sở hữu \textbf{đơn vị và thang đo khác biệt}, chẳng hạn như sự chênh lệch giữa chỉ số Glucose so với BMI hay Tuổi, các mô hình dựa trên \textbf{khoảng cách} hoặc \textbf{tối ưu hóa lề} như KNN và SVM, cùng với Hồi quy Logistic có điều chuẩn $L_2$, đều trở nên rất \textbf{nhạy cảm} với biên độ của dữ liệu. Do đó, nghiên cứu áp dụng kỹ thuật \textbf{chuẩn hóa Z-score} (\texttt{StandardScaler}) nhằm đưa các đặc trưng về phân phối có kỳ vọng bằng 0 và độ lệch chuẩn bằng 1. Các tham số của bộ chuẩn hóa được ước lượng \emph{duy nhất từ tập huấn luyện}, sau đó mới thực hiện phép biến đổi (\texttt{transform}) lên tập kiểm thử. Phương pháp này giúp loại bỏ triệt để rủi ro rò rỉ thông tin từ dữ liệu chưa biết vào quá trình huấn luyện.\cite{pedregosa2011}

Đối với mô hình Rừng ngẫu nhiên, mặc dù việc co giãn đặc trưng không ảnh hưởng đến bản chất của các phép phân tách dựa trên ngưỡng, nghiên cứu vẫn duy trì \textbf{nhất quán một luồng xử lý} (\textit{imputer + scaler}) cho tất cả các mô hình. Việc này nhằm đảm bảo \textbf{tính đồng nhất của giao thức thực nghiệm}, đồng thời tạo điều kiện đối chiếu công bằng giữa các thuật toán.

\section{Phân đoạn dữ liệu và Xác thực chéo nội bộ (\textit{Data Splitting \& Inner CV})}

Để tối ưu hóa lượng dữ liệu huấn luyện trong khi vẫn duy trì được một tập mẫu \textbf{độc lập} cho giai đoạn đánh giá cuối cùng, nghiên cứu thực hiện phân tách dữ liệu theo phương pháp \textbf{hold-out} với tỷ lệ \textbf{80/20}. Quá trình này áp dụng kỹ thuật \textbf{phân tầng} (\textit{stratification}) nhằm bảo toàn nguyên vẹn tỷ lệ phân bố lớp giữa các tập con. Trong đó, \textbf{tập kiểm thử} (\textit{test set}) được giữ hoàn toàn biệt lập và tuyệt đối không tham gia vào bất kỳ giai đoạn nào từ gán giá trị thiếu, chuẩn hóa cho đến tinh chỉnh siêu tham số. Mọi quyết định tối ưu hóa theo đó chỉ được thực hiện dựa trên dữ liệu thuộc \textbf{tập huấn luyện}.\cite{pedregosa2011}

Trong phạm vi tập huấn luyện, quy trình tinh chỉnh siêu tham số sử dụng \textbf{xác thực chéo $k$-fold phân tầng} với $k=5$, có kết hợp xáo trộn dữ liệu và hạt giống cố định. Phương pháp này giúp khai thác tối đa nguồn dữ liệu hạn chế, đồng thời giảm thiểu \textbf{phương sai sai số ước lượng} so với việc chỉ phân chia tập dữ liệu đơn lẻ. Sau khi xác định được cấu hình tối ưu, mô hình sẽ được huấn luyện lại trên toàn bộ tập huấn luyện trước khi tiến hành đo lường hiệu năng thực tế trên tập kiểm thử (Mục~\ref{subsec:models}).

\textbf{Chiến lược kiểm soát rò rỉ dữ liệu}: Mọi bước ước lượng tham số từ dữ liệu, bao gồm tính trung vị để điền giá trị thiếu, xác định kỳ vọng và độ lệch chuẩn để chuẩn hóa cho đến quá trình tìm kiếm siêu tham số, đều chỉ được thực hiện (\texttt{fit}) trong phạm vi tập huấn luyện. Nghiên cứu tuyệt đối \textbf{không} tính toán thống kê trên toàn bộ tập dữ liệu trước khi phân tách nhằm đảm bảo tính khách quan cũng như khả năng tổng quát hóa của kết quả đánh giá.

Quy trình triển khai thực tế được cụ thể hóa qua 05 bước chính, bắt đầu từ việc phân tách tập huấn luyện và kiểm thử có phân tầng, sau đó thiết lập các bộ tiền xử lý như \textit{Imputer} và \textit{Scaler} trong luồng xử lý (\textit{pipeline}) dựa trên tập huấn luyện. Tiếp theo, quá trình tối ưu hóa siêu tham số được thực hiện thông qua xác thực chéo nội bộ, làm cơ sở để tái huấn luyện mô hình với cấu hình tối ưu, trước khi tiến hành đánh giá hiệu năng trên tập kiểm thử độc lập (Mục~\ref{subsec:protocol}).

\section{Lựa chọn đặc trưng}
\label{subsec:features}

Trong nghiên cứu này, các kỹ thuật \textbf{giảm chiều} và \textbf{lựa chọn đặc trưng tự động} chủ đích \textbf{không được áp dụng}. Quyết định này nhằm đảm bảo sự cân bằng giữa hiệu năng dự báo, \textbf{khả năng diễn giải} (\textit{interpretability}) cùng tính \textbf{tối giản của giao thức} đối chiếu giữa các thuật toán.

\textbf{Thứ nhất}, không gian đầu vào chỉ bao gồm \textbf{08 biến định lượng} (Mục~\ref{subsec:dataset}), do đó bài toán không gặp phải hiện tượng ``lời nguyền đa chiều'' (\textit{curse of dimensionality}). Các phương pháp như PCA hoặc lọc tương quan trước huấn luyện có thể không mang lại lợi ích hiệu năng rõ rệt, trái lại còn có nguy cơ làm mất đi những thông tin hữu ích cho quá trình phân tích.

\textbf{Thứ hai}, mỗi biến số trong tập dữ liệu đều mang \textbf{ý nghĩa lâm sàng trực tiếp}. Việc áp dụng các phép chiếu tuyến tính như PCA sẽ tạo ra các thành phần tổng hợp khó ánh xạ ngược về đơn vị đo y khoa, từ đó gây \textbf{suy giảm tính minh bạch}, vốn là một yếu tố then chốt trong các ứng dụng hỗ trợ quyết định y tế. Việc bảo toàn không gian đặc trưng gốc giúp các chỉ số như hệ số hồi quy, tầm quan trọng đặc trưng trong RF và tầm quan trọng hoán vị (\textit{permutation importance}) duy trì được mối liên kết trực tiếp với các biến số thực tế (Mục~\ref{subsec:interpret}).

\textbf{Thứ ba}, các mô hình tổ hợp như \textbf{Rừng ngẫu nhiên} vốn đã sở hữu cơ chế tự khai thác những \textbf{tương tác phi tuyến}, đồng thời cung cấp thứ hạng đặc trưng nội tại. Do đó, nhu cầu về một giai đoạn lựa chọn đặc trưng độc lập là không thực sự cấp thiết đối với một bài toán có số lượng biến hạn chế như nghiên cứu này.

Tóm lại, nghiên cứu ưu tiên tính tinh gọn và độ tin cậy của giao thức thực nghiệm thay vì tìm cách tối đa hóa chỉ số thông qua các phép biến đổi đặc trưng phức tạp. Mọi phân tích về mức độ đóng góp của biến số nhằm giải quyết câu hỏi nghiên cứu RQ2 đều được thực hiện gián tiếp dựa trên kết quả từ các mô hình đã huấn luyện trong luồng xử lý chuẩn hóa.\cite{pedregosa2011}

\section{Mô hình hóa quy trình huấn luyện}
\label{subsec:models}
\label{subsec:training}

Nghiên cứu thiết lập một \textbf{hệ thống các bộ phân loại} đại diện cho những hướng tiếp cận khác nhau trong việc xử lý dữ liệu y sinh dạng bảng. Mục tiêu cốt lõi của việc này không chỉ dừng lại ở việc xếp hạng hiệu năng định lượng, mà còn nhằm làm rõ \textbf{ảnh hưởng của thiên kiến cảm dẫn} (\textit{inductive bias}) đối với kết quả dự báo trong cùng một luồng xử lý (\textit{pipeline}) nhất quán. Điều này giúp giải thích các giả định nội tại về hình dạng ranh giới quyết định cũng như cấu trúc phụ thuộc giữa các biến số (Mục~\ref{subsec:preprocess}, \ref{subsec:metrics}).

Cụ thể, bốn thuật toán được đưa vào thực nghiệm đối chiếu bao gồm: \textbf{Hồi quy Logistic} (Logistic Regression), \textbf{$k$-láng giềng gần nhất} (KNN), \textbf{Máy vectơ hỗ trợ} (SVM) và \textbf{Rừng ngẫu nhiên} (Random Forest). Hệ thống này bao phủ phổ giải pháp từ \textbf{tuyến tính} đến \textbf{phi tuyến}, đồng thời đi từ phương thức quyết định \textbf{dựa trên khoảng cách} đến \textbf{tổ hợp cây quyết định}. Cách tiếp cận này giúp khảo sát toàn diện không gian lời giải trên tập dữ liệu Pima trong phạm vi các mô hình truyền thống có tính tái lập cao.\cite{pedregosa2011,sisodia2018}

\begin{itemize}
  \item \textbf{Hồi quy Logistic} đóng vai trò là \textbf{mô hình cơ sở tuyến tính} (\textit{linear baseline}) dựa trên giả định rằng log-odds của lớp dương tính là một hàm tuyến tính của các biến đặc trưng. Ưu điểm nổi bật của mô hình này nằm ở sự tinh gọn cùng các \textbf{hệ số hồi quy có khả năng diễn giải trực tiếp} về mặt xác suất (Mục~\ref{subsec:interpret}).
  
  \item \textbf{KNN} là phương pháp \textbf{học dựa trên thực thể} (\textit{instance-based}) và không thực hiện ước lượng một hàm tham số cố định trên toàn bộ không gian mẫu. Thay vào đó, nhãn của dữ liệu mới sẽ được suy diễn từ các láng gềng gần nhất dựa trên hàm đo khoảng cách. Tuy KNN sở hữu khả năng nắm bắt các \textbf{cấu trúc phi tuyến cục bộ}, song thuật toán này lại đặc biệt \textbf{nhạy cảm} với thang đo đặc trưng cũng như nhiễu dữ liệu.
  
  \item \textbf{SVM} thuộc họ \textbf{tối ưu hóa lề} (\textit{margin-based}) với mục tiêu tìm kiếm siêu phẳng phân tách sao cho khoảng cách lề (\textit{margin}) giữa các lớp đạt giá trị lớn nhất. Thông qua việc sử dụng \textbf{nhân RBF} (\textit{Radial Basis Function}), SVM có khả năng thiết lập nên các ranh giới quyết định phi tuyến phức tạp ngay trong không gian đặc trưng gốc. 
  
  \item \textbf{Rừng ngẫu nhiên} là mô hình \textbf{tổ hợp} (\textit{ensemble}) dựa trên nhiều cây quyết định thông qua kỹ thuật lấy mẫu \textit{bootstrap} và lựa chọn tập con đặc trưng ngẫu nhiên. Phương pháp này giúp \textbf{giảm phương sai} sai số dự báo, đồng thời tỏ ra hiệu quả trong việc khai thác các \textbf{tương tác phi tuyến} giữa các biến số và cung cấp các ước lượng tin cậy về tầm quan trọng đặc trưng.
\end{itemize}

Ngoài ra, nghiên cứu sử dụng \texttt{DummyClassifier} với chiến lược \textbf{lớp đa số} (Majority) nhằm thiết lập một \textbf{đường cơ sở tham chiếu bổ sung}. Mô hình này được tích hợp trực tiếp vào luồng xử lý để xác định ngưỡng hiệu năng tối thiểu, từ đó làm căn cứ đánh giá ý nghĩa thực tiễn của việc áp dụng các thuật toán học máy. Toàn bộ thực nghiệm trong nghiên cứu đều được triển khai dựa trên nền tảng thư viện \textbf{scikit-learn}.\cite{pedregosa2011}

Quy trình huấn luyện được thiết kế chặt chẽ nhằm đảm bảo \textbf{tính khách quan tuyệt đối}, đồng thời hạn chế tối đa hiện tượng \textbf{quá khớp trong giai đoạn tinh chỉnh} (\textit{tuning-overfitting}). Quy trình này bao gồm hai giai đoạn logic trọng tâm, trong đó giai đoạn đầu tiên tập trung vào việc \textbf{tối ưu hóa siêu tham số} dựa trên không gian lưới tham số (Phụ lục~\ref{app:hyper}) và chỉ thực hiện trong phạm vi tập huấn luyện. Giai đoạn tiếp theo thực hiện \textbf{tái huấn luyện mô hình} với cấu hình tối ưu trên toàn bộ dữ liệu huấn luyện, làm tiền đề trước khi tiến hành đánh giá độc lập trên tập \textbf{hold-out} (Mục~\ref{subsec:metrics}).

\section{Tối ưu hóa siêu tham số }

Đối với mỗi mô hình, không gian siêu tham số cần tối ưu hóa được thiết lập nhằm phản ánh chính xác đặc thù vận hành của từng thuật toán. Các thông số này bao gồm \textbf{cường độ điều chuẩn} trong Hồi quy Logistic, \textbf{số lượng láng giềng $k$} của KNN, các hệ số \textbf{\texttt{C}} và \textbf{\texttt{gamma}} đối với SVM nhân RBF, cùng với \textbf{số lượng cây} và \textbf{độ sâu tối đa} trong mô hình Rừng ngẫu nhiên. Các giá trị cụ thể được xác định dựa trên lưới tham số cố định được trình bày chi tiết tại Phụ lục~\ref{app:hyper} cũng như trong mã nguồn đính kèm.

Quá trình tối ưu hóa sử dụng phương pháp \textbf{tìm kiếm lưới} (\texttt{GridSearchCV}) kết hợp với \textbf{xác thực chéo $k$-fold phân tầng} ($k=5$, sử dụng \texttt{StratifiedKFold} có thực hiện xáo trộn dữ liệu và cố định hạt giống ngẫu nhiên cho vòng lặp nội bộ). Quy trình này chỉ được triển khai trên \emph{phần dữ liệu huấn luyện} sau khi đã tách biệt tập kiểm thử độc lập theo tỷ lệ 80/20. Trong mỗi \textit{fold}, luồng xử lý (\textit{pipeline}) được thiết lập (\texttt{fit}) trên phần huấn luyện, sau đó mới thực hiện đo lường trên phần kiểm định tương ứng. Với mỗi tổ hợp siêu tham số, \textbf{điểm trung bình} trên các \textit{fold} sẽ được dùng làm căn cứ để lựa chọn cấu hình tối ưu, qua đó giúp giảm thiểu phương sai ước lượng so với các phương pháp phân chia dữ liệu đơn lẻ.\cite{pedregosa2011}

Tiêu chí tối ưu hóa chủ đạo được lựa chọn là điểm \textbf{$F_1$} (\texttt{scoring='f1'}) nhằm đảm bảo sự cân bằng giữa Precision và Recall trong bối cảnh dữ liệu có hiện tượng mất cân bằng lớp nhẹ. Việc lựa chọn này đồng thời giúp phản ánh đúng mục tiêu ưu tiên khả năng nhận diện lớp dương tính trong bài toán dự báo lâm sàng.

\section{Tái huấn luyện mô hình tối ưu}

Sau khi xác định được cấu hình siêu tham số tối ưu thông qua điểm xác thực chéo trung bình, bước tiếp theo là thực hiện \textbf{tái huấn luyện} toàn bộ luồng xử lý (\texttt{Pipeline}). Quy trình này bao gồm các công đoạn từ gán giá trị thiếu (\textit{imputer}), chuẩn hóa đặc trưng (\textit{scaler}) cho đến vận hành bộ phân loại trên \textbf{toàn bộ} dữ liệu thuộc tập huấn luyện (80\%). Phương pháp này cho phép mô hình khai thác tối đa lượng dữ liệu khả dụng trước khi tiến hành đánh giá hiệu năng trên \textbf{tập kiểm thử} độc lập (20\%), vốn là tập dữ liệu được giữ biệt lập hoàn toàn khỏi mọi công đoạn tinh chỉnh trước đó.

\textbf{Kiểm soát rò rỉ dữ liệu}: Quy trình gán giá trị thiếu và chuẩn hóa đặc trưng được tích hợp chặt chẽ trong cùng một luồng xử lý (\texttt{Pipeline}). Các thông số của giai đoạn tiền xử lý này chỉ được ước lượng (\texttt{fit}) dựa trên dữ liệu huấn luyện, cụ thể là trong từng \textit{fold} huấn luyện khi thực hiện xác thực chéo và thực hiện thêm một lần nữa trên toàn bộ tập huấn luyện trước khi dự báo trên tập kiểm thử. Nghiên cứu đảm bảo \textbf{tuyệt đối không} sử dụng bất kỳ đặc trưng thống kê hay nhãn mục tiêu nào từ tập kiểm thử để điều chỉnh siêu tham số hoặc các bước tiền xử lý.

Việc cố định hạt giống ngẫu nhiên cho các vòng lặp nội bộ và các mô hình có tính ngẫu nhiên như Rừng ngẫu nhiên, kết hợp với giao thức lặp lại trên nhiều lần phân đoạn dữ liệu, đã được tóm tắt chi tiết tại Mục~\ref{subsec:repro} và Mục~\ref{subsec:protocol}. Toàn bộ các thông số thực nghiệm cùng phiên bản thư viện đều được đồng bộ hóa nhằm đảm bảo \textbf{khả năng tái lập} hoàn toàn cho nghiên cứu.

\section{Giao thức đánh giá}
\label{subsec:metrics}

\noindent\textbf{Hệ thống chỉ số đánh giá.} Để thực hiện đánh giá hiệu năng \textbf{đa chiều}, nghiên cứu báo cáo đồng thời các chỉ số bao gồm \textbf{Accuracy} (Độ chính xác tổng thể), \textbf{Precision} (Độ chính xác dương tính), \textbf{Recall} (Độ nhạy), \textbf{F1-score}, \textbf{ROC--AUC} và \textbf{PR--AUC} (Diện tích dưới đường cong Precision--Recall) trên tập kiểm thử độc lập (\textit{hold-out}). Các chỉ số này được tổng hợp và tính toán giá trị kỳ vọng thông qua quy trình lặp lại trên nhiều hạt giống phân đoạn dữ liệu (Mục~\ref{subsec:protocol}). Cách tiếp cận này giúp phản ánh toàn diện các khía cạnh của bài toán phân loại nhị phân, đặc biệt là trong bối cảnh dữ liệu \textbf{mất cân bằng lớp} (Mục~\ref{subsec:dataset}).

\textbf{Accuracy} đo lường tỷ lệ dự báo đúng trên tổng số mẫu. Chỉ số này cung cấp cái nhìn sơ bộ về hiệu năng tổng quát nhưng lại có nguy cơ \textbf{che khuất} sự hạn chế của mô hình trên lớp thiểu số. Cụ thể là một mô hình thiên kiến về lớp đa số vẫn có thể đạt Accuracy cao dù thực tế đã \textbf{bỏ sót} hầu hết các ca dương tính.

\textbf{Precision} biểu thị tỷ lệ chính xác trong số các trường hợp được dự báo là dương tính, trong khi \textbf{Recall} phản ánh năng lực \textbf{phát hiện} các ca bệnh thực tế. Trong bối cảnh \textbf{tầm soát và tiên lượng nguy cơ}, sai lầm loại II vốn là việc bỏ sót ca bệnh thường dẫn đến hệ quả lâm sàng nghiêm trọng hơn so với sai lầm loại I liên quan đến các báo động giả. Do đó, Recall và các chỉ số AUC được xem là các tiêu chí \textbf{quan trọng tương đương Accuracy}. Đồng thời, chỉ số Precision vẫn được duy trì nhằm đánh giá áp lực lên hệ thống y tế phát sinh từ các báo động giả.

\textbf{F1-score} là \textbf{trung bình điều hòa} (\textit{harmonic mean}) của Precision và Recall, vốn có vai trò giúp tối ưu hóa \textbf{đồng thời} cả hai loại sai số. Đối với dữ liệu mất cân bằng, F1-score được xem là chỉ số tổng hợp có độ tin cậy cao hơn hẳn so với Accuracy đơn thuần. Đây cũng chính là \textbf{tiêu chí mục tiêu} (\texttt{scoring='f1'}) trong quy trình \texttt{GridSearchCV} (Mục~\ref{subsec:models}), đóng vai trò làm \textbf{điểm cân bằng có trọng số} khi lựa chọn cấu hình mô hình tối ưu.

\textbf{ROC--AUC} đánh giá khả năng phân hạng (\textit{ranking capability}) của mô hình trên \textbf{mọi ngưỡng cắt}, từ đó giúp so sánh hiệu năng tổng quát thay vì chỉ giới hạn tại một điểm cố định. Chỉ số này bổ trợ đắc lực cho các phép đo tại ngưỡng mặc định, tuy nhiên lại \textbf{không} thể thay thế hoàn toàn việc phân tích Precision--Recall trong trường hợp tỷ lệ lớp có độ lệch lớn hoặc khi chi phí sai số mang tính không đối xứng.

\textbf{Precision--Recall và PR--AUC.} Bên cạnh ROC--AUC, nghiên cứu còn sử dụng \textbf{đường cong Precision--Recall} cùng \textbf{diện tích dưới đường cong} tương ứng bao gồm PR--AUC hoặc \textit{Average Precision}. Trong điều kiện lớp dương tính là thiểu số hoặc trở thành \textbf{đối tượng ưu tiên}, PR--AUC thường tỏ ra \textbf{nhạy} hơn ROC--AUC trong việc phản ánh sự thay đổi hiệu năng cũng như sự đánh đổi giữa độ chính xác và độ nhạy, qua đó đáp ứng tốt mục tiêu của công tác sàng lọc y tế.

\textbf{Tính hiệu chuẩn và độ tin cậy.} Ngoài các chỉ số phân loại, \textbf{độ tin cậy của xác suất dự báo} còn được định lượng thông qua chỉ số Brier score và đường cong hiệu chuẩn (\textit{calibration curve}). Nghiên cứu chủ đích \textbf{không} áp dụng các phương pháp hiệu chuẩn hậu nghiệm như \textit{Platt scaling} nhằm phản ánh một cách khách quan năng lực hiệu chuẩn nội tại của từng thuật toán.

\noindent\textbf{Chiến lược phân đoạn và Đánh giá Hold-out.} Quy trình này được thiết lập nhằm triệt tiêu \textbf{thiên kiến lạc quan} (\textit{optimistic bias}) và tăng cường \textbf{tính minh bạch}. Theo đó, mọi quyết định lựa chọn cấu hình chỉ dựa trên thông tin từ \textbf{tập huấn luyện}, đồng thời tuyệt đối \textbf{không} tiếp cận tập kiểm thử \textit{hold-out} trong suốt quá trình tinh chỉnh.\cite{pedregosa2011} Việc này đòi hỏi sự phân biệt rạch ròi giữa hai giai đoạn là \textbf{lựa chọn mô hình} (\textit{model selection}) và \textbf{đánh giá khả năng tổng quát hóa} (\textit{model evaluation}) trên dữ liệu độc lập.

\textbf{Vòng lặp tối ưu nội bộ (\textit{Inner loop}).} Với mỗi thuật toán, các tổ hợp siêu tham số sẽ được đánh giá thông qua \textbf{xác thực chéo $k$-fold phân tầng} ($k=5$) trên 80\% dữ liệu huấn luyện. Cấu hình đạt \textbf{điểm F1 trung bình} cao nhất sẽ được lựa chọn, vốn là một phương pháp giúp giảm thiểu sự phụ thuộc vào các biến động ngẫu nhiên phát sinh trong quá trình phân chia dữ liệu.

\textbf{Tái huấn luyện và Kiểm thử độc lập.} Sau khi xác định được cấu hình tối ưu, mô hình sẽ được \textbf{tái huấn luyện trên toàn bộ} 80\% dữ liệu huấn luyện, sau đó thực hiện \textbf{đánh giá duy nhất một lần} trên 20\% dữ liệu kiểm thử. Tập kiểm thử này vốn được giữ biệt lập hoàn toàn và không tham gia vào bất kỳ bước tiền xử lý hay tinh chỉnh nào, nhờ đó đóng vai trò là cơ sở để ước lượng khách quan cho năng lực dự báo thực tế.

\textbf{Tính khách quan trong đối chiếu.} Mọi mô hình và phương pháp cơ sở (\textit{baseline}) đều được áp dụng đồng nhất một lược đồ chia dữ liệu, đi kèm với luồng xử lý (\textit{pipeline}) tiền xử lý và hệ tiêu chí đánh giá chung. Việc thiết lập quy trình nhất quán này nhằm đảm bảo tính \textbf{tương quan so sánh} cũng như khả năng tái lập hoàn toàn cho nghiên cứu.\cite{pedregosa2011}

\paragraph{Hiệu chuẩn xác suất (\textit{Probability Calibration}).}
Bên cạnh các chỉ số phân hạng như ROC--AUC, \textbf{độ tin cậy của xác suất dự báo} cần được định lượng cụ thể thông qua phân tích hiệu chuẩn. Trong bối cảnh hỗ trợ ra quyết định lâm sàng, các xác suất ước lượng thường được diễn giải trực tiếp nhằm phục vụ việc phân tầng nguy cơ hoặc thiết lập ngưỡng sàng lọc, thay vì chỉ dừng lại ở vai trò xếp hạng thứ tự cho các mẫu dữ liệu.

Nghiên cứu sử dụng hai công cụ bổ trợ sau đây để đánh giá toàn diện khía cạnh này:

\textit{(i) Brier score} là một quy tắc tính điểm phù hợp (\textit{proper scoring rule}) dùng để đo lường sai số bình phương trung bình giữa xác suất dự báo lớp dương tính và nhãn thực tế theo công thức:
\begin{equation}
\mathrm{Brier} = \frac{1}{N}\sum_{i=1}^{N}(\hat{p}_i - y_i)^2,
\end{equation}
trong đó $\hat{p}_i$ đại diện cho xác suất dự báo lớp dương tính và $y_i \in \{0,1\}$ là nhãn quan sát thực tế. Giá trị Brier score \textbf{càng thấp} sẽ thể hiện xác suất dự báo càng tiệm cận với tần suất xuất hiện thực tế của biến cố y sinh.

\textit{(ii) Đường cong hiệu chuẩn} (\textit{Reliability Diagram}): Công cụ này thực hiện phân chia các khoảng xác suất dự báo thành các nhóm (\textit{bins}), từ đó đối chiếu xác suất trung bình trong mỗi nhóm với tỷ lệ dương tính quan sát được tương ứng. Trong đồ thị này, đường chéo $y=x$ sẽ đại diện cho trạng thái hiệu chuẩn hoàn hảo của mô hình.

Nghiên cứu chủ đích \textbf{không} áp dụng các kỹ thuật hiệu chuẩn hậu nghiệm như \textit{Platt scaling} hay \textit{Isotonic regression} trong thực nghiệm chính. Việc này nhằm đánh giá một cách khách quan \textbf{năng lực hiệu chuẩn nội tại} của từng thuật toán dưới cùng một quy trình huấn luyện và tinh chỉnh tham số đã được mô tả ở trên.

\paragraph{Khoảng tin cậy Bootstrap cho ROC--AUC (\textit{Uncertainty Quantification}).}
Để định lượng \textbf{độ bất định} của chỉ số ROC--AUC trên một tập kiểm thử cố định, khoảng tin cậy (CI) 95\% sẽ được ước lượng thông qua phương pháp \textbf{Bootstrap} tái chọn mẫu có hoàn lại. Quy trình thực hiện cụ thể bao gồm việc lấy mẫu ngẫu nhiên $N$ quan sát từ tập kiểm thử ban đầu và tính toán lại ROC--AUC cho mỗi lần lặp, trong đó quy trình này được thực hiện $B=1000$ lần với hạt giống ngẫu nhiên cố định. Sau đó, phân vị thứ 2,5\% và 97,5\% của phân phối Bootstrap sẽ được dùng làm căn cứ để xác định khoảng tin cậy 95\%.

\textbf{Lưu ý}: Khoảng tin cậy này phản ánh sự biến thiên do \textbf{sai số chọn mẫu} bên trong một tập kiểm thử cụ thể, vốn \textbf{không} thay thế cho kiểm định DeLong hay các biến động phát sinh từ quy trình \textbf{phân đoạn dữ liệu} khác nhau. Giá trị trung bình $\pm$ độ lệch chuẩn của ROC--AUC qua 05 lượt chạy với các hạt giống khác nhau (Mục~\ref{subsec:tablesfigures}) sẽ được sử dụng nhằm mô tả độ nhạy của mô hình đối với việc phân chia tập huấn luyện và tập kiểm thử.

\section{Khả năng diễn giải mô hình}
\label{subsec:interpret}

Trong lĩnh vực học máy y sinh, tính \textbf{minh bạch} của mô hình thường được coi trọng tương đương với hiệu năng dự báo, bởi lẽ các kết quả đầu ra cần có khả năng giải thích rõ ràng và đảm bảo tính \textbf{nhất quán} với các kiến thức chuyên môn lâm sàng (Mục~\ref{subsec:features}). Nghiên cứu này tập trung phân tích khả năng diễn giải ở mức \textbf{tổng thể} (\textit{global interpretability}) thông qua sự kết hợp của các phương pháp mang tính \textbf{bổ trợ} lẫn nhau. Các phương pháp này bao gồm việc diễn giải đặc thù theo từng họ mô hình (\textit{model-specific}) kết hợp với \textbf{tầm quan trọng hoán vị} (\textit{Permutation Importance}), vốn là một hướng tiếp cận \textbf{phi mô hình} (\textit{model-agnostic}) giúp đánh giá thống nhất trên mọi thuật toán phân loại. Trong khi đó, việc diễn giải \textbf{cục bộ} (\textit{local interpretability}) cho từng trường hợp bệnh nhân cụ thể, ví dụ như các trị số SHAP cá thể hóa, hiện nằm ngoài phạm vi thực nghiệm chính của đề tài.

\paragraph{Diễn giải đặc thù mô hình (\textit{Model-specific Interpretability}).}

\textbf{Hồi quy Logistic (LR)}: Sau quy trình tiền xử lý và chuẩn hóa, các hệ số hồi quy tương ứng với từng biến số sẽ phản ánh \textbf{chiều hướng} cũng như \textbf{mức độ ảnh hưởng tương đối} lên tỷ số \textit{log-odds} của lớp dương tính. Do các đặc trưng đã được đưa về cùng một thang đo (\textit{Z-score}), độ lớn của các hệ số này có thể được đối chiếu \textbf{trực tiếp} nhằm xác định tầm quan trọng của từng biến số cụ thể.

\textbf{Rừng ngẫu nhiên (RF)}: Mô hình này sử dụng chỉ số \texttt{feature\_importances\_}, vốn dựa trên mức giảm độ không thuần nhất trung bình (\textit{Mean Decrease Impurity}). Chỉ số này cung cấp thứ hạng đóng góp \textbf{tương đối} của các biến, mặc dù cần lưu ý về \textbf{thiên kiến} tiềm ẩn trong trường hợp các đặc trưng có độ tương quan cao hoặc có sự khác biệt về số lượng điểm cắt (\textit{split points}) giữa các biến số với nhau.

\textbf{KNN và SVM}: Do không sở hữu các tham số dạng hệ số trực tiếp như Hồi quy Logistic, phương pháp tầm quan trọng hoán vị sẽ được áp dụng để định lượng mức độ đóng góp của các biến. Việc này giúp đảm bảo tính \textbf{khách quan trong đối chiếu} với các họ mô hình khác trong cùng hệ thống thực nghiệm.

\paragraph{Diễn giải đặc thù mô hình (\textit{Model-specific Interpretability}).}

\textbf{Hồi quy Logistic (LR)}: Sau quy trình tiền xử lý và chuẩn hóa, các hệ số hồi quy tương ứng với từng biến số sẽ phản ánh \textbf{chiều hướng} cũng như \textbf{mức độ ảnh hưởng tương đối} lên tỷ số \textit{log-odds} của lớp dương tính. Do các đặc trưng đã được đưa về cùng một thang đo (\textit{Z-score}), độ lớn của các hệ số này có thể được đối chiếu \textbf{trực tiếp} nhằm xác định tầm quan trọng của từng biến số cụ thể.

\textbf{Rừng ngẫu nhiên (RF)}: Mô hình này sử dụng chỉ số \texttt{feature\_importances\_}, vốn dựa trên mức giảm độ không thuần nhất trung bình (\textit{Mean Decrease Impurity}). Chỉ số này cung cấp thứ hạng đóng góp \textbf{tương đối} của các biến, mặc dù cần lưu ý về \textbf{thiên kiến} tiềm ẩn trong trường hợp các đặc trưng có độ tương quan cao hoặc có sự khác biệt về số lượng điểm cắt (\textit{split points}) giữa các biến số với nhau.

\textbf{KNN và SVM}: Do không sở hữu các tham số dạng hệ số trực tiếp như Hồi quy Logistic, phương pháp tầm quan trọng hoán vị sẽ được áp dụng để định lượng mức độ đóng góp của các biến. Việc này giúp đảm bảo tính \textbf{khách quan trong đối chiếu} với các họ mô hình khác trong cùng hệ thống thực nghiệm.

\paragraph{Đối chiếu các phương pháp đánh giá biến số.}
So với hệ số hồi quy của LR hay tầm quan trọng dựa trên độ không thuần nhất của RF, \textbf{Permutation Importance} thường mang lại kết quả \textbf{ổn định hơn} khi thực hiện so sánh giữa các mô hình có kiến trúc khác biệt. Tuy nhiên, phương pháp này vẫn phụ thuộc vào phân phối dữ liệu cùng quy mô tập kiểm thử. Với kích thước mẫu $n_{\mathrm{test}}$ hạn chế, một số biến số có thể ghi nhận mức đóng góp trung bình xấp xỉ bằng $0$ hoặc mang \textbf{giá trị âm nhỏ} vốn là hệ quả của nhiễu ước lượng. Mặc dù các phương pháp tiên tiến như \textbf{SHAP} (\textit{SHapley Additive exPlanations}) có thể cung cấp góc nhìn toàn diện hơn ở cả cấp độ tổng thể lẫn cục bộ, song kỹ thuật này hiện \textbf{không} nằm trong quy trình thực nghiệm chính của đề tài.\cite{kirbas2025shap} Thay vào đó, SHAP được xem là một hướng mở rộng triển vọng cho các nghiên cứu tiếp theo (Mục~\ref{subsec:disc-future}).

\textbf{Sự tương quan với kiến thức chuyên môn.} Thứ tự ưu tiên của các đặc trưng nhìn chung phù hợp với các kỳ vọng lâm sàng, tiêu biểu là chỉ số Glucose thường đóng vai trò tiên lượng then chốt, theo sau là ảnh hưởng đáng kể của BMI và tác động ổn định của độ tuổi. Kết quả này hoàn toàn nhất quán với các công bố trước đó trên tập dữ liệu Pima cũng như các phân tích diễn giải mô hình gần đây.\cite{chang2022pima,kirbas2025shap} Sự tương đồng này góp phần củng cố tính tin cậy định tính, đồng thời cho thấy mô hình đã khai thác được các tín hiệu y sinh có ý nghĩa thực chất thay vì các nhiễu thống kê thuần túy.

\textbf{Giới hạn và đạo đức trong diễn giải kết quả.} Thứ hạng của các đặc trưng vốn phụ thuộc chặt chẽ vào \textbf{kiến trúc mô hình}, tương quan biến và phân phối mẫu, do đó các kết quả này \textbf{không} đồng nhất với \emph{quan hệ nhân quả} lâm sàng. Các kết quả diễn giải cần được xem là \textbf{bằng chứng hỗ trợ} cho quá trình thảo luận và kiểm chứng, đồng thời tuyệt đối không thay thế cho các quyết định chuyên môn của đội ngũ y bác sĩ.

\section{Khả năng tái lập}
\label{subsec:repro}

Khả năng \textbf{tái lập} vốn là tiêu chí cốt lõi trong học máy thực nghiệm, bởi lẽ đặc tính này cho phép cộng đồng nghiên cứu có thể \textbf{kiểm chứng độc lập} các kết quả và thực hiện đối chiếu công bằng dưới cùng một hệ thống điều kiện kiểm soát. Trong nghiên cứu này, tính tái lập được đảm bảo thông qua ba cấp độ trọng tâm bao gồm \textbf{môi trường phần mềm} (phiên bản thư viện), \textbf{kiểm soát tính ngẫu nhiên} (hạt giống và quy tắc xáo trộn dữ liệu), cùng việc \textbf{đặc tả giao thức thực nghiệm} từ phân đoạn dữ liệu, không gian siêu tham số cho đến tiêu chí đánh giá và phương pháp tổng hợp kết quả.

Toàn bộ quy trình được triển khai dựa trên thư viện \textbf{scikit-learn} với cấu trúc \texttt{Pipeline} nhất quán cho cả giai đoạn tiền xử lý và huấn luyện, qua đó giúp triệt tiêu rủi ro sai lệch quy trình giữa các mô hình khác nhau.\cite{pedregosa2011} Các thông tin về phiên bản Python, NumPy và scikit-learn đều được ghi nhận chi tiết tại Mục~\ref{subsec:protocol} cũng như trong tệp \texttt{manifest.json} thuộc thư mục kết quả. Việc lưu trữ dữ liệu này nhằm hạn chế tối đa các sai số hệ thống có thể phát sinh từ sự khác biệt giữa các phiên bản thư viện trong quá trình thực thi.

\textbf{Kiểm soát tính ngẫu nhiên}: Các vòng lặp thực nghiệm bên trong sử dụng \texttt{StratifiedKFold(..., shuffle=True, random\_state=42)} và \texttt{GridSearchCV(..., random\_state=42)}. Đồng thời, các mô hình có thành phần ngẫu nhiên như Rừng ngẫu nhiên đều được thiết lập cố định giá trị \texttt{random\_state=42}. Quy trình \textbf{phân đoạn tập huấn luyện và tập kiểm thử} được lặp lại với năm giá trị hạt giống khác nhau bao gồm \texttt{random\_state} $\in\{0,1,2,3,4\}$ nhằm đánh giá chính xác độ ổn định của mô hình. Cần lưu ý rằng thuật toán \textbf{KNN} không sở hữu tham số ngẫu nhiên nội tại, do đó tính tái lập của thuật toán này phụ thuộc hoàn toàn vào cấu trúc dữ liệu cũng như các tham số khoảng cách đã được cố định trước đó (Phụ lục~\ref{app:seeds}).

\textbf{Minh bạch thực nghiệm}: Các yếu tố như tỷ lệ dữ liệu \textit{hold-out}, phương pháp phân tầng (\textit{stratify}), không gian tìm kiếm siêu tham số (Phụ lục~\ref{app:hyper}), tiêu chí tối ưu \texttt{scoring='f1'} (Mục~\ref{subsec:metrics}), cùng quy cách báo cáo kết quả theo định dạng trung bình $\pm$ độ lệch chuẩn vốn là \textbf{những thành phần then chốt} giúp báo cáo này có thể được kiểm chứng một cách đầy đủ.

\textbf{Mã nguồn và Tài nguyên}: Tập tin mã nguồn \texttt{resources/diabetes.ipynb} đóng vai trò là thành phần thực thi chính, trong khi thư mục \texttt{resources/outputs/} chịu trách nhiệm lưu trữ các bảng kết quả đồng bộ với nội dung trình bày tại Chương~\ref{sec:results}. Mã nguồn còn đi kèm các tệp quản lý môi trường bao gồm \texttt{pyproject.toml} và \texttt{uv.lock} cùng Phụ lục~\ref{app:code} nhằm hỗ trợ tối đa cho việc tái lập và kiểm chứng độc lập.

\paragraph{Danh mục đặc tả báo cáo (\textit{Reporting Checklist}).}
\begin{itemize}
  \item \textbf{Hạt giống và phân đoạn dữ liệu}: Nghiên cứu sử dụng phương pháp \textit{hold-out} phân tầng (\textit{stratified}) đi kèm với việc cố định hạt giống (\textit{seed}) cho các vòng lặp tối ưu nội bộ. Đồng thời, quy trình chia tập dữ liệu (\texttt{train\_test\_split}) được thực hiện lặp lại với năm hạt giống khác nhau nhằm đánh giá chính xác độ ổn định của các ước lượng.
  \item \textbf{Tính lặp lại thực nghiệm}: Quy trình xác định rõ số lượt thực thi bao gồm 05 lượt phân đoạn dữ liệu độc lập, kết hợp với phương pháp tổng hợp kết quả dựa trên giá trị trung bình và độ lệch chuẩn tương ứng.
  \item \textbf{Khoảng tin cậy (CI)}: Chỉ số F1-score được báo cáo kèm theo khoảng tin cậy 95\% dựa trên phân phối $t$-Student qua các lượt chạy thực nghiệm. Ngoài ra, các phương pháp xác định CI khác như Bootstrap cho AUC cũng được đặc tả rõ ràng về trình tự và quy trình thực hiện.
  \item \textbf{Môi trường thực thi}: Nhằm đảm bảo tính minh bạch, nghiên cứu ghi nhận đầy đủ phiên bản của các công cụ bao gồm Python, scikit-learn và NumPy, đồng thời cung cấp thông số hệ điều hành cũng như các phương thức quản lý môi trường như \texttt{uv} hoặc môi trường ảo.
\end{itemize}

Trong trường hợp chênh lệch hiệu năng giữa các mô hình \textbf{không đáng kể}, nghiên cứu thực hiện bổ sung các kiểm định cặp (\textit{paired tests}) cho những đối chiếu trọng tâm. Cụ thể, kiểm định \textbf{Wilcoxon signed-rank} được áp dụng trên điểm F1 theo từng lượt chia dữ liệu, ví dụ như khi so sánh giữa mô hình RF và KNN. Các phương pháp bổ trợ khác bao gồm kiểm định \textbf{McNemar} dựa trên ma trận nhầm lẫn hoặc \textbf{paired $t$-test} trong điều kiện thỏa mãn các giả định về phân phối chuẩn.

\paragraph{Độ lớn hiệu ứng (\textit{Effect Size -- Cliff's Delta}).}
Bên cạnh giá trị $p$-value, việc đánh giá \textbf{độ lớn hiệu ứng} giúp mô tả \emph{mức độ chênh lệch thực tế} giữa các mô hình, qua đó tránh việc phụ thuộc cực đoan vào ý nghĩa thống kê thuần túy. Nghiên cứu sử dụng chỉ số \textbf{Cliff's delta} ($\delta$), vốn là một phép đo phi tham số phù hợp để so sánh hai nhóm điểm số F1 với mỗi điểm tương ứng cho một lần chia tập độc lập. Với hai vectơ điểm số có độ dài $n$ (ở đây $n=5$ lượt chia), Cliff's delta được xác định theo công thức:
\[
\hat{\delta} = \frac{\#\{(i,j): x_i > y_j\} - \#\{(i,j): x_i < y_j\}}{n_x n_y},
\]
trong đó $x$ và $y$ lần lượt là điểm F1 của hai mô hình được so sánh theo thứ tự hạt giống. Giá trị $\delta$ nằm trong khoảng $[-1, 1]$, đồng thời dựa trên các ngưỡng diễn giải chuẩn, mức độ ảnh hưởng sẽ được phân loại thành các mức bao gồm: \textbf{không đáng kể} ($|\hat{\delta}| < 0,147$), \textbf{nhỏ} ($< 0,33$), \textbf{vừa} ($< 0,474$) và \textbf{lớn} ($\geq 0,474$). Kết quả chi tiết của phân tích này được trình bày tại Mục~\ref{subsec:uncertainty-effect}.

\noindent\textbf{Tổng kết thiết kế nghiên cứu.} Quy trình thực nghiệm được thiết lập nhằm đảm bảo \emph{tính khách quan} khi đối chiếu các thuật toán, đồng thời giúp \emph{triệt tiêu rò rỉ dữ liệu} (\textit{data leakage}) trong giai đoạn tiền xử lý và tối ưu hóa siêu tham số (Mục~\ref{subsec:models}). Bên cạnh đó, quy trình còn cung cấp các \emph{ước lượng độc lập} thông qua giao thức \textit{hold-out} (Mục~\ref{subsec:metrics}) với khả năng \emph{tái lập} cao (Mục~\ref{subsec:repro}). Những hạn chế về quy mô dữ liệu cũng như phạm vi mô hình sẽ được thảo luận chi tiết tại Chương~\ref{sec:discussion}.

\chapter{Kết quả thực nghiệm}
\label{sec:results}

\section{Giao thức thực nghiệm}
\label{subsec:protocol}

Nội dung này trình bày chi tiết về \textbf{thiết lập thực nghiệm} được áp dụng để thu thập các kết quả định lượng trong hệ thống bảng biểu và hình vẽ của chương. Quy trình thực hiện được xây dựng nhằm bảo đảm tính thống nhất với phương pháp nghiên cứu đã đề xuất tại Chương~\ref{sec:methods}, đồng thời cho phép kiểm chứng độc lập thông qua tệp mã nguồn \texttt{resources/diabetes.ipynb} cùng các tài nguyên (\textit{artifacts}) đi kèm đã công bố.\cite{pedregosa2011}

\textbf{Phân đoạn dữ liệu.} Tập dữ liệu Pima được phân tách theo phương pháp \textbf{hold-out với tỷ lệ 80/20}. Trong đó, hàm xử lý được cấu hình với tham số \texttt{stratify=y} nhằm duy trì nguyên vẹn tỷ lệ phân bố giữa các lớp trên cả tập huấn luyện (\textit{train}) và tập kiểm thử (\textit{test}). Với mục đích ước lượng \textbf{độ biến thiên} của mô hình trước các phương thức phân chia dữ liệu ngẫu nhiên, toàn bộ luồng xử lý (\textit{pipeline}) được \textbf{lặp lại 05 lần}. Tại mỗi lượt thực thi, giá trị hạt giống (\texttt{random\_state}) sẽ được thay đổi tuần tự trong tập $\{0,1,2,3,4\}$ để tạo ra năm bộ dữ liệu huấn luyện và kiểm thử độc lập. Lưu ý rằng các bảng kết quả chi tiết trong chương này minh họa cho một lượt chạy đơn lẻ với hạt giống số 0, còn các số liệu thống kê tổng hợp sẽ được tính toán dựa trên kết quả trung bình của cả năm lượt thực nghiệm nói trên.

\textbf{Tối ưu hóa và tái huấn luyện.} Đối với mỗi tập huấn luyện cụ thể, quy trình lựa chọn siêu tham số tối ưu được thực hiện thông qua công cụ \texttt{GridSearchCV} kết hợp cùng phương pháp \textbf{xác thực chéo 5-fold phân tầng} (\textit{StratifiedKFold}). Nghiên cứu sử dụng tiêu chí \textbf{F1-score} (\texttt{scoring='f1'}) làm hàm mục tiêu tối ưu hóa nhằm đảm bảo sự cân bằng giữa hai chỉ số Precision và Recall, đặc biệt là trong bối cảnh dữ liệu có hiện tượng lệch lớp nhẹ. Để đảm bảo tính khách quan, nghiên cứu cam kết \textbf{triệt tiêu hoàn toàn hiện tượng rò rỉ dữ liệu} bằng cách không sử dụng bất kỳ thông tin nào từ tập kiểm thử cho các bước tiền xử lý hay tinh chỉnh tham số. Sau khi xác định được cấu hình tối ưu, mô hình sẽ được \textbf{tái huấn luyện} trên toàn bộ dữ liệu của tập huấn luyện đó trước khi tiến hành đánh giá \textbf{duy nhất một lần} trên tập kiểm thử tương ứng. Chi tiết về không gian siêu tham số được trình bày cụ thể tại Phụ lục~\ref{app:hyper}.

\textbf{Hệ chỉ số và phương pháp tổng hợp.} Hiệu năng thực nghiệm trên mỗi tập kiểm thử được định lượng thông qua các chỉ số bao gồm: \textbf{Accuracy}, \textbf{Precision}, \textbf{Recall}, \textbf{F1-score}, \textbf{ROC--AUC}, \textbf{PR--AUC} và \textbf{Brier score} (chi tiết tại Mục~\ref{subsec:metrics}). Các kết quả cuối cùng được trình bày dưới dạng \textbf{trung bình $\pm$ độ lệch chuẩn} tính toán qua năm lượt phân đoạn dữ liệu. Riêng đối với chỉ số F1-score, nghiên cứu còn cung cấp thêm \textbf{khoảng tin cậy 95\%} dựa trên phân phối $t$-Student với 04 bậc tự do. Song song đó, các quy trình tính toán đường hiệu chuẩn, PR--AUC, Bootstrap ROC--AUC và độ lớn hiệu ứng Cliff's delta đều được thực thi thông qua hệ thống tập lệnh chuyên biệt trong thư mục \texttt{resources/}.

\textbf{Mô hình cơ sở (\textit{Baseline}).} Nghiên cứu thiết lập \texttt{DummyClassifier} với chiến lược \textbf{dự báo lớp đa số} để làm đường cơ sở tham chiếu. Nhằm bảo đảm tính khách quan và công bằng trong quá trình đối chiếu hiệu năng, mô hình này cũng được đặt trong \textbf{cùng} một luồng xử lý tiền xử lý tương tự như các thuật toán chính được đề xuất.

\textbf{Môi trường thực thi.} Toàn bộ quá trình thực nghiệm được triển khai trên nền tảng Python 3.13.12 với các thư viện cốt lõi bao gồm scikit-learn 1.8.0 và NumPy 2.4.3, trong đó các gói phụ thuộc được quản lý chặt chẽ thông qua công cụ \texttt{uv}. Các thông tin chi tiết về phiên bản phần mềm, cấu hình hệ thống cũng như danh sách các hạt giống ngẫu nhiên đã sử dụng hiện được lưu trữ tại tệp \texttt{manifest.json} và trình bày thêm tại Phụ lục~\ref{app:seeds}.

\section{Phân tích kết quả định lượng}
\label{subsec:tablesfigures}

Mục này trình bày hệ thống các kết quả định lượng thu được từ thiết lập thực nghiệm đã được mô tả chi tiết tại Mục~\ref{subsec:protocol}. Phương pháp trình bày kết quả được xây dựng nhằm phản ánh đồng thời hai khía cạnh quan trọng là \textbf{ước lượng điểm} (\textit{point estimate}) dựa trên một lượt phân đoạn cụ thể và \textbf{độ biến thiên} (\textit{variability}) thông qua nhiều lần phân chia dữ liệu độc lập. Trong bối cảnh cỡ mẫu nghiên cứu còn hạn chế, việc kết hợp hai góc nhìn này là thực sự thiết yếu nhằm mang lại cái nhìn toàn diện và khách quan nhất.

\textbf{Hệ thống nhận định.} Sự khác biệt về hiệu năng giữa các mô hình không chỉ được thể hiện thông qua trị số trung bình mà còn nằm ở độ ổn định trước những thay đổi của dữ liệu huấn luyện. Trên thực tế, một mô hình có hiệu năng kỳ vọng cao nhưng đi kèm độ dao động lớn thường không phải là lựa chọn tối ưu khi triển khai vào môi trường thực tiễn. Chính vì vậy, các kết quả thu được sẽ được diễn giải dựa trên ba tầng thông tin bổ trợ lẫn nhau, bao gồm: \textbf{hiệu năng kỳ vọng}, \textbf{độ bất định} và cuối cùng là \textbf{ngữ cảnh ứng dụng} cụ thể.

\textbf{Tính nhạy cảm của giao thức thực nghiệm.} Kết quả cho thấy thứ hạng giữa các thuật toán có sự thay đổi nhất định tùy thuộc vào tiêu chí đánh giá được lựa chọn (chẳng hạn như F1-score so với ROC--AUC) cũng như phương thức phân đoạn dữ liệu. Điều này minh chứng rằng việc xếp hạng mô hình không phải là một đặc tính bất biến của bản thân thuật toán, mà thực chất là kết quả của sự tương tác phức hợp giữa kiến trúc mô hình, chỉ số đo lường và giao thức thực nghiệm được áp dụng.

\textbf{Hiệu năng trên một lượt phân đoạn đại diện.} Bảng~\ref{tab:results-point} ghi nhận kết quả từ một lần thực thi cụ thể với giá trị hạt giống được cố định tại số 0 (\texttt{random\_state=0}). Dữ liệu này cung cấp cái nhìn trực quan về thứ hạng tương đối giữa các phương pháp trong một kịch bản cố định với đầy đủ hệ thống chỉ số đo lường. Trong phân đoạn này, mô hình SVM (nhân RBF) cùng với Rừng ngẫu nhiên ghi nhận mức F1-score tương đương và cho thấy sự vượt trội so với các mô hình còn lại, trong khi đó Hồi quy Logistic lại đạt được mức ROC--AUC tối ưu nhất.

\textbf{Hiệu năng tổng hợp đa tầng dữ liệu.} Bảng~\ref{tab:results-meanstd} thực hiện tổng hợp hiệu năng dựa trên \textbf{05} lượt phân đoạn độc lập và được trình bày dưới dạng giá trị trung bình kèm theo độ lệch chuẩn. Quy cách báo cáo này không chỉ giúp nhận diện mức độ \textbf{dao động} của mô hình mà còn hỗ trợ việc đưa ra kết luận khách quan hơn, tránh phụ thuộc vào các kết quả ngẫu nhiên đơn lẻ. Bên cạnh đó, cột F1-score được bổ sung thêm \textbf{khoảng tin cậy 95\%} dựa trên ước lượng từ phân phối $t$-Student nhằm định lượng độ bất định thống kê một cách tổng quát và khoa học.

\textbf{Diễn giải kết quả tổng hợp.} Xét trên giá trị F1 trung bình, mô hình \textbf{Rừng ngẫu nhiên} thể hiện ưu thế nhẹ, trong khi đó \textbf{Hồi quy Logistic} lại dẫn đầu về các chỉ số Accuracy và ROC--AUC trung bình. Sự khác biệt này phản ánh rõ rệt tính đánh đổi tùy thuộc vào mục tiêu tối ưu hóa được ưu tiên. Đáng chú ý là chênh lệch về chỉ số F1 giữa các mô hình thuộc nhóm dẫn đầu là \textbf{tương đối nhỏ}. Kết quả kiểm định \textbf{Wilcoxon signed-rank} trên điểm F1 giữa hai mô hình tốt nhất là Rừng ngẫu nhiên và KNN cho giá trị $W=4,00$ với $p=0,438$ (dựa trên $n=5$ cặp quan sát), điều này cho thấy chưa đủ cơ sở để bác bỏ giả thuyết không có sự khác biệt ở mức ý nghĩa thống kê thông thường. Tuy nhiên, cần lưu ý rằng với cỡ mẫu $n=5$, \textbf{lực thống kê} (\textit{statistical power}) của phép kiểm định còn ở mức thấp. Do đó, kết quả này không khẳng định sự tương đương tuyệt đối giữa các mô hình mà chỉ cho thấy hiện tại \textbf{chưa đủ bằng chứng} để xác nhận sự khác biệt có ý nghĩa trong thiết lập thực nghiệm này. Quan sát này gợi ý rằng những biến thiên ghi nhận được có thể xuất phát từ phương thức phân tầng mẫu thay vì ưu thế nội tại từ kiến trúc của thuật toán.

\textbf{Ý nghĩa thống kê và thực tiễn.} Với quy mô mẫu hiện tại, các kết quả kiểm định cho thấy sự cạnh tranh rất sít sao giữa các mô hình dẫn đầu. Chính vì vậy, việc lựa chọn mô hình triển khai trong thực hành lâm sàng nên cân nhắc thêm các yếu tố về \textbf{độ lớn hiệu ứng} và \textbf{tính ổn định}, thay vì chỉ dựa vào bảng xếp hạng các trị số trung bình đơn thuần.

\textbf{Phân tích sai số và khả năng phân tách.} Hình~\ref{fig:roc-confusion} minh họa đường cong ROC cùng ma trận nhầm lẫn (\textit{confusion matrix}) của mô hình Rừng ngẫu nhiên. Trong khi đường ROC phản ánh năng lực phân tách của mô hình trên toàn bộ phổ ngưỡng quyết định, ma trận nhầm lẫn lại giúp đặc tả chi tiết các sai lầm loại I và loại II. Đây được coi là những tham số then chốt trong bài toán \emph{tầm soát y tế}, nơi đòi hỏi sự cân bằng tối ưu giữa rủi ro bỏ sót ca bệnh và áp lực vận hành phát sinh do các trường hợp báo động giả.

\textbf{Các phân tích bổ trợ.} Bên cạnh các kết quả chính, hiệu năng ước lượng xác suất bao gồm chỉ số Brier score và đường hiệu chuẩn được trình bày chi tiết tại Mục~\ref{subsec:calibration}. Song song với đó, các phân tích chuyên sâu về đường cong Precision--Recall được thực hiện tại Mục~\ref{subsec:prauc}. Các phép đo về độ bất định ROC--AUC cùng với độ lớn hiệu ứng Cliff's delta cũng được mô tả cụ thể tại Mục~\ref{subsec:uncertainty-effect}. Cuối cùng, kết quả định lượng về tầm quan trọng của các đặc trưng thông qua kỹ thuật hoán vị được tổng hợp đầy đủ tại Mục~\ref{subsec:permimp}.

\noindent\textit{Tóm lại,} hệ thống bảng biểu và hình ảnh nêu trên đã cung cấp những bằng chứng thực nghiệm khách quan, đảm bảo tính nhất quán với giao thức đánh giá đã thiết lập từ trước. Những thảo luận chuyên sâu nhằm giải quyết các câu hỏi nghiên cứu (RQ) cũng như các liên hệ trong thực tế lâm sàng sẽ được trình bày chi tiết tại Chương~\ref{sec:discussion}.

\begin{table}[htbp]
  \centering
  \caption{Hiệu năng thực nghiệm trên một lượt phân đoạn đại diện (Hold-out 20\%, \texttt{stratify}, \texttt{random\_state=0}).}
  \label{tab:results-point}
  \begin{tabular}{lccccc}
    \toprule
    Mô hình & Accuracy & Precision & Recall & F1-score & ROC--AUC \\
    \midrule
    Logistic Regression & 0,766 & 0,737 & 0,519 & 0,609 & 0,870 \\
    SVM (RBF)             & 0,773 & 0,711 & 0,593 & 0,646 & 0,855 \\
    Random Forest         & 0,779 & 0,738 & 0,574 & 0,646 & 0,858 \\
    KNN                   & 0,740 & 0,659 & 0,537 & 0,592 & 0,794 \\
    Majority (Baseline)   & 0,649 & 0,000 & 0,000 & 0,000 & 0,500 \\
    \bottomrule
  \end{tabular}
\end{table}

\noindent\textbf{Đánh giá bối cảnh áp dụng.} Các kết quả ghi nhận tại Bảng~\ref{tab:results-point} gợi ý rằng những mô hình phi tuyến như Rừng ngẫu nhiên (RF) và SVM-RBF có khả năng khai thác các tương tác đặc trưng phức tạp hiệu quả hơn trên lượt phân đoạn đại diện, từ đó dẫn đến chỉ số F1-score vượt trội hơn so với mô hình Hồi quy Logistic (LR). Tuy nhiên, mô hình LR vẫn duy trì được mức ROC--AUC cao nhất trong cùng lượt thực thi này. Quan sát này cho thấy trong trường hợp mục tiêu ưu tiên là chất lượng phân hạng xác suất thay vì tối ưu hóa chỉ số F1 tại một ngưỡng cố định, các mô hình tuyến tính truyền thống vẫn là một phương án có tính cạnh tranh rất cao.

\begin{table}[htbp]
  \centering
  \small
  \setlength{\tabcolsep}{4pt}
  \caption{Kết quả tổng hợp (Trung bình $\pm$ Độ lệch chuẩn) qua 05 lượt phân đoạn; Khoảng tin cậy (CI) 95\% cho F1-score.}
  \label{tab:results-meanstd}
  \begin{tabular}{l p{2.8cm} p{3.6cm} p{2.8cm}}
    \toprule
    Mô hình & Acc. (mean $\pm$ std) & F1 (mean $\pm$ std) [95\% CI] & AUC (mean $\pm$ std) \\
    \midrule
    Random Forest         & $0,743 \pm 0,024$ & $0,614 \pm 0,025$ \quad [$0,583$; $0,645$] & $0,823 \pm 0,025$ \\
    KNN                   & $0,729 \pm 0,033$ & $0,595 \pm 0,058$ \quad [$0,523$; $0,668$] & $0,778 \pm 0,023$ \\
    Logistic Regression   & $0,748 \pm 0,018$ & $0,590 \pm 0,017$ \quad [$0,570$; $0,611$] & $0,824 \pm 0,031$ \\
    SVM (RBF)             & $0,736 \pm 0,035$ & $0,572 \pm 0,046$ \quad [$0,515$; $0,629$] & $0,805 \pm 0,059$ \\
    Majority (Baseline)   & $0,649 \pm 0,000$ & $0,000 \pm 0,000$ \quad [$0,000$; $0,000$] & $0,500 \pm 0,000$ \\
    \bottomrule
  \end{tabular}
\end{table}

\noindent\textbf{Phân tích biến thiên và độ ổn định.} Kết quả từ Bảng~\ref{tab:results-meanstd} cho thấy mặc dù mô hình Rừng ngẫu nhiên (RF) dẫn đầu về chỉ số F1-score trung bình, song Hồi quy Logistic (LR) lại ghi nhận độ lệch chuẩn thấp hơn đáng kể. Hiện tượng này phản ánh rõ nét sự đánh đổi giữa \textbf{năng lực tối ưu hiệu năng} và \textbf{tính ổn định} trước những thay đổi của dữ liệu huấn luyện, cụ thể là các mô hình có độ linh hoạt cao thường đạt được kết quả đỉnh tốt hơn nhưng lại nhạy cảm hơn trước các biến thiên của mẫu dữ liệu. Một quan sát quan trọng khác nằm ở sự khác biệt về độ ổn định giữa các thuật toán, trong đó RF dao động ở mức vừa phải, còn KNN và SVM lại cho thấy sự biến thiên mạnh mẽ tùy thuộc vào thành phần dữ liệu đầu vào. Điều này khẳng định rằng \textbf{độ nhạy với phương thức phân đoạn dữ liệu} chính là một thuộc tính thực nghiệm quan trọng cần được cân nhắc song hành cùng các giá trị trung bình kỳ vọng khi đánh giá mô hình.

\begin{table}[htbp]
  \centering
  \caption{Tổng hợp sự đánh đổi trong đặc tính vận hành của các mô hình dưới cùng một giao thức thực nghiệm.}
  \label{tab:model-tradeoff-summary}
  \begin{tabular}{lp{5cm}p{5cm}}
    \toprule
    Mô hình & Ưu điểm chủ đạo & Hạn chế và rủi ro tiềm ẩn \\
    \midrule
    Rừng ngẫu nhiên & Khai thác tốt các tương tác phi tuyến và đạt chỉ số F1-score tối ưu. & Độ biến thiên theo lượt phân đoạn lớn đồng thời chi phí diễn giải kỹ thuật cao. \\
    Hồi quy Logistic & Tính ổn định cao với hệ số mô hình tường minh, tạo điều kiện thuận lợi cho diễn giải lâm sàng. & Rủi ro thiếu khớp (\textit{underfitting}) trong trường hợp dữ liệu tồn tại các quan hệ phi tuyến phức tạp. \\
    \bottomrule
  \end{tabular}
\end{table}

\begin{figure}[htbp]
  \centering
  \includegraphics[width=\textwidth]{image/roc_confusion_test_seed0.png}
  \caption{Đường cong ROC (AUC = 0,858) và ma trận nhầm lẫn của mô hình Rừng ngẫu nhiên trên tập kiểm thử (\texttt{split\_seed=0}, $n_{\mathrm{test}}=154$). Các thông số chi tiết bao gồm: TN=89, FP=11, FN=23, TP=31.}
  \label{fig:roc-confusion}
\end{figure}

\section{Phân tích hiệu chuẩn xác suất}
\label{subsec:calibration}

Bên cạnh các chỉ số phân loại truyền thống, nghiên cứu tiến hành định lượng \textbf{độ tin cậy của xác suất dự báo} lớp dương tính thông qua chỉ số Brier score và đường hiệu chuẩn (chi tiết tại Mục~\ref{subsec:metrics}). Bảng~\ref{tab:results-brier} thực hiện tổng hợp giá trị trung bình kèm theo độ lệch chuẩn của Brier score dựa trên \textbf{05} lượt phân đoạn dữ liệu độc lập, vốn được trích xuất từ kịch bản xử lý \texttt{resources/compute\_calibration.py} và tệp dữ liệu \texttt{brier\_mean\_std.csv}.

Về mặt lý thuyết, chỉ số Brier score càng thấp sẽ chứng tỏ các xác suất dự báo càng tiệm cận với nhãn thực tế theo tiêu chí sai số bình phương trung bình. Trong thiết lập thực nghiệm này, mô hình \textbf{Rừng ngẫu nhiên} đạt được giá trị Brier trung bình thấp nhất và theo sát ngay sau đó là \textbf{Hồi quy Logistic}. Các mô hình bao gồm \textbf{SVM (RBF)} và \textbf{KNN} lần lượt xếp ở các vị trí tiếp theo, trong đó SVM bộc lộ biên độ biến thiên lớn hơn đáng kể giữa các lượt thực thi, điều này được thể hiện rõ qua giá trị độ lệch chuẩn ở mức cao. Tại lượt phân đoạn đại diện ứng với hạt giống số 0 (\texttt{split\_seed=0}), Hồi quy Logistic cùng với Rừng ngẫu nhiên cho thấy kết quả Brier tương đồng, hoàn toàn nhất quán với xu hướng hội tụ của đường hiệu chuẩn về phía đường chéo lý tưởng được minh họa trong Hình~\ref{fig:calibration}.

Kết quả đối chiếu giữa các chỉ số ROC--AUC, PR--AUC và Brier score cho thấy thứ hạng của các mô hình \textbf{không hoàn toàn có sự nhất quán giữa các hệ chỉ số khác nhau}. Thực tế này minh chứng rằng một mô hình sở hữu năng lực phân hạng (\textit{ranking}) tốt chưa chắc đã đạt được mức hiệu chuẩn xác suất tối ưu. Quan sát trên một lần nữa tái khẳng định rằng quy trình lựa chọn chỉ số đánh giá cần phải bám sát mục tiêu đặc thù của từng tác vụ lâm sàng cụ thể thay vì chỉ dựa vào một thước đo duy nhất.

Hình~\ref{fig:calibration} minh họa đường hiệu chuẩn (\textit{reliability diagram}) được xây dựng theo chiến lược phân nhóm đều (\texttt{uniform}, 10 bins) cho bốn mô hình trên tập kiểm thử ứng với hạt giống số 0 (\texttt{split\_seed=0}). Trong đó, đường đứt đoạn đóng vai trò đại diện cho trạng thái hiệu chuẩn hoàn hảo ($y=x$). Quan sát biểu đồ cho thấy Hồi quy Logistic và Rừng ngẫu nhiên có xu hướng bám sát đường lý tưởng tại hầu hết các phân đoạn xác suất. Ngược lại, mô hình KNN lại bộc lộ sự dao động mạnh hơn, điều này xuất phát từ bản chất dự báo dựa trên lân cận cục bộ cũng như độ nhạy cao đối với nhiễu dữ liệu. Kết quả phân tích nhấn mạnh một nhận định quan trọng rằng mặc dù \textbf{khả năng phân loại mạnh} (thể hiện qua F1-score và ROC--AUC) cùng \textbf{độ hiệu chuẩn xác suất tốt} có thể xuất hiện đồng thời, song việc báo cáo tách biệt hai khía cạnh này vẫn là yêu cầu bắt buộc trong các ứng dụng hỗ trợ quyết định y khoa.

\begin{table}[htbp]
  \centering
  \caption{Kết quả Brier score (Trung bình $\pm$ Độ lệch chuẩn) qua 05 lượt phân đoạn dữ liệu huấn luyện/kiểm thử.}
  \label{tab:results-brier}
  \begin{tabular}{lc}
    \toprule
    Mô hình & Brier score (mean $\pm$ std) \\
    \midrule
    Random Forest         & $0,163 \pm 0,012$ \\
    Logistic Regression   & $0,165 \pm 0,013$ \\
    SVM (RBF)             & $0,171 \pm 0,023$ \\
    KNN                   & $0,183 \pm 0,014$ \\
    \bottomrule
  \end{tabular}
\end{table}

\begin{figure}[htbp]
  \centering
  \includegraphics[width=0.72\textwidth]{image/calibration_curves_test_seed0.png}
  \caption{Đường hiệu chuẩn (\textit{Reliability Diagram}) của bốn mô hình trên tập kiểm thử (\texttt{split\_seed=0}); đường đứt đoạn biểu thị trạng thái hiệu chuẩn hoàn hảo.}
  \label{fig:calibration}
\end{figure}

\section{Đường cong Precision--Recall và chỉ số PR--AUC (\textit{PR Analysis})}
\label{subsec:prauc}

Trong điều kiện lớp dương tính chiếm tỷ lệ thấp như đã trình bày tại Mục~\ref{subsec:dataset}, việc sử dụng đường cong \textbf{Precision--Recall} cùng chỉ số \textbf{PR--AUC} mang lại những thông tin bổ trợ then chốt bên cạnh ROC--AUC. Cụ thể là phân tích PR sẽ tập trung sâu vào sự đánh đổi giữa độ chính xác dương tính (Precision) và độ nhạy (Recall) trên nhóm đối tượng mục tiêu (chi tiết tại Mục~\ref{subsec:metrics}). Bảng~\ref{tab:results-prauc} thực hiện tổng hợp giá trị PR--AUC trung bình kèm theo độ lệch chuẩn thông qua 05 lượt phân đoạn dữ liệu độc lập, dựa trên dữ liệu được trích xuất từ tệp \texttt{pr\_auc\_mean\_std.csv}.

Về kết quả thực nghiệm, \textbf{Rừng ngẫu nhiên} đạt được chỉ số PR--AUC trung bình cao nhất và theo sát ngay sau đó là mô hình \textbf{Hồi quy Logistic}. Các mô hình \textbf{SVM (RBF)} cùng \textbf{KNN} xếp ở những vị trí tiếp theo, trong đó SVM tiếp tục ghi nhận mức độ lệch chuẩn lớn, hoàn toàn tương ứng với sự biến động đã được quan sát trước đó ở chỉ số ROC--AUC. Hình~\ref{fig:pr-curves} minh họa trực quan các đường cong Precision--Recall trên tập kiểm thử ứng với lượt chia hạt giống số 0 (\texttt{split\_seed=0}). Tại đây, đường ngang nét đứt biểu thị tỷ lệ mẫu dương tính thực tế trong tập kiểm thử, đóng vai trò là đường cơ sở (\textit{baseline}) cho một bộ phân loại ngẫu nhiên. Trong kịch bản thực tế này, các đường cong của RF và LR duy trì được vị thế ưu thế rõ rệt so với KNN, điều này hoàn toàn nhất quán với những kết luận rút ra từ giá trị PR--AUC trung bình.

Dưới góc độ diễn giải, điểm mấu chốt không chỉ nằm ở việc mô hình Rừng ngẫu nhiên (RF) chiếm ưu thế về chỉ số PR--AUC, mà còn thể hiện ở sự hoán đổi thứ hạng giữa các mô hình khi được đánh giá qua các hệ chỉ số khác nhau như ROC--AUC, PR--AUC hay Brier score. Thực tế này càng củng cố thêm luận điểm rằng giao thức thực nghiệm cùng việc lựa chọn chỉ số đo lường có tầm ảnh hưởng trực tiếp đến quy trình xác định đâu là "mô hình tối ưu" trong một bài toán cụ thể.

\begin{table}[htbp]
  \centering
  \caption{Kết quả PR--AUC (Trung bình $\pm$ Độ lệch chuẩn) qua 05 lượt phân đoạn dữ liệu huấn luyện/kiểm thử.}
  \label{tab:results-prauc}
  \begin{tabular}{lc}
    \toprule
    Mô hình & PR--AUC (mean $\pm$ std) \\
    \midrule
    Random Forest         & $0,701 \pm 0,046$ \\
    Logistic Regression   & $0,692 \pm 0,050$ \\
    SVM (RBF)             & $0,672 \pm 0,097$ \\
    KNN                   & $0,615 \pm 0,057$ \\
    \bottomrule
  \end{tabular}
\end{table}

\begin{figure}[htbp]
  \centering
  \includegraphics[width=0.72\textwidth]{image/pr_curves_test_seed0.png}
  \caption{Đường cong Precision--Recall của bốn mô hình trên tập kiểm thử (\texttt{split\_seed=0}). Lưu ý: AP tương ứng với giá trị PR--AUC thực nghiệm; đường nét đứt biểu thị tỷ lệ lớp dương tính thực tế.}
  \label{fig:pr-curves}
\end{figure}

\section{Độ bất định của ROC--AUC và độ lớn hiệu ứng}
\label{subsec:uncertainty-effect}

\textbf{Phân tích độ bất định của ROC--AUC.} Bảng~\ref{tab:results-auc-bootstrap} trình bày chỉ số ROC--AUC trên tập kiểm thử ứng với lượt phân đoạn đại diện (\texttt{split\_seed=0}), đi kèm với khoảng tin cậy (CI) 95\% được ước lượng thông qua phương pháp Bootstrap với $B=1000$ lần lặp (chi tiết tại Mục~\ref{subsec:metrics}). Kết quả cho thấy khoảng tin cậy của các mô hình thuộc nhóm dẫn đầu bao gồm Hồi quy Logistic (LR), Rừng ngẫu nhiên (RF) và SVM ghi nhận sự \textbf{chồng lấn (\textit{overlap})} ở mức độ đáng kể. Quan sát này hàm ý rằng sự khác biệt về ước lượng điểm giữa các mô hình cần phải được xem xét song hành với biên độ biến thiên nội tại của tập dữ liệu, thay vì chỉ dựa vào các trị số trung bình đơn thuần để đưa ra kết luận.

\begin{table}[htbp]
  \centering
  \small
  \caption{Chỉ số ROC--AUC trên tập kiểm thử (\texttt{split\_seed=0}) và khoảng tin cậy 95\% ước lượng bằng Bootstrap ($B=1000$).}
  \label{tab:results-auc-bootstrap}
  \begin{tabular}{lcc}
    \toprule
    Mô hình & ROC--AUC (Lượt 0) & 95\% CI Bootstrap \\
    \midrule
    Hồi quy Logistic    & $0,870$ & $[0,810\,;\,0,922]$ \\
    Rừng ngẫu nhiên     & $0,858$ & $[0,796\,;\,0,913]$ \\
    SVM (nhân RBF)      & $0,855$ & $[0,788\,;\,0,917]$ \\
    KNN                 & $0,794$ & $[0,714\,;\,0,864]$ \\
    \bottomrule
  \end{tabular}
\end{table}

\noindent\textbf{Biện giải kết quả.} Việc các khoảng CI Bootstrap của các mô hình LR, RF và SVM chồng lấn lên nhau cho thấy ưu thế về chỉ số ROC--AUC giữa các phương pháp này có thể không thực sự ổn định trước những thay đổi của mẫu kiểm thử. Trong các tình huống thực tiễn đòi hỏi sự vận hành bền vững, kết quả phân tích này ủng hộ việc lựa chọn mô hình dựa trên \textbf{hồ sơ năng lực tổng thể}, bao gồm hiệu năng kỳ vọng, độ bất định và khả năng diễn giải, thay vì chỉ tập trung tối ưu hóa cho một chỉ số đơn lẻ.

\textbf{Độ lớn hiệu ứng Cliff's delta trên chỉ số F1.} Bảng~\ref{tab:cliffs-delta} báo cáo chỉ số Cliff's delta ($\hat{\delta}$) giữa các cặp mô hình dựa trên tập hợp giá trị F1 thu được từ 05 lượt phân đoạn dữ liệu độc lập (chi tiết tại Mục~\ref{subsec:repro}). Kết quả đối chiếu giữa RF với LR và giữa RF với SVM ghi nhận giá trị tuyệt đối $|\hat{\delta}|$ ở mức \textbf{lớn} theo các ngưỡng diễn giải chuẩn, trong khi đó sự khác biệt giữa RF và KNN chỉ đạt mức \textbf{nhỏ} với $\hat{\delta}=0,20$. Kết quả này hoàn toàn nhất quán với việc kiểm định Wilcoxon không cho thấy ý nghĩa thống kê giữa RF và KNN như đã trình bày tại Mục~\ref{subsec:tablesfigures}. Ngược lại, phép so sánh giữa LR và KNN ghi nhận $\hat{\delta}=-0,24$, tương ứng với mức độ ảnh hưởng nhỏ. Theo định nghĩa, điều này đồng nghĩa với việc F1-score của KNN có xu hướng chiếm ưu thế hơn so với LR trong một số lượt chia dữ liệu nhất định, điều này cũng phù hợp với sự tiệm cận về giá trị F1 trung bình giữa hai mô hình này.

\textbf{Ý nghĩa thực tiễn của độ lớn hiệu ứng.} Dưới góc độ triển khai thực tế, những cặp mô hình có độ lớn hiệu ứng chỉ ở mức \textbf{nhỏ} thường khó tạo ra sự khác biệt rõ rệt trong quá trình vận hành, vì vậy các mô hình sở hữu độ ổn định cao hoặc có khả năng kiểm toán tốt nên được ưu tiên lựa chọn. Ngược lại, trong trường hợp mức độ ảnh hưởng đạt ngưỡng \textbf{vừa hoặc lớn}, các nhà nghiên cứu sẽ có cơ sở khoa học vững chãi hơn để cân nhắc việc chuyển đổi thuật toán. Tuy nhiên, quyết định lựa chọn sau cùng vẫn cần được đối chiếu thận trọng với các yêu cầu về khả năng diễn giải, chi phí sai số đồng thời xem xét biên độ biến thiên của mô hình qua các lượt phân đoạn dữ liệu khác nhau.

\begin{table}[htbp]
  \centering
  \small
  \caption{Chỉ số Cliff's delta (dựa trên F1-score của 05 lượt phân đoạn dữ liệu). Giá trị $\hat{\delta} > 0$ cho thấy mô hình tại cột bên trái có xu hướng đạt F1-score vượt trội so với mô hình bên phải.}
  \label{tab:cliffs-delta}
  \begin{tabular}{lcc}
    \toprule
    Cặp đối chiếu & $\hat{\delta}$ & Mức độ ảnh hưởng ($|\hat{\delta}|$) \\
    \midrule
    RF vs LR   & $0,52$ & Lớn \\
    RF vs KNN  & $0,20$ & Nhỏ \\
    RF vs SVM  & $0,52$ & Lớn \\
    LR vs SVM  & $0,48$ & Lớn \\
    LR vs KNN  & $-0,24$ & Nhỏ \\
    SVM vs KNN & $-0,36$ & Vừa \\
    \bottomrule
  \end{tabular}
\end{table}

\section{Phân tích tầm quan trọng đặc trưng}
\label{subsec:permimp}

Mục này trình bày các kết quả phân tích diễn giải \textbf{phi mô hình} (\textit{model-agnostic}) mang tính bổ trợ như đã mô tả tại Mục~\ref{subsec:interpret}. Đối với mỗi lượt phân đoạn dữ liệu và từng mô hình đã được tối ưu hóa, kỹ thuật \texttt{permutation\_importance} được áp dụng nhằm định lượng mức độ \textbf{suy giảm chỉ số F1} trên tập kiểm thử khi thực hiện hoán vị ngẫu nhiên từng biến đầu vào với $n_{\mathrm{rep}}=30$. Các kết quả sau đó được \textbf{tổng hợp giá trị trung bình qua 05 lượt phân đoạn} độc lập và được đặc tả chi tiết cho từng thuật toán tại Hình~\ref{fig:perm-importance}.

Dựa trên kết quả tổng hợp từ tất cả các mô hình và lượt phân đoạn, hai đặc trưng \textbf{Glucose} và \textbf{BMI} ghi nhận tầm quan trọng hoán vị trung bình cao nhất, theo sau là các biến \textbf{Age} và \textbf{Insulin} ở mức độ thấp hơn. Trong khi đó, một số biến số khác như Huyết áp hoặc Độ dày nếp gấp da cho kết quả trung bình tiệm cận mức $0$ hoặc thậm chí đạt \textbf{giá trị âm nhỏ}. Hiện tượng này thường phát sinh trong bối cảnh quy mô tập kiểm thử còn hạn chế hoặc khi việc hoán vị vô tình loại bỏ nhiễu ngẫu nhiên thay vì phá vỡ các cấu trúc thông tin hữu ích, do đó các kết quả này cần được đối chiếu thận trọng với hệ số hồi quy của mô hình LR cũng như tầm quan trọng dựa trên độ tạp chất (\textit{impurity}) của RF. Đáng chú ý là thứ hạng của các biến số chủ đạo bao gồm Glucose và BMI luôn duy trì tính \textbf{ổn định} qua các phương thức phân đoạn dữ liệu khác nhau, điều này minh chứng rằng đây là các tín hiệu dự báo thực chất thay vì là hệ quả ngẫu nhiên của một phương án phân hoạch dữ liệu cụ thể nào.

\noindent\textbf{Đối chiếu thực nghiệm.} Kết quả phân tích thứ hạng dựa trên độ không thuần nhất (\textit{impurity}) của mô hình RF và các hệ số hồi quy của mô hình LR trên lượt phân đoạn đại diện (RQ2) ghi nhận sự \textbf{nhất quán} cao với phương pháp hoán vị đối với các biến số then chốt. Những sai lệch nhỏ giữa các phép đo được xác định là hệ quả tất yếu từ đặc tính thiên kiến vốn có của phương pháp \textit{impurity} đối với các biến có độ tương quan cao hoặc sở hữu nhiều điểm cắt dữ liệu.

\begin{table}[htbp]
  \centering
  \small
  \caption{Tầm quan trọng hoán vị đồng thuận (Trung bình mức suy giảm F1 trên bốn mô hình qua 05 lượt phân đoạn dữ liệu).}
  \label{tab:perm-consensus}
  \begin{tabular}{lcc}
    \toprule
    Đặc trưng & Giá trị trung bình & Độ lệch chuẩn (qua các hạt giống) \\
    \midrule
    Glucose & $0,168$ & $0,018$ \\
    BMI     & $0,038$ & $0,020$ \\
    Age     & $0,012$ & $0,015$ \\
    Insulin & $0,005$ & $0,006$ \\
    \bottomrule
  \end{tabular}
\end{table}

\begin{figure}[htbp]
  \centering
  \includegraphics[width=\textwidth]{image/permutation_importance_mean5splits.png}
  \caption{Mức suy giảm F1 trung bình do hoán vị đặc trưng (\textit{Permutation Importance}, $n_{\mathrm{rep}}=30$), đo lường trên \textbf{tập kiểm thử} và tổng hợp qua 05 lượt phân đoạn dữ liệu huấn luyện/kiểm thử.}
  \label{fig:perm-importance}
\end{figure}

\section{Phân tích thành phần}
\label{subsec:ablation}

Mục này đóng vai trò bổ trợ cho các kết quả thực nghiệm chính tại Mục~\ref{subsec:tablesfigures} bằng cách \textbf{kiểm chứng thực nghiệm} vai trò của các thành phần trong luồng xử lý (\textit{pipeline}), cụ thể là bước \textbf{chuẩn hóa đặc trưng}. Thông qua đó, nghiên cứu tái khẳng định cơ sở \textbf{lý thuyết và phương pháp luận} của quy trình đã thiết lập tại Chương~\ref{sec:methods}. Mục tiêu cốt lõi của phân tích này không nhằm lặp lại toàn bộ quá trình tối ưu hóa, mà để xác nhận rằng hiệu năng của mô hình thực sự xuất phát từ một giao thức hệ thống thay vì phụ thuộc vào các yếu tố ngẫu nhiên đơn lẻ.

\textbf{Thực nghiệm đối chứng với chuẩn hóa dữ liệu.} Trên cùng một lượt phân đoạn dữ liệu (\texttt{split\_seed=0}), nghiên cứu thực hiện so sánh hai cấu hình Hồi quy Logistic với tham số điều chuẩn cố định \texttt{C=1,0}, bao gồm: (i) luồng xử lý đầy đủ với các bước gán giá trị thiếu và chuẩn hóa bằng \texttt{StandardScaler}; và (ii) cấu hình lược bỏ hoàn toàn bước chuẩn hóa. Kết quả ghi nhận trên tập kiểm thử dựa trên tệp \texttt{ablation\_lr\_scaler.csv} cho thấy chỉ số F1-score đạt mức \textbf{tương đồng} gần như tuyệt đối với giá trị xấp xỉ $0,624$. Kết quả này phản ánh rằng, xét riêng với mức điều chuẩn cụ thể nói trên, bước chuẩn hóa Z-score không làm thay đổi trực tiếp hiệu năng của mô hình LR. Tuy nhiên, điều này \textbf{không đồng nghĩa} với việc khẳng định chuẩn hóa là bước dư thừa, bởi khi tham số \texttt{C} được tối ưu hóa thông qua \texttt{GridSearchCV} (chi tiết tại Phụ lục~\ref{app:hyper}), sự tương tác giữa thang đo đặc trưng và cường độ điều chuẩn sẽ trở nên phức tạp cũng như mang tính quyết định hơn đối với kết quả cuối cùng.

\textbf{Đặc thù của các mô hình nhạy cảm với thang đo.} Đối với các thuật toán như KNN và SVM (nhân RBF), việc chuẩn hóa dữ liệu là một \textbf{yếu tố tiên quyết} về mặt lý thuyết, bởi các phương pháp này xác lập ranh giới quyết định dựa trên hàm khoảng cách hoặc quy trình tối ưu hóa lề (\textit{margin}) như đã trình bày tại Mục~\ref{subsec:preprocess}. Chính vì vậy, nghiên cứu chủ đích \textbf{không} thực hiện đối chứng việc lược bỏ bước chuẩn hóa cho các mô hình này do hiệu năng chắc chắn sẽ suy giảm nghiêm trọng và không phản ánh đúng năng lực thực chất của thuật toán. Theo lý thuyết học máy, \textbf{KNN và SVM} thuộc nhóm mô hình \textbf{nhạy cảm đặc biệt} với biên độ của đặc trưng, do đó được dự báo sẽ \textbf{chịu mức sụt giảm hiệu năng lớn nhất} nếu thiếu bước chuẩn hóa, trong khi các mô hình như RF và LR thường có mức độ nhạy cảm thấp hơn. Lập luận này củng cố vững chắc nền tảng lý thuyết cho các kết quả thực nghiệm đã được công bố.

\textbf{Vai trò của tối ưu hóa siêu tham số.} Mọi mô hình được trình bày tại các Bảng~\ref{tab:results-point} đến Bảng~\ref{tab:results-meanstd} đều đã trải qua quy trình tinh chỉnh nghiêm ngặt thông qua \texttt{GridSearchCV}. Việc tối ưu hóa này được xem là một \textbf{thành phần cấu thành bắt buộc} nhằm đảm bảo tính đối chiếu công bằng giữa các thuật toán, bởi việc sử dụng tham số mặc định thường dẫn đến kết quả thiếu ổn định và gây khó khăn cho quá trình tái lập. \emph{Xét về mặt phương pháp luận:} việc bỏ qua bước xác thực chéo nội bộ có thể làm \textbf{gia tăng phương sai} của các ước lượng hiệu năng, đồng thời gây ra sự \textbf{dao động thứ hạng} không mong muốn giữa các lượt phân đoạn dữ liệu khác nhau. Nhận định này được minh chứng rõ nét khi đối chiếu tính ổn định giữa kết quả thu được từ một lượt chạy đơn lẻ tại Bảng~\ref{tab:results-point} và kết quả tổng hợp qua năm hạt giống ngẫu nhiên tại Bảng~\ref{tab:results-meanstd}.

\textbf{Giao thức thực nghiệm đa hạt giống.} Sự khác biệt rõ rệt giữa kết quả chạy đơn lẻ và kết quả tổng hợp chính là minh chứng trực quan cho mức độ \textbf{biến thiên} của mô hình. Đây được coi là cơ sở thực nghiệm quan trọng nhằm cảnh báo những rủi ro khi đưa ra kết luận chỉ dựa trên \textbf{một} lượt phân đoạn dữ liệu ngẫu nhiên, vốn rất dễ dẫn đến hiện tượng \textbf{phóng đại hiệu năng hoặc làm đảo ngược thứ hạng} thực chất của các thuật toán.

\textbf{Đường cơ sở tham chiếu (\textit{Baseline}).} Mô hình \texttt{DummyClassifier} đóng vai trò là \textbf{mốc tham chiếu tối thiểu}, từ đó cung cấp ngữ cảnh cần thiết để định lượng chính xác giá trị thực chất của các chỉ số Precision, Recall và F1-score (như đã trình bày tại Mục~\ref{subsec:models}).

\noindent\textit{Tóm lại,} các phân tích thành phần đã làm nổi bật một \textbf{chuỗi hệ quả phương pháp luận} chặt chẽ, trong đó việc chuẩn hóa đặc trưng là điều kiện tiên quyết cho các mô hình KNN và SVM, đồng thời quy trình xác thực chéo nội bộ giúp kiểm soát hiệu quả phương sai siêu tham số. Bên cạnh đó, việc lặp lại thực nghiệm trên nhiều hạt giống ngẫu nhiên giúp loại bỏ các sai lệch về thứ hạng do yếu tố ngẫu nhiên gây ra. Những thành phần này khi kết hợp cùng mô hình tham chiếu sẽ tạo nên một khung đánh giá minh bạch, nhất quán và sở hữu khả năng tái lập cao.

\textbf{Tổng kết Chương 4.} Tổng hợp các dữ liệu thực nghiệm thu được cho thấy một sự đánh đổi (\textit{trade-off}) rõ nét giữa \textbf{hiệu năng tối ưu} và \textbf{độ ổn định}. Trong khi mô hình Rừng ngẫu nhiên thường chiếm ưu thế về các chỉ số phân loại thì Hồi quy Logistic lại mang đến sự bền vững cùng khả năng diễn giải trực tiếp. Chính vì vậy, nghiên cứu khẳng định không tồn tại một mô hình tối ưu tuyệt đối cho mọi mục tiêu mà việc lựa chọn sau cùng cần căn cứ vào ưu tiên của từng ứng dụng cụ thể, dù đó là tối đa hóa năng lực dự báo hay ưu tiên tính ổn định và sự minh bạch trong thực hành lâm sàng. Chương Thảo luận tiếp theo sẽ tiếp tục mở rộng các hàm ý này dưới góc độ chuyên môn y sinh sâu sắc hơn.

\chapter{Thảo luận}
\label{sec:discussion}

\noindent Chương này tập trung vào việc \textbf{biện giải các cơ chế phương pháp luận} đứng sau những quan sát thực nghiệm, thay vì chỉ dừng lại ở việc tóm lược các số liệu định lượng đơn thuần. Trục phân tích chính của chương sẽ xoay quanh ba vấn đề cốt lõi, bao gồm: \textbf{(i) Nguyên nhân} dẫn đến sự khác biệt về đặc tính vận hành giữa các mô hình; \textbf{(ii) Ngữ cảnh ứng dụng} tối ưu cho từng thuật toán dựa trên đặc thù của dữ liệu; và cuối cùng là \textbf{(iii) Sự đánh đổi} tất yếu giữa năng lực dự báo, độ ổn định cùng khả năng diễn giải mô hình.

\section{Biện giải và phân tích kết quả}
\label{subsec:disc-results}

Các luận điểm được trình bày dưới đây nhằm thiết lập mối liên kết trực tiếp giữa các kết quả thực nghiệm tại Chương~\ref{sec:results} và những giả định của lý thuyết học thống kê. Mục tiêu trọng tâm là chuyển dịch từ việc đơn thuần ``xếp hạng chỉ số'' sang việc ``giải thích cơ chế'' vận hành trong cùng một giao thức thực nghiệm thống nhất. Theo đó, mạch lập luận sẽ được triển khai theo trình tự logic: từ \textbf{Quan sát thực nghiệm} đến \textbf{Biện giải cơ chế}, từ đó rút ra các \textbf{Hàm ý ứng dụng} và xác định các \textbf{Giới hạn tổng quát hóa} của nghiên cứu.

\textbf{Đặc tính mô hình và sự khác biệt hóa về hiệu năng.}
Mô hình \textbf{Rừng ngẫu nhiên (RF)} đạt chỉ số F1-score trung bình cao nhất (Bảng~\ref{tab:results-meanstd}), qua đó phản ánh ưu thế của kiến trúc \textbf{tổ hợp cây} (\textit{ensemble learning}) trong việc khai thác các \textbf{tương tác phi tuyến} giữa các đặc trưng đầu vào. Trên thực tế đối với tập dữ liệu Pima, mối tương quan giữa các biến số lâm sàng bao gồm đường huyết, BMI, độ tuổi và tiền sử gia đình vốn khó có thể được đặc tả trọn vẹn bằng một ranh giới tuyến tính đơn giản. Cơ chế phân tách dữ liệu theo cấu trúc phân cấp của cây quyết định cho phép RF nắm bắt được các tương tác phức tạp này, từ đó cải thiện đáng kể độ nhạy (\textit{Recall}) và tối ưu hóa chỉ số F1-score. \emph{Điểm mấu chốt nằm ở chỗ:} cả hệ số hồi quy của LR và chỉ số \textit{Permutation Importance} đều xác định \textbf{Glucose} cùng với \textbf{BMI} là những biến số quan trọng nhất. RF sở hữu khả năng mô hình hóa các \textbf{tổ hợp điều kiện}, chẳng hạn như việc xác định vùng rủi ro khi cả Glucose và BMI cùng vượt ngưỡng, điều mà một mặt phẳng \textit{log-odds} tuyến tính khó có thể biểu diễn đầy đủ. Đây chính là lý do giải thích vì sao RF thường đạt hiệu năng F1-score nhỉnh hơn, mặc dù mô hình này nhạy cảm hơn trước những biến thiên của dữ liệu.

Ngược lại, \textbf{Hồi quy Logistic (LR)} vận hành dựa trên giả định về mối quan hệ \textbf{tuyến tính trong không gian log-odds}. Mặc dù giả định này có thể hạn chế khả năng biểu diễn các mẫu dữ liệu phức tạp, song nó lại thiết lập được một \textbf{thiên kiến cảm dẫn} (\textit{inductive bias}) mạnh mẽ, giúp kiểm soát phương sai hiệu quả trong điều kiện cỡ mẫu hạn chế. Đặc tính này củng cố cho quan sát thực nghiệm rằng LR đạt Accuracy và ROC--AUC trung bình cao nhất (Bảng~\ref{tab:results-meanstd}). Khả năng \textbf{phân hạng tổng thể} (\textit{ranking}) của LR tỏ ra rất ổn định trên toàn bộ phổ xác suất dự báo, đồng thời độ lệch chuẩn của chỉ số F1 ở LR thấp hơn đáng kể so với RF ($0,017$ so với $0,025$), minh chứng cho \textbf{tính ổn định} cao hơn trước các biến thiên của phân đoạn dữ liệu. Như vậy, kết quả thực nghiệm cho thấy một \textbf{sự đánh đổi} điển hình, đó là RF chiếm ưu thế về F1-score nhờ cấu trúc phi tuyến linh hoạt, trong khi LR lại vượt trội về AUC và độ ổn định nhờ kiến trúc mô hình tinh gọn. Đây có thể coi là một mẫu hình đặc trưng đối với các tập dữ liệu y sinh dạng bảng có quy mô vừa và nhỏ.

\textbf{Sự đánh đổi thiên kiến và phương sai (\textit{Bias--Variance Tradeoff}).}
Sự khác biệt về hiệu năng có thể được giải thích một cách rõ nét qua lăng kính \textbf{Bias--Variance}, trong đó mô hình LR có thiên kiến cao hơn nhưng phương sai lại thấp hơn, do đó thường tỏ ra \textbf{vững chãi} (\textit{robust}) hơn trước các nhiễu dữ liệu. Ngược lại, RF sở hữu tính linh hoạt cao và có khả năng khớp tốt với các đặc trưng phi tuyến, song mô hình này lại phụ thuộc nhiều vào các phép phân tách mang tính cục bộ, dẫn đến sự dao động mạnh hơn giữa các lượt thực thi. Kết quả ghi nhận qua 05 hạt giống ngẫu nhiên cho thấy mặc dù RF đạt được các đỉnh hiệu năng cao hơn, nhưng lại thiếu đi tính ổn định khi so sánh với LR. Điều này khẳng định sự đánh đổi tất yếu giữa hiệu năng kỳ vọng và \textbf{khả năng tái lập} (\textit{reproducibility}) của các kết quả thực nghiệm.

\textbf{Ngữ cảnh ưu tiên mô hình.}
Việc lựa chọn thuật toán cần phải gắn liền với \textbf{ngữ cảnh ứng dụng} lâm sàng thay vì chỉ dựa vào một chỉ số đơn lẻ. Cụ thể, RF là lựa chọn phù hợp khi mục tiêu ưu tiên là tối đa hóa F1-score hoặc Recall trên lớp dương tính đồng thời hệ thống có khả năng chấp nhận một mô hình có độ phức tạp cao. Trong khi đó, LR lại là phương án tối ưu khi yêu cầu đặt nặng \textbf{tính minh bạch} của các hệ số, độ bền vững qua các lần phân đoạn dữ liệu và một quy trình kiểm chứng đơn giản. Các hàm ý về triển khai thực tế sẽ được làm rõ chi tiết hơn tại Mục~\ref{subsec:disc-implications}.

Mặc dù RF dẫn đầu về chỉ số F1-score trung bình, song \textbf{chênh lệch hiệu năng thực tế giữa các mô hình là tương đối nhỏ}. Kết quả kiểm định \textbf{Wilcoxon signed-rank} trên chỉ số F1 giữa RF và KNN cho giá trị $p \approx 0,44$ (như đã nêu tại Mục~\ref{subsec:tablesfigures}), điều này cho thấy \textbf{không} có đủ cơ sở thống kê để khẳng định một thuật toán vượt trội hoàn toàn ở mức ý nghĩa $0,05$. \textbf{Lưu ý về phương pháp luận:} với quy mô $n=5$ cặp mẫu, \textbf{lực thống kê} (\textit{statistical power}) của phép kiểm định đang ở mức thấp. Do đó, việc thiếu ý nghĩa thống kê \textbf{không} đồng nghĩa với việc các mô hình có hiệu năng tương đương, mà chỉ chỉ ra rằng bằng chứng thực nghiệm hiện tại \textbf{chưa đủ mạnh} để bác bỏ giả thuyết không ($H_0$). Quan sát này củng cố thêm nhận định rằng ưu thế ghi nhận được chịu ảnh hưởng đáng kể từ \textbf{biến thiên chọn mẫu} và các sai số ước lượng phát sinh khi cỡ mẫu còn hạn chế.

Đối với các mô hình còn lại, \textbf{KNN} vận hành dựa trên cơ chế \textbf{lân cận cục bộ}, do đó thuật toán này cực kỳ nhạy cảm với việc chuẩn hóa thang đo, mật độ dữ liệu cũng như nhiễu nhãn tại các vùng biên quyết định. Trong điều kiện quy mô mẫu hạn chế, chỉ một thay đổi nhỏ của các điểm dữ liệu kiểm thử theo hạt giống (\textit{seed}) cũng có thể làm thay đổi hoàn toàn cấu trúc láng giềng đa số, từ đó dẫn đến hiện tượng \textbf{dao động} hiệu năng mạnh mẽ như đã ghi nhận tại Bảng~\ref{tab:results-meanstd}. Trong khi đó, \textbf{SVM (nhân RBF)} thực hiện việc tối ưu hóa \textbf{biên phân tách lớp} trong không gian nhân (\textit{kernel space}). Hiệu năng của thuật toán này phụ thuộc chặt chẽ vào việc tinh chỉnh cặp siêu tham số $(C, \gamma)$ cùng cấu trúc dữ liệu sau khi ánh xạ, điều này vô tình gây ra sự biến động lớn qua các lượt phân đoạn dữ liệu khác nhau. Trong một luồng xử lý (\textit{pipeline}) thống nhất, mặc dù hai mô hình này cho thấy tính cạnh tranh trên một số lượt chạy cụ thể tại Bảng~\ref{tab:results-point}, song thứ hạng trung bình của chúng thường \textbf{kém ổn định} hơn đáng kể so với mô hình LR. Kết quả này hoàn toàn phù hợp với kỳ vọng về mức phương sai cao của các mô hình phi tuyến khi thực hiện trên dữ liệu quy mô nhỏ.

\textbf{Trả lời RQ1.} Xét theo chỉ số $F_1$-score trung bình, mô hình RF chiếm ưu thế dẫn đầu với giá trị $0{,}614 \pm 0{,}025$, theo sau lần lượt là KNN và LR; trong khi đó, SVM (RBF) lại ghi nhận kết quả thấp nhất trong không gian tham số khảo sát. Ngược lại, các chỉ số Accuracy và ROC--AUC trung bình cao nhất lại thuộc về mô hình LR với các giá trị tương ứng là $0{,}748 \pm 0{,}018$ và $0{,}824 \pm 0{,}031$. Tại lượt phân đoạn đại diện được trình bày ở Bảng~\ref{tab:results-point}, SVM và RF đạt chỉ số $F_1$ tương đương nhau ở mức xấp xỉ $0{,}646$ và đều vượt trội hơn so với LR. Những kết quả này là minh chứng rõ nét cho sự \textbf{đánh đổi} giữa các tiêu chí đánh giá khác nhau cũng như phản ánh \textbf{độ nhạy của thuật toán đối với phương thức phân đoạn mẫu} được áp dụng.

\textbf{Trả lời RQ2.} Dựa trên các mô hình đã huấn luyện với thiết lập hạt giống số 0 (\texttt{split\_seed=0}), hai biến số \textbf{Glucose} và \textbf{BMI} thể hiện tầm quan trọng vượt trội, đồng thời được xác nhận một cách thống nhất qua hệ số hồi quy của mô hình LR cũng như tầm quan trọng đặc trưng của mô hình RF. Kết quả phân tích \textbf{Permutation Importance} trung bình qua 05 lượt thực thi tại Mục~\ref{subsec:permimp} tiếp tục củng cố vai trò của Glucose và BMI là hai yếu tố có tác động mạnh mẽ nhất theo phương pháp đo lường phi mô hình (\textit{agnostic}). Thứ tự ưu tiên này hoàn toàn nhất quán với các bằng chứng lâm sàng hiện nay về nồng độ đường huyết và tình trạng béo phì trong quy trình tầm soát đái tháo đường loại 2. Tuy nhiên, cần lưu ý rằng các kết quả này chỉ mang tính chất diễn giải bổ trợ và không đại diện cho mối quan hệ nhân quả trực tiếp \cite{kirbas2025shap,chang2022pima}.

\textbf{Năng lực phân tách và hiệu chuẩn xác suất.} Trong thực nghiệm nghiên cứu, việc phân biệt rạch ròi giữa \textbf{năng lực phân tách} (\textit{discrimination}) và \textbf{độ hiệu chuẩn xác suất} (\textit{calibration}) là thực sự cần thiết. Xét theo lý thuyết học máy, các mô hình có ranh giới quyết định linh hoạt thường đạt được chỉ số $F_1$ cao nhưng xác suất dự báo lại có thể kém hiệu chuẩn, trong khi các mô hình tuyến tính thường được ưu tiên nhờ cung cấp xác suất dự báo dễ diễn giải hơn. Trong nghiên cứu này, mô hình RF không chỉ dẫn đầu về giá trị $F_1$ trung bình mà còn đạt được chỉ số Brier thấp nhất và theo sát ngay sau đó là mô hình LR. Kết quả này cho thấy hai tiêu chí nêu trên không nhất thiết phải mâu thuẫn với nhau nhưng cũng \textbf{không} mặc định có sự tương quan thuận, vì vậy việc báo cáo đồng thời chỉ số Brier và đường hiệu chuẩn tại Mục~\ref{subsec:calibration} là yêu cầu bắt buộc khi sử dụng xác suất để phân tầng nguy cơ. Trong trường hợp ưu tiên tính minh bạch của các tham số cùng sự bền vững của hệ thống, mô hình LR vẫn là lựa chọn tối ưu mặc dù có hiệu năng $F_1$ thấp hơn so với RF.

\textbf{Tính ổn định của hiệu năng.} Việc bổ sung chỉ số PR--AUC tại Mục~\ref{subsec:prauc} đã giúp làm rõ năng lực dự báo trên lớp dương tính trong bối cảnh dữ liệu có hiện tượng mất cân bằng nhẹ, trong đó mô hình RF và LR tiếp tục khẳng định vị thế dẫn đầu. Kết quả phân tích \textbf{khoảng tin cậy (CI) Bootstrap} cho chỉ số ROC--AUC tại Bảng~\ref{tab:results-auc-bootstrap} cho thấy sự chồng lấn (\textit{overlap}) ở mức độ đáng kể giữa các mô hình mạnh. Quan sát này là lời nhắc nhở quan trọng rằng những chênh lệch về điểm số trung bình \textbf{không} đồng nghĩa với ưu thế tuyệt đối, nhất là khi biên độ sai số vẫn còn ở mức rộng. Bên cạnh đó, chỉ số \textbf{Cliff's delta} tại Bảng~\ref{tab:cliffs-delta} xác nhận một số cặp mô hình có độ lớn hiệu ứng (\textit{effect size}) đạt mức ``lớn'', trong khi đối chiếu giữa RF và KNN chỉ dừng lại ở mức ``nhỏ''. Kết quả này hoàn toàn tương ứng với việc kiểm định Wilcoxon không tìm thấy khác biệt có ý nghĩa thống kê trên chỉ số $F_1$. Nhìn chung, việc phân tích song song giữa \emph{ý nghĩa thống kê}, \emph{khoảng tin cậy} và \emph{độ lớn hiệu ứng} là yêu cầu bắt buộc nhằm tránh việc phóng đại những cải thiện ở biên độ hẹp trên các bộ dữ liệu có cỡ mẫu hạn chế.

\textbf{Hàm ý phương pháp luận.} Từ các bằng chứng thực nghiệm về sự biến thiên theo phân đoạn dữ liệu, sự chồng lấn của khoảng tin cậy cùng độ lớn hiệu ứng thấp ở một số cặp đối chiếu, nghiên cứu rút ra hàm ý then chốt là khi làm việc với dữ liệu y sinh quy mô nhỏ, việc thiết lập một \textbf{giao thức đánh giá vững chãi} (\textit{robust evaluation protocol}) có tầm quan trọng lớn hơn nhiều so với việc theo đuổi những sai khác nhỏ về mặt điểm số thuần túy.

\textbf{Giá trị $p$ và ý nghĩa thực tiễn.} Một nguyên tắc phương pháp luận trọng yếu cần được lưu ý chính là sự khác biệt có ý nghĩa thống kê không phải lúc nào cũng đồng nhất với sự khác biệt có giá trị trong ứng dụng thực tiễn. Do đó, quy trình xác định mô hình tối ưu trong nghiên cứu này dựa trên sự tổng hòa của các bằng chứng thực nghiệm bao gồm giá trị $p$, độ lớn hiệu ứng thực tế cùng độ ổn định qua các lượt phân đoạn dữ liệu, thay vì chỉ phụ thuộc vào một chỉ số đơn lẻ duy nhất.

\textbf{Tính bền vững của thứ hạng đặc trưng.} Việc bổ sung phân tích \textbf{tầm quan trọng hoán vị} tại Mục~\ref{subsec:permimp} đóng vai trò quan trọng trong việc củng cố tính khách quan cho phần diễn giải mô hình. Đây là một phương pháp tiếp cận \textbf{phi mô hình} (\textit{agnostic}) nên sẽ ít phụ thuộc vào cấu trúc thuật toán hơn so với chỉ số \textit{impurity} của RF, từ đó giúp giảm thiểu rủi ro sai lệch khi đánh giá tầm quan trọng của một biến số do đặc tính chia nhánh đặc thù của cây quyết định. Sự nhất quán ghi nhận được giữa hệ số hồi quy của LR, độ quan trọng đặc trưng của RF cùng kết quả hoán vị về vai trò chủ đạo của \textbf{Glucose} và \textbf{BMI} đã gia tăng đáng kể độ tin cậy định tính cho câu hỏi nghiên cứu RQ2. Bên cạnh đó, những sai lệch nhỏ giữa phương pháp \textit{impurity} và hoán vị thường phát sinh trong trường hợp các biến số có sự tương quan cao, đây cũng chính là lý do nghiên cứu khuyến nghị nên tiếp cận theo hướng \textbf{đa phương pháp} khi thực hiện diễn giải dữ liệu y sinh.

\textbf{Phân tích sai số và ranh giới quyết định.} Các phân tích định tính cho thấy những lỗi \textbf{âm tính giả} (\textit{false negative}) thường tập trung tại các \textbf{trường hợp cận biên} (\textit{borderline cases}), cụ thể là những nơi mà các đặc trưng nằm rất sát với ranh giới quyết định. Tại các vùng dữ liệu này, một mô hình có độ linh hoạt cao như RF thể hiện khả năng thích nghi tốt hơn so với ranh giới tuyến tính có phần cứng nhắc của LR, điều này dẫn đến việc chỉ số \textit{recall} (độ nhạy) của RF thường ưu thế hơn trên một số phân đoạn dữ liệu nhất định. Tuy nhiên, chính sự linh hoạt này đôi khi lại làm gia tăng tỷ lệ \textbf{dương tính giả} (\textit{false positive}), phản ánh sự đánh đổi kinh điển giữa độ nhạy và độ đặc hiệu trong bài toán phân loại. Trong khi đó, mô hình LR với ranh giới mượt hơn trong không gian \textit{log-odds} thường có xu hướng ít cực đoan hơn tại các vùng có nhiều nhiễu, song lại có thể bỏ sót những tương tác phi tuyến tinh vi giữa các biến số. Chi tiết về ma trận nhầm lẫn trên một lượt phân đoạn đại diện đã được trình bày cụ thể tại Mục~\ref{subsec:disc-error-related} và Hình~\ref{fig:roc-confusion}.

\textbf{Hệ quả trong tầm soát y tế.} Trong bối cảnh tầm soát lâm sàng, sai số \textbf{âm tính giả} dẫn đến việc bỏ sót bệnh nhân thường mang lại hệ quả nghiêm trọng, bởi điều này có thể gây ra tình trạng \textbf{trì hoãn can thiệp} và làm gia tăng nguy cơ biến chứng dài hạn. Ngược lại, sai số \textbf{dương tính giả} chủ yếu gây ra các \textbf{xét nghiệm bổ sung cùng gánh nặng về chi phí và tâm lý} không cần thiết cho người bệnh. Chính vì vậy, việc đánh giá các chỉ số FP và FN không chỉ đơn thuần là phân tích số liệu thống kê mà phải luôn gắn liền với \textbf{hàm chi phí mục tiêu} của từng chương trình tầm soát cụ thể. Bên cạnh đó, cần tái khẳng định rằng các phân tích trong luận văn này chủ yếu mang tính minh chứng về mặt phương pháp luận, hoàn toàn không có giá trị thay thế cho các quy trình chẩn đoán lâm sàng chuyên môn.

\textbf{Tổng kết nội dung thảo luận.} Kết quả nghiên cứu cho thấy không có mô hình đơn lẻ nào chiếm ưu thế tuyệt đối trên mọi phương diện. Trong khi RF tỏ ra ưu việt về năng lực phân loại thông qua các chỉ số $F_1$ và PR--AUC, thì LR lại mang lại độ bền vững cao, khả năng phân hạng xác suất (AUC) tốt cùng cấu trúc diễn giải tường minh. Cả hai mô hình này đều cần được xem xét cẩn trọng kèm theo các chỉ số hiệu chuẩn và biên độ bất định thống kê. Điều này một lần nữa tái khẳng định tầm quan trọng của việc thực hiện đánh giá \textbf{đa chiều} thay vì lệ thuộc vào một thước đo duy nhất trong nghiên cứu học máy y sinh.

\textbf{Hàm ý tổng quát hóa.} Một quan sát then chốt rút ra từ nghiên cứu là sự khác biệt về hiệu năng giữa các thuật toán luôn phụ thuộc mật thiết vào đặc thù của dữ liệu cũng như giao thức thực nghiệm được áp dụng. Do đó, mọi kết luận về một ``mô hình tối ưu'' cần được diễn giải một cách thận trọng, bởi thứ hạng giữa các phương pháp hoàn toàn có thể bị hoán đổi khi triển khai trên các tập dữ liệu hoặc thiết lập đánh giá khác nhau. Giá trị cốt lõi của nghiên cứu này không chỉ dừng lại ở kết quả định lượng mà còn nằm ở việc xây dựng thành công một quy trình đánh giá chuẩn hóa, có khả năng tái lập cao và triệt tiêu hoàn toàn hiện tượng rò rỉ dữ liệu (\textit{data leakage}), từ đó đảm bảo tính khách quan tối đa trong việc đối chiếu các phương pháp.

Xét ở phạm vi rộng hơn, mẫu hình kết quả này hoàn toàn có khả năng lặp lại trên các bộ dữ liệu y sinh dạng bảng có quy mô nhỏ khác, trong đó các mô hình linh hoạt thường đạt được đỉnh hiệu năng cao hơn, song song với đó là các mô hình đơn giản sẽ duy trì được độ ổn định tốt hơn. Tuy nhiên, năng lực chuyển giao của các mô hình này vẫn cần phải được xác nhận thông qua các bước \textbf{kiểm chứng độc lập} (\textit{external validation}) trước khi có thể thực hiện tổng quát hóa cho các quần thể bệnh nhân mới.

\section{Hàm ý thực tiễn}
\label{subsec:disc-implications}

Dựa trên các phân tích về đặc tính vận hành của mô hình tại Mục~\ref{subsec:disc-results}, trong bối cảnh không có thuật toán nào chiếm ưu thế tuyệt đối trên mọi tiêu chí, quy trình \textbf{lựa chọn mô hình} cần phải được gắn kết chặt chẽ với \textbf{mục tiêu của từng tác vụ cụ thể}. Trường hợp ưu tiên tối ưu hóa chỉ số $F_1$ cùng năng lực nắm bắt các cấu trúc phi tuyến phức tạp, Rừng ngẫu nhiên (RF) sẽ là một lựa chọn phù hợp, điều này cũng tương ứng với mô hình đã được lưu trữ dựa trên điểm $F_1$ trung bình tại tệp \texttt{resources/diabetes\_prediction\_model.pkl}. Ngược lại, nếu yêu cầu đặt nặng \textbf{tính minh bạch} và khả năng diễn giải trực tiếp thông qua hệ số của các biến số, mô hình Hồi quy Logistic (LR) sẽ tỏ ra ưu thế hơn khi vẫn duy trì được mức ROC--AUC trung bình cao nhất. Bên cạnh đó, sự chênh lệch không đáng kể giữa các thuật toán dẫn đầu cũng gợi ý rằng các yếu tố như \textbf{chi phí tính toán}, \textbf{thời gian huấn luyện} và \textbf{khả năng triển khai thực tế}, bao gồm tốc độ suy luận hay quy trình kiểm định phần mềm, nên được xem xét với trọng số tương đương với các chỉ số hiệu năng, nhất là khi các sai khác này vẫn nằm trong phạm vi sai số thống kê.

\textbf{Định hướng lựa chọn ở cấp độ ra quyết định.}
Trong kịch bản \textbf{tầm soát sớm} vốn ưu tiên việc giảm thiểu tỷ lệ \textbf{âm tính giả} thông qua việc chú trọng chỉ số \textit{Recall} hoặc Độ nhạy, các mô hình đạt hiệu năng cao như \textbf{Rừng ngẫu nhiên} nên được ưu tiên thử nghiệm chuyên sâu, đồng thời kết hợp với việc hiệu chỉnh ngưỡng quyết định theo chi phí lâm sàng thực tế. Ngược lại, khi trọng tâm được đặt vào \textbf{khả năng diễn giải, tính kiểm toán quyết định cùng độ ổn định} qua các lượt phân đoạn dữ liệu, \textbf{Hồi quy Logistic} lại là phương án \textit{đáng tin cậy hơn} nhờ sở hữu hệ số định lượng rõ ràng và chỉ số ROC--AUC dẫn đầu. Cần nhấn mạnh rằng những nhận định này không đóng vai trò như một khuyến cáo lâm sàng trực tiếp mà thực chất là một \textbf{khung phương pháp luận} nhằm gắn kết thuật toán với \textit{mục tiêu triển khai}, cụ thể là việc cân bằng giữa nhiệm vụ giảm FN và kiểm soát FP trước khi xác lập phương án thực thi cuối cùng.

Về mặt \textbf{ứng dụng y tế}, các mô hình trong nghiên cứu này hiện đóng vai trò \textbf{hỗ trợ nghiên cứu hoặc giáo dục} và không có giá trị thay thế cho các quy trình chẩn đoán chuyên khoa hiện hành. Việc triển khai vào thực tế đòi hỏi phải có quy trình xác thực độc lập, hệ thống quản trị rủi ro nghiêm ngặt cùng sự tuân thủ tuyệt đối các quy định về bảo mật dữ liệu y tế như đã đề cập tại Mục~\ref{subsec:disc-limitations}.

\textbf{Khuyến nghị thực thi.} Dựa trên các kết quả thực nghiệm thu được, nghiên cứu rút ra hai nguyên tắc cốt lõi: thứ nhất, đối với các tác vụ tầm soát ưu tiên độ phủ, người thực hiện có thể ưu tiên các thiết lập đạt chỉ số \textit{Recall} cao dù phải chấp nhận sự sụt giảm nhất định về mặt \textit{Precision}; thứ hai, trong bối cảnh cần sự minh bạch để phục vụ việc \textbf{kiểm toán quyết định}, các mô hình tuyến tính như LR sẽ là lựa chọn tối ưu nhờ khả năng diễn giải trực tiếp. Mọi quyết định lựa chọn cần dựa trên kết quả đánh giá thông qua \textbf{nhiều lần phân đoạn dữ liệu} thay vì chỉ dựa vào một lần chia tập duy nhất, nhằm giảm thiểu tối đa rủi ro từ các biến động ngẫu nhiên của dữ liệu.

Ở mức độ quy trình, nghiên cứu đề xuất cần tách biệt rõ rệt hai giai đoạn bao gồm: \textbf{lựa chọn mô hình} dựa trên giao thức tái lập và \textbf{xác lập ngưỡng quyết định} dựa trên hàm chi phí lâm sàng đặc thù của từng đơn vị, cụ thể là việc so sánh chi phí giữa sai số FN và FP. Đối với mô hình LR, các hệ số thu được còn có thể chuyển đổi thành biểu đồ định mức nguy cơ (\textit{nomogram}), từ đó hỗ trợ việc truyền thông thông tin y tế một cách minh bạch và hiệu quả trong thực hành lâm sàng.

\textbf{Hướng dẫn ra quyết định (Tóm lược):}
\begin{itemize}
  \item \textbf{Trong trường hợp ưu tiên khả năng diễn giải và tính kiểm toán:} Mô hình \textbf{Hồi quy Logistic (LR)} được khuyến nghị sử dụng nhờ cấu trúc tham số tường minh.
  \item \textbf{Trong trường hợp ưu tiên năng lực phân loại tối ưu trên lớp bệnh lý:} Nên ưu tiên mô hình \textbf{Rừng ngẫu nhiên (RF)}, đồng thời cần thiết lập các cơ chế kiểm soát chặt chẽ biên độ biến thiên giữa các lượt phân đoạn dữ liệu.
  \item \textbf{Đối với các tập dữ liệu quy mô nhỏ và có độ nhiễu cao:} Cần ưu tiên lựa chọn các cấu hình sở hữu \textbf{độ ổn định cao} trước khi tiến hành thực hiện các bước tối ưu hóa hiệu năng biên.
\end{itemize}

\section{Phân tích sai số và đối chiếu với các nghiên cứu liên quan}
\label{subsec:disc-error-related}

Dựa trên kết quả thực nghiệm với mô hình RF tại lượt phân đoạn hạt giống số 0 (\texttt{split\_seed=0}) như minh họa tại Hình~\ref{fig:roc-confusion}, ma trận nhầm lẫn ghi nhận \textbf{23 trường hợp âm tính giả} (\textit{false negative}) và \textbf{11 trường hợp dương tính giả} (\textit{false positive}), tương ứng với độ nhạy (\textit{Recall}) trên lớp dương tính đạt mức xấp xỉ $57\%$. Việc tỷ lệ FN cao hơn FP đã phản ánh một sự đánh đổi thường gặp trong bài toán \textbf{tầm soát y tế}, nơi đòi hỏi sự cân đối khắt khe giữa rủi ro bỏ sót ca bệnh và hệ quả từ các báo động giả. Xét về mặt hệ quả thực tiễn, mỗi trường hợp FN có thể dẫn đến việc \textbf{bỏ lỡ cơ hội can thiệp sớm} cho bệnh nhân, trong khi đó mỗi ca FP lại có nguy cơ gây ra những \textbf{xét nghiệm bổ sung hoặc các can thiệp y khoa không cần thiết}. Mức độ chấp nhận các loại sai số này phụ thuộc mật thiết vào ngưỡng quyết định cùng hàm chi phí mục tiêu do từng cơ sở y tế thiết lập, thay vì chỉ dựa vào các kết quả định lượng thuần túy từ mô hình. So với đường cơ sở (\textit{majority baseline}) vốn luôn dự báo lớp đa số dẫn đến chỉ số \textit{Recall} và $F_1$ bằng không, các mô hình được đề xuất đã thể hiện năng lực phân tách rõ rệt. Tuy nhiên, khi diễn giải kết quả cho các cộng đồng khác, cần đặc biệt lưu ý về \textbf{yếu tố thiên lệch và tính đặc thù quần thể} của bộ dữ liệu Pima \cite{uci-pima}.

Một hướng phân tích chuyên sâu hơn là thực hiện đặc tả \textbf{nhóm mẫu bị phân loại sai} dựa theo các cụm đặc trưng cụ thể. Theo đó, các trường hợp âm tính giả thường có xu hướng tập trung tại vùng biên của không gian đặc trưng, nơi mà tín hiệu lâm sàng giữa hai lớp không đủ mức độ tách biệt rõ ràng. Quan sát này gợi ý nhu cầu về việc bổ sung thêm các biến số mới hoặc tiến hành mô hình hóa các tương tác phức tạp hơn đối với những ca bệnh khó chẩn đoán. Ở chiều ngược lại, sự tồn tại bền vững của các mẫu dữ liệu thuộc nhóm ``khó'' này cũng phản ánh giới hạn về lượng thông tin nội tại chứa đựng trong dữ liệu quan sát, hơn là chỉ đơn thuần xuất phát từ những hạn chế của thuật toán học máy.

Về mặt kỹ thuật, việc thực hiện phân tích sai số theo nhóm, chẳng hạn như dựa trên dải xác suất dự báo hoặc theo phân vị của các biến Glucose và BMI, sẽ giúp xác định chính xác những vùng dữ liệu mà mô hình thường gặp thất bại. Từ những nhận định đó, nhà nghiên cứu có thể định hướng các biện pháp can thiệp cụ thể như thiết lập lại ngưỡng quyết định, bổ sung thêm các thuộc tính mới hoặc ưu tiên thu thập dữ liệu bổ sung cho các nhóm mẫu hiện đang nằm tại vùng ranh giới quyết định.

Khi tiến hành đối chiếu với các công trình nghiên cứu liên quan tại Bảng~\ref{tab:related-numbers}, sự chênh lệch về hiệu năng thường phản ánh tầm ảnh hưởng từ \textbf{quy trình tiền xử lý và giao thức thực nghiệm}, vốn đóng vai trò quan trọng tương đương với bản thân thuật toán. Chính vì vậy, các kết quả thu được từ nghiên cứu này cần được hiểu rõ trong khuôn khổ của một quy trình \textbf{đối chiếu nội bộ khách quan}, dựa trên các thiết lập hệ thống đã được trình bày chi tiết tại Chương~\ref{sec:methods} \cite{zhang2025pima,kirbas2025shap,smith1988}.

\section{Hạn chế và các mối đe dọa đối với tính giá trị}
\label{subsec:disc-limitations}

\textbf{Phạm vi thiết kế.} Nghiên cứu này hiện được \textbf{giới hạn có chủ đích} trong việc khảo sát bốn thuật toán phân loại truyền thống cùng mô hình cơ sở (\textit{majority baseline}), theo đó không bao gồm các phương pháp thúc đẩy (\textit{boosting}) hay các kiến trúc mạng thần kinh sâu. Toàn bộ thực nghiệm mới chỉ được triển khai duy nhất trên tập dữ liệu Pima với quy mô mẫu còn hạn chế, đồng thời \textbf{chưa} tiến hành các bước kiểm chứng độc lập (\textit{external validation}) trên những nguồn dữ liệu khác. Mặc dù những giới hạn này không làm giảm đi giá trị cốt lõi của quy trình thực nghiệm, song chúng cần được cân nhắc một cách kỹ lưỡng khi đánh giá khả năng tổng quát hóa các kết quả thu được trên những quần thể bệnh nhân khác nhau \cite{pedregosa2011}.

\textbf{Các mối đe dọa đối với tính giá trị} hiện được phân tích dựa trên hệ thống khung lý thuyết chuẩn như sau:

\begin{itemize}
  \item \textbf{Tính giá trị nội tại (\textit{Internal Validity}):} Các rủi ro tiềm ẩn bao gồm hiện tượng quá khớp (\textit{overfitting}) phát sinh trong giai đoạn tối ưu hóa siêu tham số, cùng với nguy cơ rò rỉ dữ liệu (\textit{data leakage}) nếu quy trình tiền xử lý không tuân thủ nghiêm ngặt cấu trúc luồng xử lý (\textit{pipeline}). Bên cạnh đó, độ nhạy của các mô hình như RF và SVM trước các hạt giống ngẫu nhiên (\textit{seed}) cũng là một yếu tố cần được kiểm soát chặt chẽ.
  \item \textbf{Tính giá trị ngoại tại (\textit{External Validity}):} Vấn đề này liên quan đến khả năng mở rộng kết quả thực nghiệm cho các khu vực địa lý, thời điểm hoặc các đặc điểm nhân khẩu học khác nhau. Ngoài ra, khoảng tin cậy (CI) ghi nhận được có thể vẫn còn ở mức rộng do những giới hạn nhất định của việc chỉ thực hiện 05 lượt phân đoạn dữ liệu.
  \item \textbf{Tính giá trị cấu trúc (\textit{Construct Validity}):} Khía cạnh này xem xét mức độ phản ánh mục tiêu lâm sàng thực tế thông qua việc gán nhãn nhị phân, đồng thời đánh giá sự tương thích giữa hệ chỉ số đã lựa chọn tại Mục~\ref{subsec:metrics} đối với mục đích tầm soát trong thực hành y khoa thực tiễn.
\end{itemize}

\noindent\textbf{Đạo đức và quản trị dữ liệu.} Nghiên cứu đòi hỏi sự minh bạch về các giới hạn của tính tổng quát, đặc biệt là những rủi ro liên quan đến \textbf{tính công bằng} (\textit{fairness}) khi áp dụng mô hình cho các nhóm dân tộc nằm ngoài quần thể Pima, song song với các yêu cầu khắt khe về \textbf{bảo mật dữ liệu cá nhân}. Mặc dù nghiên cứu đã nỗ lực triệt tiêu việc tối ưu hóa lặp lại trên cùng một tập kiểm thử, tuy nhiên với cỡ mẫu nhỏ, việc sử dụng \textbf{05 phân đoạn} dữ liệu dù là một bước tiến đáng kể nhưng vẫn có thể \textbf{chưa bao quát hết} mọi biến thiên tiềm ẩn trong dữ liệu. Do đó, các nghiên cứu tiếp theo nên cân nhắc sử dụng kỹ thuật \textit{bootstrapping} hoặc gia tăng số lượng vòng lặp thực nghiệm để nâng cao độ tin cậy.

Những hạn chế nêu trên không chỉ ảnh hưởng đến tính tổng quát mà còn tác động trực tiếp đến quá trình diễn giải kết quả thực nghiệm. Việc thiếu vắng bước kiểm chứng độc lập kéo theo hệ quả là hiệu năng quan sát được có thể sẽ không duy trì tính ổn định trên các quần thể có sự khác biệt về điều kiện thu thập dữ liệu. Do đó, các kết quả của luận văn này nên được tiếp nhận chủ yếu như một \textbf{minh chứng về phương pháp luận} đối với quy trình đánh giá mô hình, thay vì được xem là một bằng chứng đã sẵn sàng cho việc triển khai lâm sàng trong thực tế.

Trong nghiên cứu này, lựa chọn phương pháp \textit{hold-out} lặp lại được ưu tiên nhờ vào tính minh bạch cùng chi phí tính toán hợp lý. Tuy nhiên, kỹ thuật xác thực chéo lồng nhau (\textit{nested cross-validation}) sẽ là một hướng mở rộng phù hợp nhằm ước lượng khả năng tổng quát hóa của mô hình một cách chặt chẽ hơn. Bên cạnh đó, việc cỡ mẫu chỉ dừng lại ở mức $n=768$ đã làm gia tăng rủi ro về \textbf{phương sai ước lượng}. Điều này hoàn toàn nhất quán với mức độ biến thiên theo hạt giống ngẫu nhiên đã quan sát được, từ đó nhấn mạnh nhu cầu cấp thiết của việc xem xét hiệu năng đi kèm với các chỉ số bất định thống kê thay vì chỉ dựa vào trị số trung bình đơn thuần.

\section{Định hướng nghiên cứu tương lai (\textit{Future Work})}
\label{subsec:disc-future}

Các hướng nghiên cứu tiếp theo sẽ tập trung vào ba trụ cột chính bao gồm: \textbf{tính tổng quát hóa}, \textbf{năng lực mô hình} và \textbf{khả năng diễn giải}. Cụ thể là về tính tổng quát hóa, cần thực hiện quy trình \textbf{kiểm chứng độc lập} (\textit{external validation}) trên các nguồn dữ liệu đa trung tâm nhằm củng cố giá trị ngoại tại của mô hình. Về năng lực mô hình, nghiên cứu cần mở rộng khảo sát các thuật toán thuộc họ \textbf{Gradient Boosting} như XGBoost, LightGBM, CatBoost cùng các kiến trúc học sâu dựa trên cùng một giao thức triệt tiêu rò rỉ dữ liệu nghiêm ngặt. Đối với khả năng vận hành, nghiên cứu sẽ áp dụng các kỹ thuật \textbf{hiệu chuẩn hậu nghiệm} (\textit{post-hoc calibration}) trên một tập hiệu chỉnh biệt lập, song song với việc \textbf{tối ưu hóa ngưỡng quyết định} theo hàm chi phí lâm sàng thực tế. Ngoài ra, việc tích hợp các phương pháp như \textbf{SHAP} hoặc LIME sẽ cung cấp cái nhìn toàn diện hơn về cơ chế diễn giải ở cả cấp độ tổng thể lẫn cục bộ. Một vấn đề nghiên cứu trọng yếu cần tiếp tục làm rõ chính là mức độ tác động của giao thức đánh giá đến độ ổn định thứ hạng của mô hình trên các tập dữ liệu y sinh quy mô nhỏ. Đồng thời, việc lượng hóa sai số âm tính giả (FN) theo từng phân khúc thuộc tính lâm sàng sẽ giúp xác định chính xác các vùng dữ liệu mà mô hình chưa đạt hiệu năng kỳ vọng để từ đó đưa ra những điều chỉnh phù hợp \cite{kirbas2025shap}.

\noindent\textbf{Tổng kết Chương Thảo luận.} Tổng hợp lại, các phát hiện của luận văn không ủng hộ quan điểm về một ``thuật toán tối ưu tuyệt đối'' mà nhấn mạnh rằng việc lựa chọn mô hình phải dựa trên sự tương thích giữa \textbf{đặc thù dữ liệu} và \textbf{mục tiêu ứng dụng cụ thể}. Quan trọng hơn, thứ hạng của các mô hình không phải là một đặc tính bất biến mà phụ thuộc đáng kể vào giao thức thực nghiệm, phương thức phân đoạn dữ liệu và hệ thống chỉ số đo lường, đặc biệt là khi biên độ bất định thống kê vẫn còn lớn. Giá trị cốt lõi của nghiên cứu nằm ở việc xây dựng một khung đối chiếu minh bạch và có khả năng tái lập, đồng thời chỉ rõ những sự đánh đổi tất yếu khi chuyển dịch từ môi trường thực nghiệm sang triển khai ứng dụng. Những kết luận này đã giải quyết trực tiếp khoảng trống về tính thiếu nhất quán và khó so sánh giữa các công trình nghiên cứu trước đây như đã trình bày tại Chương~\ref{sec:literature}.

\chapter{Kết luận và kiến nghị}
\label{sec:conclusion}

Nghiên cứu này đã thiết lập thành công một \textbf{quy trình thực nghiệm hệ thống} nhằm đối chiếu khách quan hiệu năng của các thuật toán phân loại truyền thống trên bộ dữ liệu đái tháo đường Pima Indians. Khung thực nghiệm được xây dựng dựa trên sự kết hợp chặt chẽ giữa \textbf{phương pháp hold-out lặp lại} trên nhiều hạt giống ngẫu nhiên và quy trình \textbf{xác thực chéo phân tầng} để tinh chỉnh siêu tham số. Song song với đó, nghiên cứu áp dụng kỹ thuật \textbf{\texttt{GridSearchCV}} nhằm tối ưu hóa theo tiêu chí $F_1$, kết hợp với một luồng xử lý (\textit{pipeline}) tiền xử lý nhất quán. Những nỗ lực này nhằm bảo đảm tối đa tính \textbf{khách quan trong đối chiếu}, triệt tiêu hoàn toàn hiện tượng rò rỉ dữ liệu (\textit{data leakage}) và tăng cường \textbf{khả năng tái lập} cho toàn bộ công trình nghiên cứu (chi tiết tại Chương~\ref{sec:methods} và Mục~\ref{subsec:protocol}) \cite{pedregosa2011}.

Các kết quả thực nghiệm cho thấy những mô hình được khảo sát đạt mức hiệu năng \textbf{tương đương và cạnh tranh sát nút} trên hệ thống chỉ số đánh giá. Cụ thể là trong khi \textbf{Rừng ngẫu nhiên (RF)} chiếm ưu thế về chỉ số $F_1$-score trung bình, thì \textbf{Hồi quy Logistic (LR)} lại thể hiện ưu thế về độ chính xác (\textit{Accuracy}) cùng chỉ số ROC--AUC. Thực tế này minh chứng cho sự đánh đổi tất yếu giữa năng lực mô hình hóa phi tuyến và tính ổn định của các kiến trúc tuyến tính trong điều kiện cỡ mẫu hạn chế. Kết quả kiểm định Wilcoxon trên chỉ số $F_1$ giữa RF và KNN \textbf{không} cho thấy sự khác biệt có ý nghĩa thống kê tại ngưỡng $0,05$ với 05 lượt phân đoạn dữ liệu như đã trình bày tại Mục~\ref{subsec:disc-results}. Tuy nhiên, cần lưu ý rằng quy mô lượt chạy hạn chế có thể làm \textbf{suy giảm lực thống kê} (\textit{statistical power}), do đó việc chưa đạt được ý nghĩa thống kê không đồng nghĩa với sự đồng nhất hoàn toàn về bản chất giữa các thuật toán. Kết quả này một lần nữa tái khẳng định rủi ro khi chỉ dựa vào một lần phân chia dữ liệu duy nhất để xếp hạng mô hình, đồng thời nhấn mạnh nhu cầu cấp thiết của việc sử dụng các khoảng tin cậy khi thực hiện so sánh những thuật toán có kết quả tiệm cận nhau.

Bên cạnh đó, các công đoạn liên quan đến \textbf{tiền xử lý}, \textbf{tối ưu hóa siêu tham số} và \textbf{giao thức đánh giá lặp} được xác định là những nhân tố then chốt, đóng vai trò quan trọng tương đương với việc lựa chọn thuật toán. Kết quả phân tích thành phần (\textit{ablation study}) tại Mục~\ref{subsec:ablation} đã củng cố thêm quan điểm cho rằng hiệu năng thực nghiệm cần phải được đánh giá trong toàn bộ bối cảnh của luồng xử lý, thay vì chỉ xem xét các bước thực hiện một cách biệt lập.

Mặc dù đã hoàn thành các mục tiêu đề ra về mặt phương pháp luận, song nghiên cứu vẫn tồn tại một số hạn chế nhất định liên quan đến \textbf{quy mô dữ liệu}, \textbf{tính đặc thù của quần thể} cũng như \textbf{phạm vi của các họ mô hình} được khảo sát. Các hướng mở rộng tiềm năng bao gồm việc kiểm chứng độc lập, sử dụng thuật toán tổ hợp mạnh hơn và diễn giải mô hình chuyên sâu hiện đã được thảo luận tại Mục~\ref{subsec:disc-future}, song song với đó là các vấn đề về tính giá trị (\textit{validity}) đã được đặc tả chi tiết tại Mục~\ref{subsec:disc-limitations}.

Quan trọng hơn hết, nghiên cứu chỉ ra rằng chênh lệch hiệu năng giữa các thuật toán thường \textbf{thấp hơn so với kỳ vọng} khi áp dụng các giao thức đánh giá nghiêm ngặt. Phát hiện này nhấn mạnh rằng chất lượng của thiết kế thực nghiệm đóng vai trò tiên quyết, thậm chí còn quan trọng hơn cả việc lựa chọn một thuật toán cụ thể nào đó trong quá trình xây dựng hệ thống dự báo.

\textbf{Thông điệp chủ đạo:} Đối với các tập dữ liệu y sinh dạng bảng có kích thước hạn chế, việc \textbf{thiết kế giao thức thực nghiệm} cụ thể là các bước phân đoạn dữ liệu, tiền xử lý gắn liền với tập huấn luyện và tối ưu hóa siêu tham số có khả năng \textbf{thay đổi thứ hạng mô hình} tương đương với việc thay đổi thuật toán cốt lõi. Chính vì vậy, mọi công bố kết quả thực nghiệm cần phải được gắn liền với các mô tả chi tiết về quy trình (\textit{protocol}) thực hiện cụ thể.

Tóm lại, đóng góp chính của luận văn này là xây dựng được một \textbf{khung đánh giá thực nghiệm minh bạch, có khả năng tái lập và ngăn chặn rò rỉ dữ liệu} nhằm mục tiêu đối chiếu công bằng giữa các mô hình học máy trên dữ liệu y sinh quy mô vừa và nhỏ. Mô hình hiện đang được lưu trữ tại tệp \texttt{resources/diabetes\_prediction\_model.pkl} đóng vai trò phục vụ mục đích minh chứng phương pháp luận, đồng thời \textbf{không} mang giá trị triển khai lâm sàng khi chưa trải qua các bước kiểm chứng độc lập nghiêm ngặt. Việc thực hành đánh giá một cách cẩn trọng và minh bạch chính là điều kiện tiên quyết để đảm bảo độ tin cậy trong các nghiên cứu học máy y tế hiện nay.

\appendix
\chapter*{Phụ lục}
\addcontentsline{toc}{chapter}{Phụ lục}

\noindent\textit{Quy ước:} Các tiểu mục trong phần này được đánh ký hiệu bằng các chữ cái in hoa như A, B, C nhằm phân biệt rõ ràng với hệ thống số hiệu của các chương nội dung chính, cụ thể là sẽ sử dụng cách gọi ``Phụ lục A'' thay vì tiếp nối theo số thứ tự chương.

\noindent Phần Phụ lục này tập hợp toàn bộ các \textbf{thông số kỹ thuật chi tiết} phục vụ mục đích tái lập và kiểm chứng độc lập. Các nội dung được trình bày bao gồm không gian tìm kiếm siêu tham số, hệ thống kết quả thực nghiệm chi tiết theo từng phân đoạn dữ liệu, cấu hình môi trường thực thi cùng các tham số tái lập hệ thống như mã nguồn, danh sách thư viện phụ thuộc và các hạt giống ngẫu nhiên đã sử dụng.

\phantomsection
\label{app:hyper}
\section*{Phụ lục A. Không gian tìm kiếm siêu tham số (\textit{Hyperparameter Search Space})}

Không gian tìm kiếm siêu tham số được trình bày tại Bảng~\ref{tab:hyper} hiện đã được thiết lập cố định trong tệp mã nguồn \texttt{resources/diabetes.ipynb} và được áp dụng một cách nhất quán thông qua quy trình \texttt{GridSearchCV}. Các giá trị trong lưới tham số này phản ánh những thực hành phổ biến đối với dữ liệu dạng bảng quy mô vừa, qua đó giúp cân bằng giữa \textbf{chi phí tính toán} và \textbf{độ bao phủ} các cấu hình mô hình khác biệt. Bên cạnh đó, một số tham số cụ thể, ví dụ như \texttt{kernel='rbf'} đối với mô hình SVM, đã được cố định nhằm duy trì phạm vi đối chiếu trong cùng một họ mô hình như mô tả tại Chương~\ref{sec:methods}. Lưới tham số này được giữ \textbf{nguyên trạng} cho toàn bộ các lượt thực thi, đồng thời mọi thay đổi phát sinh đều được đồng bộ hóa chặt chẽ giữa nội dung bảng này và các tệp kết quả lưu trữ (\textit{artifacts}).

\begin{table}[htbp]
  \centering
  \small
  \caption{Chi tiết không gian siêu tham số trong quy trình \texttt{GridSearchCV}. Lưu ý: Tham số \texttt{class\_weight=None} được áp dụng thống nhất cho LR, SVM và RF.}
  \label{tab:hyper}
  \begin{tabular}{lp{8.5cm}}
    \toprule
    \textbf{Thuật toán} & \textbf{Lưới tham số / Cấu hình cố định} \\
    \midrule
    Hồi quy Logistic & \texttt{max\_iter=10000}, \texttt{solver='lbfgs'}, \texttt{C} $\in \{0{,}01; 0{,}1; 1; 10\}$ \\
    SVM (nhân RBF) & \texttt{kernel='rbf'}, \texttt{probability=True}, \texttt{C} $\in \{0{,}1; 1; 10\}$, \texttt{gamma} $\in \{\text{'scale'}; 0{,}01; 0{,}1\}$ \\
    Rừng ngẫu nhiên & \texttt{n\_estimators} $\in \{100; 200\}$, \texttt{max\_depth} $\in \{\text{None}; 8; 16\}$, \texttt{min\_samples\_leaf} $\in \{1; 2\}$ \\
    $k$-Láng giềng gần nhất & \texttt{n\_neighbors} $\in \{3; 5; 7; 9\}$, \texttt{weights} $\in \{\text{'uniform'}; \text{'distance'}\}$, \texttt{p=2} \\
    Đường cơ sở (Majority) & \texttt{DummyClassifier(strategy='most\_frequent')} \\
    \bottomrule
  \end{tabular}
\end{table}

\phantomsection
\label{app:per-split}
\section*{Phụ lục B. Kết quả chi tiết theo từng phân đoạn (\textit{Per-split Outputs})}

Sau mỗi lượt thực thi hoàn chỉnh, hệ thống sẽ tự động kết xuất tệp tin \texttt{results\_by\_split.csv} và lưu trữ tại thư mục \texttt{resources/outputs/latest/}, đồng thời tạo một bản sao lưu định danh theo thời gian thực tại \texttt{resources/outputs/runs/<timestamp>/}. Tệp tin này ghi nhận đầy đủ các chỉ số bao gồm \texttt{accuracy}, \texttt{precision}, \texttt{recall}, \texttt{f1} và \texttt{roc\_auc} tương ứng với từng giá trị hạt giống ngẫu nhiên (\texttt{split\_seed}). Việc công bố dữ liệu chi tiết \textbf{theo từng hạt giống} nhằm minh bạch hóa \textbf{biên độ biến thiên} của mô hình, điều này không chỉ giúp người đọc theo dõi sự ổn định của thuật toán mà còn cung cấp dữ liệu nền tảng cho bảng thống kê tổng hợp tại Chương~\ref{sec:results} (Bảng~\ref{tab:results-meanstd}). Bảng~\ref{tab:appendix-f1-by-split} dưới đây trích dẫn riêng chỉ số $F_1$-score để phục vụ mục đích tra cứu nhanh, trong khi các thông số chi tiết khác có thể được truy xuất trực tiếp từ tệp mã nguồn tương ứng.

\begin{table}[htbp]
  \centering
  \small
  \setlength{\tabcolsep}{6pt}
  \caption{Chỉ số $F_1$-score trên tập kiểm thử theo từng \texttt{split\_seed} (trích xuất từ \texttt{results\_by\_split.csv}).}
  \label{tab:appendix-f1-by-split}
  \begin{tabular}{lccccc}
    \toprule
    \textbf{Thuật toán} & \textbf{Seed 0} & \textbf{Seed 1} & \textbf{Seed 2} & \textbf{Seed 3} & \textbf{Seed 4} \\
    \midrule
    Hồi quy Logistic      & 0{,}609 & 0{,}584 & 0{,}565 & 0{,}596 & 0{,}598 \\
    SVM (nhân RBF)        & 0{,}646 & 0{,}584 & 0{,}535 & 0{,}543 & 0{,}550 \\
    Rừng ngẫu nhiên       & 0{,}646 & 0{,}583 & 0{,}602 & 0{,}606 & 0{,}633 \\
    $k$-Láng giềng gần nhất & 0{,}592 & 0{,}667 & 0{,}505 & 0{,}598 & 0{,}615 \\
    Đường cơ sở (Majority) & 0{,}000 & 0{,}000 & 0{,}000 & 0{,}000 & 0{,}000 \\
    \bottomrule
  \end{tabular}
\end{table}

\noindent Hệ thống các tệp tin bổ trợ đi kèm trong thư mục kết quả bao gồm: \texttt{summary\_mean\_std.csv} phục vụ thống kê mô tả, \texttt{f1\_ci95\_by\_model.csv} cung cấp khoảng tin cậy $95\%$, cùng các tệp \texttt{brier\_by\_split.csv} và \texttt{brier\_mean\_std.csv} dùng để định lượng chỉ số hiệu chuẩn. Bên cạnh đó, các tệp \texttt{pr\_auc\_by\_split.csv} và \texttt{pr\_auc\_mean\_std.csv} ghi nhận diện tích dưới đường cong PR, trong khi \texttt{roc\_auc\_bootstrap\_ci\_seed0.csv} lưu trữ ước lượng bootstrap cho AUC và \texttt{cliffs\_delta\_f1\_pairs.csv} mô tả độ lớn hiệu ứng. Danh mục này còn bao gồm các kết quả về tầm quan trọng đặc trưng như \texttt{permutation\_importance\_by\_split.csv} và \texttt{permutation\_importance\_consensus.csv}, tệp kiểm định thống kê \texttt{wilcoxon\_f1.txt}, thư mục đồ thị \texttt{figures/} cùng tệp cấu trúc định danh \texttt{manifest.json}.

\phantomsection
\label{app:environment}
\section*{Phụ lục C. Môi trường thực nghiệm (\textit{Experimental Environment})}

Các kết quả định lượng trong nghiên cứu được thực hiện trên cấu hình hệ thống bao gồm \textbf{Python} 3.13.12, \textbf{scikit-learn} 1.8.0 và \textbf{NumPy} 2.4.3, trong đó các thông số chi tiết đã được đặc tả tại tệp \texttt{resources/outputs/latest/manifest.json}. Toàn bộ các thư viện phụ thuộc phục vụ luồng xử lý chính như \textbf{pandas}, \textbf{SciPy} và \textbf{matplotlib} đều được khai báo trong tệp \texttt{pyproject.toml} và khóa phiên bản chặt chẽ thông qua \texttt{uv.lock}. Thực nghiệm được triển khai trên hệ điều hành \textbf{Linux} với kiến trúc CPU đa nhân, đồng thời nghiên cứu \textbf{không} sử dụng tăng tốc phần cứng GPU trong quy trình thực thi chính thức. Song song với đó, việc quản trị môi trường và cài đặt phần mềm được thực hiện đồng nhất thông qua công cụ \texttt{uv} tại thư mục gốc của kho mã nguồn.

\phantomsection
\label{app:pipeline}
\section*{Phụ lục D. Tóm tắt luồng công việc (\textit{End-to-End Workflow})}

Trình tự logic của nghiên cứu tuân thủ nghiêm ngặt giao thức thực nghiệm đã được trình bày tại Mục~\ref{subsec:protocol}, cụ thể bao gồm các giai đoạn sau: 
(1) Thực hiện chuyển đổi quy ước giá trị thiếu và gán giá trị trung vị (\textit{median imputation}) tích hợp trực tiếp trong luồng xử lý (\textit{pipeline}); 
(2) Phân đoạn tập dữ liệu Huấn luyện và Kiểm thử (\textit{Train/Test}) theo tỷ lệ $80/20$ có áp dụng phương thức phân tầng (\textit{stratify}); 
(3) Triển khai quy trình \texttt{GridSearchCV} kết hợp với \texttt{StratifiedKFold} ($k=5$) trên tập huấn luyện nhằm tối ưu hóa các mô hình dựa theo tiêu chí mục tiêu là chỉ số $F_1$-score; 
(4) Tiến hành tái huấn luyện mô hình trên toàn bộ tập dữ liệu huấn luyện tương ứng với cấu hình siêu tham số tối ưu đã tìm được; 
(5) Thực hiện đánh giá hiệu năng duy nhất một lần trên tập kiểm thử biệt lập (\textit{hold-out test set}). 

Bên cạnh đó, các tập lệnh bổ trợ phục vụ cho quá trình hậu xử lý bao gồm: \texttt{resources/compute\_calibration.py} dùng để tái lập chỉ số Brier cùng biểu đồ hiệu chuẩn; \texttt{resources/compute\_extended\_metrics.py} nhằm tái lập chỉ số PR--AUC, thực hiện \textit{bootstrap} ROC--AUC, tính toán độ lớn hiệu ứng Cliff's delta và trực quan hóa các đường cong PR; song song với đó là tệp \texttt{resources/compute\_permutation\_importance.py} phục vụ việc tổng hợp tầm quan trọng đặc trưng và xây dựng đồ thị diễn giải. Cần nhấn mạnh rằng mọi bước ước lượng tham số từ dữ liệu đều chỉ thực hiện lệnh \texttt{fit} trong phạm vi tập huấn luyện, qua đó đảm bảo tính khách quan tuyệt đối và triệt tiêu hoàn toàn hiện tượng rò rỉ dữ liệu.

\phantomsection
\label{app:repro}
\label{app:code}
\label{app:seeds}
\section*{Phụ lục E. Tính tái lập và khả năng truy cập mã nguồn (\textit{Reproducibility \& Code Availability})}

Nhằm mục đích tái lập toàn bộ kết quả thực nghiệm của luận văn, các thành phần cốt lõi hiện đã được chuẩn hóa và cung cấp một cách đầy đủ, bao gồm: (i) Notebook thực thi chính tại tệp \texttt{resources/diabetes.ipynb} cùng tập lệnh kết xuất dữ liệu \texttt{resources/export\_run.py}; (ii) Công cụ quản trị \texttt{uv} đi kèm các tệp tin cấu hình môi trường \texttt{pyproject.toml} và \texttt{uv.lock}; (iii) Kho lưu trữ mã nguồn công khai được đính kèm cùng bản báo cáo chính thức.

\textbf{Cấu hình hạt giống ngẫu nhiên (\textit{Seeds}):} 
Thiết lập \texttt{random\_state=42} được áp dụng thống nhất cho các tiến trình nội bộ như \texttt{StratifiedKFold}, \texttt{GridSearchCV} cùng các thuật toán LR, SVM và RF nhằm đảm bảo tính ổn định tối đa cho mô hình. Song song với đó, các giá trị hạt giống lần lượt thuộc tập $\{0, 1, 2, 3, 4\}$ sẽ được áp dụng cho hàm \texttt{train\_test\_split} để tạo ra năm lượt phân đoạn dữ liệu hoàn toàn độc lập. Đối với thuật toán \textbf{KNN}, do bản chất của phương pháp này không chứa các tham số ngẫu nhiên nội tại, vì vậy kết quả thu được sẽ được cố định dựa trên cấu trúc dữ liệu cùng hệ thống chỉ số đã thiết lập như đã mô tả chi tiết tại Mục~\ref{subsec:repro}.

\noindent Việc công khai toàn bộ mã nguồn, thực hiện khóa phiên bản phụ thuộc và cung cấp đầy đủ các kết quả trung gian không chỉ giúp triệt tiêu các sai lệch hệ thống nằm ngoài tầm kiểm soát, mà còn góp phần tăng cường độ tin cậy cũng như mở ra khả năng mở rộng cho các nghiên cứu trong tương lai.

\newpage
\printbibliography[title={Tài liệu tham khảo}]

\end{document}
